% Applications linéaires du plan — Mathématiques Tle TI — Cours signé OnBuch+
% Programme officiel MINESEC. Compiler avec : tectonic M19-applications-lineaires-du-plan.tex
\newcommand{\DOCMATIERE}{Mathématiques}
\newcommand{\DOCNIVEAU}{Tle TI}
\newcommand{\DOCMODULE}{Configurations et transformations élémentaires du plan}
\newcommand{\DOCLECON}{19}
\newcommand{\DOCTITRE}{Applications linéaires du plan}
% OnBuch+ — préambule « cours signés » ENRICHI (compilé avec tectonic / XeLaTeX)
% Système visuel « Pop » d'OnBuch+ : fond crème (jamais blanc), bordures encre
% épaisses, ombres « dures » décalées, titres Archivo Black en capitales.
% Version MATHÉMATIQUES : graphes (pgfplots/tikz), boîtes pédagogiques
% (définition, propriété, méthode, attention, le savais-tu, exercice résolu,
% exercice, TP, bilan), page de garde et signature OnBuch+ ancrée en bas.
%
% Macros attendues dans main.tex AVANT \input{preamble} :
%   \DOCMATIERE  \DOCNIVEAU  \DOCMODULE  \DOCLECON (numéro)  \DOCTITRE
\documentclass[11pt]{article}

\usepackage{fontspec}
\usepackage[french]{babel}
\usepackage[a4paper,margin=2cm,top=3.4cm,bottom=2.8cm,headheight=1.4cm,footskip=1.3cm]{geometry}
\usepackage{amsmath,amssymb,amsfonts}
\usepackage{mathtools} % \xmapsto, \coloneqq... souvent utilisés par les modèles
\usepackage{cancel} % simplifications barrées
% \abs{z} (module, valeur absolue) et \norm{u} (norme), utilisés par les modèles
\providecommand{\abs}{}\let\abs\relax\DeclarePairedDelimiter{\abs}{\lvert}{\rvert}
\providecommand{\norm}{}\let\norm\relax\DeclarePairedDelimiter{\norm}{\lVert}{\rVert}
\usepackage[table]{xcolor} % [table] : fournit aussi \rowcolors (alternance de lignes)
\usepackage{pagecolor}
\usepackage{tikz}
\usetikzlibrary{arrows.meta,positioning,calc,shapes.geometric,shapes.misc,shapes.arrows,decorations.pathmorphing,decorations.markings,patterns,shadows,mindmap,trees,fit,backgrounds,matrix,intersections}
\usepackage{pgfplots}
\usepackage{pgf-pie} % diagrammes circulaires \pie (statistiques)
\pgfplotsset{compat=1.18}
\usepgfplotslibrary{fillbetween,groupplots} % aires entre courbes (calcul intégral)
\usepackage{enumitem}
\usepackage{fancyhdr}
% [most] charge toutes les bibliothèques tcolorbox (listings, minted, raster,
% documentation...) et gonfle inutilement la mémoire TeX sur un document long
% (~300 boîtes) : on ne charge que ce qui est réellement utilisé.
\usepackage[skins,breakable]{tcolorbox}
\usepackage{graphicx}
\usepackage{colortbl}
\usepackage{booktabs}
\usepackage{tabularx}
\usepackage{diagbox} % case d'en-tête diagonale des tableaux à double entrée
% Colonnes extensibles centrée / gauche / droite (C, L, R), souvent utilisées par les modèles
\newcolumntype{C}{>{\centering\arraybackslash}X}
\newcolumntype{L}{>{\raggedright\arraybackslash}X}
\newcolumntype{R}{>{\raggedleft\arraybackslash}X}
\usepackage{array}
\usepackage{multicol}
\usepackage{siunitx}
% parse-numbers=false : accepte aussi \SI{2,5 \times 10^{-3}}{...}
\sisetup{locale=FR,output-decimal-marker={,},group-minimum-digits=4,parse-numbers=false}
\DeclareSIUnit\ohm{\ensuremath{\symup{\Omega}}} % U+2126 absent de la police
\usepackage{qrcode}
% \degree : commande fréquemment utilisée par les modèles pour un angle ou une
% température en dehors de \SI{}{} (non fournie par les paquets chargés).
\providecommand{\degree}{^{\circ}}
% \qed (fin de démonstration), utilisé par les modèles sans amsthm
\providecommand{\qed}{\hfill$\square$}
% \barre : double barre des valeurs interdites dans un tableau de variations (tkz-tab)
\providecommand{\barre}{\ensuremath{\|}}
% environnement proof (amsthm non chargé) utilisé par les modèles
\ifdefined\proof\else\newenvironment{proof}[1][Démonstration]{\par\noindent\textit{#1.}\ }{\qed\par}\fi
\usepackage{lastpage}
\usepackage{hyperref}
\hypersetup{hidelinks}

% ============================================================
% Palette « Pop » OnBuch+ — mêmes hex que la classe P (plus_ui.dart)
% ============================================================
\definecolor{poporange}{HTML}{F59321}
\definecolor{popdark}{HTML}{E07A0C}
\definecolor{popcream}{HTML}{FFF6E8}
\definecolor{popink}{HTML}{1C1714}
\definecolor{poppurple}{HTML}{7A5AE0}
\definecolor{popgreen}{HTML}{2E9E6A}
\definecolor{poppink}{HTML}{E0457B}
\definecolor{popblue}{HTML}{2F7DE1}
\definecolor{popgold}{HTML}{C98A12}
\definecolor{popmuted}{HTML}{8A7F76}
\definecolor{popline}{HTML}{EADFD2}
\definecolor{popgray}{HTML}{8A7F76}
\colorlet{popgrey}{popgray}
% Alias tolérants (le modèle nomme parfois les couleurs différemment)
\colorlet{popwhite}{white}
\colorlet{popblack}{popink}
\colorlet{popred}{poppink}
\colorlet{popyellow}{popgold}
% Teintes claires pour fonds de tableaux / zones
\colorlet{popblueL}{popblue!10!white}
\colorlet{poporangeL}{poporange!14!white}
\colorlet{popgreenL}{popgreen!12!white}
\colorlet{poppurpleL}{poppurple!12!white}
\colorlet{poppinkL}{poppink!12!white}
\colorlet{popgoldL}{popgold!14!white}

\pagecolor{popcream}

% ============================================================
% Typographie
% ============================================================
\defaultfontfeatures{Ligatures=TeX}
\setmainfont{PlusJakartaSans-Medium.ttf}[Path=../../../fonts/,BoldFont=PlusJakartaSans-Bold.ttf]
% Maths en sans-serif (Fira Math) avec chiffres et lettres droites de Plus
% Jakarta, pour que les formules se fondent dans le texte.
\usepackage[math-style=ISO,bold-style=ISO]{unicode-math}
\setmathfont{FiraMath-Regular.otf}[Path=../../../fonts/]
\setmathfont{PlusJakartaSans-Medium.ttf}[Path=../../../fonts/,range={up/num,up/{Latin,latin},\mathup}]
\usepackage{icomma} % virgule décimale sans espace en mode math
% Caractères mathématiques Unicode absents des polices de texte (Plus Jakarta,
% Archivo Black) : rendus par leur équivalent mathématique, en texte comme en
% formule — sinon ils s'affichent en carré vide (ex. « Arithmétique dans ℤ »).
\usepackage{newunicodechar}
\newunicodechar{ℝ}{\ensuremath{\mathbb{R}}}
\newunicodechar{ℤ}{\ensuremath{\mathbb{Z}}}
\newunicodechar{ℂ}{\ensuremath{\mathbb{C}}}
\newunicodechar{ℕ}{\ensuremath{\mathbb{N}}}
\newunicodechar{ℚ}{\ensuremath{\mathbb{Q}}}
\newunicodechar{𝔻}{\ensuremath{\mathbb{D}}}
\newunicodechar{α}{\ensuremath{\alpha}}
\newunicodechar{β}{\ensuremath{\beta}}
\newunicodechar{γ}{\ensuremath{\gamma}}
\newunicodechar{δ}{\ensuremath{\delta}}
\newunicodechar{ε}{\ensuremath{\varepsilon}}
\newunicodechar{θ}{\ensuremath{\theta}}
\newunicodechar{λ}{\ensuremath{\lambda}}
\newunicodechar{μ}{\ensuremath{\mu}}
\newunicodechar{σ}{\ensuremath{\sigma}}
\newunicodechar{τ}{\ensuremath{\tau}}
\newunicodechar{φ}{\ensuremath{\varphi}}
\newunicodechar{ψ}{\ensuremath{\psi}}
\newunicodechar{ω}{\ensuremath{\omega}}
\newunicodechar{Σ}{\ensuremath{\Sigma}}
% majuscules produites par \MakeUppercase dans les titres (α-aminés -> Α)
\newunicodechar{Α}{\ensuremath{\alpha}}
\newunicodechar{Β}{\ensuremath{\beta}}
\newunicodechar{Γ}{\ensuremath{\Gamma}}
\newunicodechar{Δ}{\ensuremath{\Delta}}
\newunicodechar{Ε}{\ensuremath{\varepsilon}}
\newunicodechar{Θ}{\ensuremath{\Theta}}
\newunicodechar{Λ}{\ensuremath{\Lambda}}
\newunicodechar{Μ}{\ensuremath{\mu}}
\newunicodechar{Τ}{\ensuremath{\tau}}
\newunicodechar{Φ}{\ensuremath{\Phi}}
\newunicodechar{Ψ}{\ensuremath{\Psi}}
\newunicodechar{Ω}{\ensuremath{\Omega}}
\newunicodechar{↦}{\ensuremath{\mapsto}}
\newunicodechar{∈}{\ensuremath{\in}}
\newunicodechar{∉}{\ensuremath{\notin}}
\newunicodechar{∧}{\ensuremath{\wedge}}
\newunicodechar{⇔}{\ensuremath{\Leftrightarrow}}
\newunicodechar{⟺}{\ensuremath{\Longleftrightarrow}}
\newunicodechar{⇒}{\ensuremath{\Rightarrow}}
\newunicodechar{⟹}{\ensuremath{\Longrightarrow}}
\newunicodechar{∘}{\ensuremath{\circ}}
\newunicodechar{∪}{\ensuremath{\cup}}
\newunicodechar{∩}{\ensuremath{\cap}}
\newunicodechar{≡}{\ensuremath{\equiv}}
\newunicodechar{⊂}{\ensuremath{\subset}}
\newunicodechar{⊥}{\ensuremath{\perp}}
\newunicodechar{∃}{\ensuremath{\exists}}
\newunicodechar{∀}{\ensuremath{\forall}}
\newunicodechar{∛}{\ensuremath{\sqrt[3]{\,}}}
\newunicodechar{∜}{\ensuremath{\sqrt[4]{\,}}}
\newunicodechar{⃗}{\ensuremath{{}^{\to}}}
\newunicodechar{ⁿ}{\ifmmode^{n}\else\textsuperscript{\ensuremath{n}}\fi}
\newunicodechar{ˣ}{\ifmmode^{x}\else\textsuperscript{\ensuremath{x}}\fi}
\newunicodechar{ᵇ}{\ifmmode^{b}\else\textsuperscript{\ensuremath{b}}\fi}
\newunicodechar{ʳ}{\ifmmode^{r}\else\textsuperscript{\ensuremath{r}}\fi}
\newunicodechar{ᵘ}{\ifmmode^{u}\else\textsuperscript{\ensuremath{u}}\fi}
\newunicodechar{ʸ}{\ifmmode^{y}\else\textsuperscript{\ensuremath{y}}\fi}
\newunicodechar{ᵃ}{\ifmmode^{a}\else\textsuperscript{\ensuremath{a}}\fi}
\newunicodechar{ᵉ}{\ifmmode^{e}\else\textsuperscript{\ensuremath{e}}\fi}
\newunicodechar{ᵏ}{\ifmmode^{k}\else\textsuperscript{\ensuremath{k}}\fi}
\newunicodechar{ᵖ}{\ifmmode^{p}\else\textsuperscript{\ensuremath{p}}\fi}
\newunicodechar{ᴺ}{\ifmmode^{N}\else\textsuperscript{\ensuremath{N}}\fi}
\newunicodechar{ᶜ}{\ifmmode^{c}\else\textsuperscript{\ensuremath{c}}\fi}
\newunicodechar{ʰ}{\ifmmode^{h}\else\textsuperscript{\ensuremath{h}}\fi}
\newunicodechar{ᵠ}{\ifmmode^{q}\else\textsuperscript{\ensuremath{q}}\fi}
\newunicodechar{ₙ}{\ifmmode_{n}\else\textsubscript{\ensuremath{n}}\fi}
\newunicodechar{ₐ}{\ifmmode_{a}\else\textsubscript{\ensuremath{a}}\fi}
\newunicodechar{ᵢ}{\ifmmode_{i}\else\textsubscript{\ensuremath{i}}\fi}
\newunicodechar{ᵣ}{\ifmmode_{r}\else\textsubscript{\ensuremath{r}}\fi}
\newunicodechar{ₖ}{\ifmmode_{k}\else\textsubscript{\ensuremath{k}}\fi}
\newunicodechar{ᵦ}{\ifmmode_{\beta}\else\textsubscript{\ensuremath{\beta}}\fi}
\newunicodechar{ₓ}{\ifmmode_{x}\else\textsubscript{\ensuremath{x}}\fi}
\newfontfamily\popdisplay{ArchivoBlack-Regular.ttf}[Path=../../../fonts/]
\newfontfamily\poplogo{SpaceGrotesk-Bold.ttf}[Path=../../../fonts/]
\newfontfamily\popsemi{PlusJakartaSans-SemiBold.ttf}[Path=../../../fonts/]
\newfontfamily\popxbold{PlusJakartaSans-ExtraBold.ttf}[Path=../../../fonts/]
\newfontfamily\popxboldls{PlusJakartaSans-ExtraBold.ttf}[Path=../../../fonts/,LetterSpace=3.0]

\setlist[enumerate]{leftmargin=1.7em,itemsep=5pt,topsep=4pt}
\setlist[itemize]{leftmargin=1.7em,itemsep=4pt,topsep=4pt,label={\color{poporange}\textbullet}}
\setlength{\parindent}{0pt}
\setlength{\parskip}{5pt}
\renewcommand{\arraystretch}{1.25}

% pgfplots : style Pop par défaut
\pgfplotsset{
  every axis/.append style={axis line style={popink,line width=1pt},tick style={popink},
    grid style={popline,line width=0.5pt},label style={font=\small},tick label style={font=\footnotesize}},
  popcurve/.style={poporange,line width=1.8pt,no markers,smooth},
}
% Même style disponible dans un \draw TikZ simple (hors pgfplots)
\tikzset{popcurve/.style={poporange,line width=1.8pt}}

% ============================================================
% Pill / squiggle
% ============================================================
\newcommand{\poppill}[3]{%
  \tikz[baseline=-0.55ex]\node[draw=popink,line width=1.2pt,rounded corners=2.6pt,%
    fill=#2,inner xsep=6.5pt,inner ysep=3pt]{%
    \popxboldls\fontsize{7}{7}\selectfont\color{#3}\MakeUppercase{#1}};%
}
\newcommand{\popsquiggle}[1]{%
  \tikz[baseline=-0.4ex]{
    \draw[#1,line width=2.6pt,line cap=round]
      plot [smooth] coordinates {(0,0.09) (0.34,0.01) (0.68,0.21) (1.02,0.01) (1.36,0.10)};
  }%
}
% Mot-clé surligné (marqueur orange sous le texte)
\newcommand{\cle}[1]{{\popxbold\color{popdark}#1}}

% ============================================================
% En-tête / pied
% ============================================================
\pagestyle{fancy}
\fancyhf{}
\renewcommand{\headrulewidth}{0pt}
\renewcommand{\footrulewidth}{0pt}
\lhead{\raisebox{-0.15ex}{\poplogo\fontsize{13}{13}\selectfont\color{popink}OnBuch\color{poporange}+}}
\rhead{\poppill{\DOCMATIERE\ \textbullet\ \DOCNIVEAU\ \textbullet\ Leçon \DOCLECON}{white}{popink}}
\lfoot{\popxbold\fontsize{7.6}{7.6}\selectfont\color{popmuted}\MakeUppercase{Cours signé OnBuch+}\ \textbullet\ {\popsemi\DOCTITRE}}
\rfoot{\tikz[baseline=-0.6ex]\node[rounded corners=8.5pt,fill=popink,minimum height=17pt,minimum width=17pt,inner xsep=5pt,inner sep=0]{\popxbold\fontsize{8}{8}\selectfont\color{popcream}\thepage\,/\,\pageref*{LastPage}};}

% ============================================================
% Titres
% ============================================================
\newcommand{\chaptitre}[2]{%
  \begin{tcolorbox}[enhanced,colback=poporange,colframe=popink,boxrule=2.6pt,arc=11pt,
      drop shadow=popink,left=15pt,right=15pt,top=12pt,bottom=11pt,before skip=3pt,after skip=18pt]
    \poppill{#2}{popink}{popcream}\par\vspace{9pt}
    {\raggedright\hyphenpenalty=10000\exhyphenpenalty=10000\popdisplay\fontsize{22}{24}\selectfont\color{popcream}\MakeUppercase{#1}\par}
  \end{tcolorbox}
}
% Section numérotée : pastille orange avec le numéro + titre Archivo
\newcounter{popsec}
\newcommand{\coursec}[1]{%
  \stepcounter{popsec}\setcounter{popsub}{0}%
  \par\vspace{16pt}\noindent
  \begin{minipage}{\linewidth}
  \tikz[baseline=-0.7ex]\node[draw=popink,line width=1.6pt,fill=poporange,rounded corners=4pt,minimum size=22pt,inner sep=0]{\popdisplay\fontsize{12}{12}\selectfont\color{popcream}\thepopsec};%
  \hspace{8pt}{\popdisplay\fontsize{15}{17}\selectfont\color{popink}\MakeUppercase{#1}}\par
  \vspace{4pt}\noindent{\color{poporange}\rule{36pt}{3.6pt}}
  \end{minipage}\par\nopagebreak\vspace{8pt}
}
\newcounter{popsub}
\newcommand{\courssub}[1]{%
  \stepcounter{popsub}%
  \par\vspace{9pt}\noindent{\popxbold\fontsize{11.5}{13}\selectfont\color{popink}{\color{poporange}\thepopsec.\thepopsub}\hspace{6pt}#1}\par\nopagebreak\vspace{4pt}
}

% ============================================================
% Boîtes pédagogiques — même recette : fond blanc, bordure épaisse colorée,
% ombre dure encre, badge plein. Titre optionnel [#1].
% ============================================================
\tcbset{popbase/.style={breakable,enhanced,colback=white,boxrule=2.3pt,arc=9pt,
  shadow={3pt}{-3pt}{0pt}{popink},left=14pt,right=14pt,top=11pt,bottom=11pt,before skip=11pt,after skip=15pt,
  fontupper=\popsemi\fontsize{10.6}{14.5}\selectfont\color{popink},
  fontlower=\popsemi\fontsize{10.6}{14.5}\selectfont\color{popink}}}
% Coche dessinée (le glyphe ✓ n'existe pas dans les polices du document)
\newcommand{\popcheck}[1]{\tikz[baseline=-0.5ex]\draw[#1,line width=1.6pt,line cap=round,line join=round](-0.13,0.0)--(-0.04,-0.09)--(0.13,0.1);}
\newcommand{\popboxhead}[3]{\poppill{#1}{#2}{white}\ifx\relax#3\relax\else\hspace{6pt}{\popxbold\fontsize{10.6}{12}\selectfont\color{popink}#3}\fi\par\vspace{7pt}}

\newtcolorbox{aretenir}[1][]{popbase,colframe=popgreen,before upper={\popboxhead{À retenir}{popgreen}{#1}}}
\newtcolorbox{exemplebox}[1][]{popbase,colframe=poporange,before upper={\popboxhead{Exemple}{poporange}{#1}}}
\newtcolorbox{definition}[1][]{popbase,colframe=popblue,before upper={\popboxhead{Définition}{popblue}{#1}}}
\newtcolorbox{propriete}[1][]{popbase,colframe=poppurple,before upper={\popboxhead{Propriété}{poppurple}{#1}}}
\newtcolorbox{methode}[1][]{popbase,colframe=popgold,before upper={\popboxhead{Méthode}{popgold}{#1}}}
\newtcolorbox{attention}[1][]{popbase,colframe=poppink,before upper={\popboxhead{Attention}{poppink}{#1}}}
\newtcolorbox{experience}[1][]{popbase,colframe=popblue,colback=popblueL,before upper={\popboxhead{Expérience}{popblue}{#1}}}
% « Le savais-tu ? » : carte violette pleine (contexte, culture, Cameroun)
\newtcolorbox{savaistu}[1][]{popbase,colback=poppurple,colframe=popink,
  fontupper=\popsemi\fontsize{10.4}{14.2}\selectfont\color{white},
  before upper={\poppill{Le savais-tu\,?}{white}{poppurple}\ifx\relax#1\relax\else\hspace{6pt}{\popxbold\color{white}#1}\fi\par\vspace{7pt}}}
% Exercice résolu : énoncé en haut, \tcblower puis la solution sur fond crème
\newcounter{popexo}\newcounter{popexr}
\newtcolorbox{exoresolu}[1][]{popbase,colframe=poporange,colbacklower=popcream,
  segmentation style={popink,line width=1.2pt,dashed},
  before upper={\stepcounter{popexr}\popboxhead{Exercice résolu \thepopexr}{poporange}{#1}},
  before lower={\poppill{Solution}{popink}{popcream}\par\vspace{6pt}}}
% Exercice (non résolu en place) — la correction est dans la partie Corrigés
\newtcolorbox{exercice}[1][]{popbase,colframe=popink,
  before upper={\stepcounter{popexo}\popboxhead{Exercice \thepopexo}{popink}{#1}}}
% Corrigé
\newtcolorbox{corrige}[1][]{popbase,colframe=popgreen,colback=popgreenL,
  before upper={\popboxhead{Corrigé}{popgreen}{#1}}}
% Objectifs / prérequis / bilan
\newtcolorbox{objectifs}{popbase,colframe=popink,colback=poporangeL,
  before upper={\popboxhead{Objectifs de la leçon}{popink}{}}}
\newtcolorbox{prerequis}{popbase,colframe=popink,colback=popblueL,
  before upper={\popboxhead{Prérequis}{popblue}{}}}
\newtcolorbox{bilan}{popbase,colframe=popink,colback=white,boxrule=2.8pt,
  before upper={\popboxhead{Fiche bilan}{poporange}{L'essentiel à connaître pour l'examen}}}
% Figure encadrée avec légende
\newtcolorbox{popfigure}[1][]{enhanced,breakable=false,colback=white,colframe=popline,boxrule=1.4pt,arc=8pt,
  left=10pt,right=10pt,top=10pt,bottom=8pt,before skip=10pt,after skip=12pt,halign=center,
  lower separated=false,halign lower=center,
  fontlower=\popsemi\fontsize{9}{11.5}\selectfont\color{popmuted},
  #1}
\newcommand{\legende}[1]{\par\vspace{6pt}{\popsemi\fontsize{9}{11.5}\selectfont\color{popmuted}#1}}

% ============================================================
% Signature OnBuch+
% ============================================================
\newcommand{\poplogoinline}[1]{{\poplogo\fontsize{#1}{#1}\selectfont\color{popink}OnBuch\color{poporange}+}}
\newcommand{\signatureonbuch}{%
  \begin{tcolorbox}[enhanced,colback=popink,colframe=popink,boxrule=2.2pt,arc=10pt,
      drop shadow=poporange,left=14pt,right=14pt,top=10pt,bottom=10pt,before skip=0pt,after skip=0pt]
    \tikz[baseline=-0.7ex]\node[circle,fill=popgreen,draw=popcream,line width=1.3pt,minimum size=22pt,inner sep=0]{\popcheck{white}};%
    \hspace{9pt}\begin{minipage}[c]{0.62\linewidth}
      {\popxbold\fontsize{12}{13}\selectfont\color{popcream}Signé\ \textbullet\ {\poplogo OnBuch\color{poporange}+}}\par\vspace{2pt}
      {\raggedright\popsemi\fontsize{8}{10.5}\selectfont\color{popline}\MakeUppercase{Équipe pédagogique OnBuch+}\\\MakeUppercase{Conforme au programme officiel MINESEC}\par}
    \end{minipage}\hfill
    {\poplogo\fontsize{22}{22}\selectfont\color{popcream}OnBuch\color{poporange}+}
  \end{tcolorbox}%
}

% ============================================================
% Page de garde
% #1 titre · #2 matière · #3 niveau · #4 module · #5 n° leçon · #6 durée ·
% #7 description / objectifs (texte ou itemize)
% ============================================================
\newcommand{\pagegarde}[7]{%
  \thispagestyle{empty}
  \newgeometry{a4paper,margin=1.7cm,top=1.6cm,bottom=1.4cm}%
  \begin{tikzpicture}[remember picture,overlay]
    \fill[poporange!18] ([xshift=-3.2cm,yshift=-3.6cm]current page.north east) circle (4.6cm);
    \fill[poppurple!14] ([xshift=2.4cm,yshift=7.6cm]current page.south west) circle (3.8cm);
  \end{tikzpicture}%
  {\poplogo\fontsize{30}{30}\selectfont\color{popink}OnBuch\color{poporange}+}\hfill
  \raisebox{0.6ex}{\poppill{Cours signé \textbullet\ Premium}{popink}{popcream}}\par
  \vspace{1pt}\popsquiggle{poppurple}\par
  \vspace{8pt}
  \begin{tcolorbox}[enhanced,colback=poporange,colframe=popink,boxrule=2.8pt,arc=13pt,
      drop shadow=popink,left=18pt,right=18pt,top=12pt,bottom=13pt,after skip=10pt]
    \poppill{#2 \textbullet\ #3}{popink}{popcream}\hspace{5pt}\poppill{#4}{white}{popink}\par\vspace{8pt}
    {\popdisplay\fontsize{14}{14}\selectfont\color{popink}LEÇON #5}\par\vspace{4pt}
    {\raggedright\hyphenpenalty=10000\exhyphenpenalty=10000\popdisplay\fontsize{25}{27}\selectfont\color{popcream}\MakeUppercase{#1}\par}
    \vspace{8pt}
    {\popxbold\fontsize{9.5}{9.5}\selectfont\color{popink}\MakeUppercase{Durée conseillée : #6}}
  \end{tcolorbox}
  \begin{tcolorbox}[enhanced,colback=white,colframe=popink,boxrule=2.2pt,arc=10pt,
      drop shadow=popink,left=16pt,right=16pt,top=10pt,bottom=10pt,after skip=8pt]
    \poppill{Dans cette leçon}{poppurple}{white}\par\vspace{5pt}
    \popsemi\fontsize{9.3}{12.6}\selectfont\color{popink}#7
  \end{tcolorbox}
  \noindent\begin{minipage}[t]{0.487\linewidth}
    \begin{tcolorbox}[enhanced,colback=popgreen,colframe=popink,boxrule=2.2pt,arc=10pt,
        drop shadow=popink,left=11pt,right=11pt,top=10pt,bottom=10pt,height=3.1cm,valign=center]
      \tcbox[colback=white,colframe=popink,boxrule=1.3pt,arc=5pt,boxsep=0pt,nobeforeafter,
        left=3pt,right=3pt,top=3pt,bottom=3pt,valign=center]{\qrcode[height=1.9cm]{https://play.google.com/store/apps/details?id=com.luvvix.onbuch&hl=fr}}%
      \hspace{8pt}\begin{minipage}[c]{0.5\linewidth}
        {\popdisplay\fontsize{11}{12.5}\selectfont\color{white}\MakeUppercase{Télécharge OnBuch}}\par\vspace{4pt}
        {\raggedright\popsemi\fontsize{8.4}{11}\selectfont\color{white}Gratuit \textbullet\ Google Play\\Tuteur IA, annales, concours, résultats\par}
      \end{minipage}
    \end{tcolorbox}
  \end{minipage}\hfill
  \begin{minipage}[t]{0.487\linewidth}
    \begin{tcolorbox}[enhanced,colback=poppurple,colframe=popink,boxrule=2.2pt,arc=10pt,
        drop shadow=popink,left=11pt,right=11pt,top=10pt,bottom=10pt,height=3.1cm,valign=center]
      \tikz[baseline=-0.7ex]\node[circle,fill=white,draw=popink,line width=1.3pt,minimum size=1.6cm,inner sep=0]{\popdisplay\fontsize{24}{24}\selectfont\color{poppurple}+};%
      \hspace{8pt}\begin{minipage}[c]{0.6\linewidth}
        {\popdisplay\fontsize{11}{12.5}\selectfont\color{white}\MakeUppercase{Passe à OnBuch+}}\par\vspace{4pt}
        {\raggedright\popsemi\fontsize{8.4}{11}\selectfont\color{white}Tous les cours signés, Tuteur IA illimité, fiches PDF — onglet OnBuch+ de l'app\par}
      \end{minipage}
    \end{tcolorbox}
  \end{minipage}
  \vspace{8pt}
  \popcontenu
  \vfill
  \signatureonbuch
  \restoregeometry
  \newpage
}

% Rangée « Ce cours contient » (page de garde)
\newcommand{\poptile}[3]{% #1 couleur #2 gros texte #3 légende
  \begin{tcolorbox}[enhanced,colback=#1,colframe=popink,boxrule=2pt,arc=9pt,shadow={2.5pt}{-2.5pt}{0pt}{popink},
      left=6pt,right=6pt,top=9pt,bottom=9pt,halign=center,valign=center,height=1.75cm,width=\linewidth,nobeforeafter]
    {\popdisplay\fontsize{12.5}{13}\selectfont\color{white}\MakeUppercase{#2}}\par\vspace{3pt}
    {\centering\popsemi\fontsize{7.8}{9.5}\selectfont\color{white}#3\par}
  \end{tcolorbox}}
\newcommand{\popcontenu}{%
  \par\noindent{\popxboldls\fontsize{7.5}{7.5}\selectfont\color{popmuted}CE COURS SIGNÉ CONTIENT}\par\vspace{7pt}
  \noindent\begin{minipage}[t]{0.235\linewidth}\poptile{popblue}{Cours}{complet, illustré, pas à pas}\end{minipage}\hfill
  \begin{minipage}[t]{0.235\linewidth}\poptile{poporange}{Méthodes}{et exercices résolus}\end{minipage}\hfill
  \begin{minipage}[t]{0.235\linewidth}\poptile{poppink}{Exercices}{type examen + corrigés}\end{minipage}\hfill
  \begin{minipage}[t]{0.235\linewidth}\poptile{popgreen}{Fiche bilan}{l'essentiel à retenir}\end{minipage}\par}

% Sommaire
\newtcolorbox{sommaire}{enhanced,colback=white,colframe=popink,boxrule=2.2pt,arc=10pt,
  drop shadow=popink,left=18pt,right=18pt,top=13pt,bottom=13pt,before skip=6pt,after skip=20pt,
  before upper={\poppill{Sommaire}{poppurple}{white}\par\vspace{9pt}\popsemi\fontsize{10.6}{14.5}\selectfont\color{popink}}}

% Signature de fin de cours — poussée tout en bas de la dernière page
\newcommand{\findecours}{%
  \par\vfill\nopagebreak
  \begin{center}\popsquiggle{poporange}\end{center}\vspace{4pt}
  \signatureonbuch
}
\AtBeginDocument{\def\llbracket{\mathopen{[\mkern-3.5mu[}}\def\rrbracket{\mathclose{]\mkern-3.5mu]}}} % intervalles d'entiers (glyphe ⟦ absent de la police)
% Glyphes absents de Fira Math : redéfinis à partir de glyphes disponibles
\AtBeginDocument{%
  \def\setminus{\mathbin{\backslash}}\let\smallsetminus\setminus
  \def\vdots{\mathord{\vbox{\baselineskip3.2pt\lineskiplimit0pt\kern4pt\hbox{.}\hbox{.}\hbox{.}}}}%
  \def\ddots{\mathinner{\mkern1mu\raise6.5pt\hbox{.}\mkern2mu\raise3.6pt\hbox{.}\mkern2mu\raise0.7pt\hbox{.}\mkern1mu}}%
  \def\triangle{\mathord{\scalebox{0.9}{$\symup{\Delta}$}}}%
  \def\nabla{\mathord{\raisebox{\depth}{\scalebox{1}[-1]{$\symup{\Delta}$}}}}%
  \def\checkmark{\popcheck{popdark}}%
  \def\star{\mathbin{\ast}}\def\bigstar{\mathord{\ast}}%
  \def\diamond{\mathbin{\scalebox{0.6}{$\Diamond$}}}%
}

%% FIGFIX-BEGIN v4 — correctifs globaux des figures (docs/onbuch-generation/tools/figfix)
\usepackage{adjustbox}\usepackage{etoolbox}
% toute figure TikZ/pgfplots plus large que la ligne est réduite à la largeur disponible
\BeforeBeginEnvironment{tikzpicture}{\begin{adjustbox}{max width=\linewidth}}
\AfterEndEnvironment{tikzpicture}{\end{adjustbox}}
% légendes pgfplots lisibles par-dessus les courbes
\pgfplotsset{every axis/.append style={legend style={fill=white,fill opacity=0.92,text opacity=1,draw=black!15,inner sep=3pt}}}
% étiquettes d'arêtes (syntaxe edge["..."]) avec fond blanc
\tikzset{every edge quotes/.append style={fill=white,inner sep=1.2pt}}
%% FIGFIX-END
\begin{document}

\pagegarde{Applications linéaires du plan}{Mathématiques}{Tle TI}{Configurations et transformations élémentaires du plan}{19}{5 h}{Cette leçon introduit les applications linéaires du plan vectoriel, outil unificateur des transformations étudiées depuis la classe de Première. L'élève apprend à reconnaître une application linéaire, à l'écrire dans une base, à déterminer son noyau et son image, et à caractériser sa bijectivité.
\vspace{4pt}
\begin{multicols}{2}\small
\begin{itemize}
\item À la fin de cette leçon, je sais définir une application linéaire du plan vectoriel et citer des exemples (homothéties, symétries, projections, rotations vectorielles).
\item À la fin de cette leçon, je sais vérifier qu'une application donnée par son écriture analytique est linéaire.
\item À la fin de cette leçon, je sais déterminer l'image d'un vecteur par une application linéaire donnée dans une base.
\item À la fin de cette leçon, je sais déterminer le noyau d'une application linéaire du plan et interpréter son équation caractéristique.
\item À la fin de cette leçon, je sais déterminer l'image d'une application linéaire du plan et en donner une base.
\item À la fin de cette leçon, je sais justifier que le noyau et l'image sont stables pour l'addition et la multiplication par un réel.
\end{itemize}\end{multicols}}

\begin{sommaire}
\begin{multicols}{2}
\begin{enumerate}
\item Définition et premiers exemples d'applications linéaires
\item Endomorphismes, isomorphismes, automorphismes
\item Expression d'une application linéaire dans une base
\item Noyau d'une application linéaire
\item Image d'une application linéaire et stabilité
\item Caractérisation de la bijectivité et synthèse
\item Activité d'intégration
\item Méthodes et exercices résolus
\item Exercices
\item Corrigés des exercices
\item Fiche bilan
\end{enumerate}
\end{multicols}
\end{sommaire}

\begin{objectifs}
\begin{itemize}
\item À la fin de cette leçon, je sais définir une application linéaire du plan vectoriel et citer des exemples (homothéties, symétries, projections, rotations vectorielles).
\item À la fin de cette leçon, je sais vérifier qu'une application donnée par son écriture analytique est linéaire.
\item À la fin de cette leçon, je sais déterminer l'image d'un vecteur par une application linéaire donnée dans une base.
\item À la fin de cette leçon, je sais déterminer le noyau d'une application linéaire du plan et interpréter son équation caractéristique.
\item À la fin de cette leçon, je sais déterminer l'image d'une application linéaire du plan et en donner une base.
\item À la fin de cette leçon, je sais justifier que le noyau et l'image sont stables pour l'addition et la multiplication par un réel.
\item À la fin de cette leçon, je sais décider si une application linéaire du plan est bijective en utilisant son écriture analytique, son noyau ou son image.
\item À la fin de cette leçon, je sais distinguer endomorphisme, isomorphisme et automorphisme.
\end{itemize}
\end{objectifs}

\begin{prerequis}
\begin{itemize}
\item Vecteurs du plan : coordonnées, colinéarité, déterminant de deux vecteurs (Première)
\item Bases du plan vectoriel et décomposition d'un vecteur dans une base (Première)
\item Résolution de systèmes de deux équations linéaires à deux inconnues (Seconde/Première)
\item Notion d'application, d'injectivité, de surjectivité et de bijectivité (début d'année de Terminale)
\item Transformations du plan : symétries, homothéties, projections (Première)
\end{itemize}
\end{prerequis}

\newpage

\coursec{Définition et premiers exemples d'applications linéaires}

\courssub{Situation d'accroche : la machine de sérigraphie de Douala}

À Douala, une entreprise de sérigraphie imprime des logos sur des maillots de football. Pour agrandir, réduire, incliner ou déformer un motif, la machine n'a pas besoin de redessiner le logo : elle applique à \emph{chaque point} du dessin une règle de calcul simple, qui transforme les coordonnées $(x\,;\,y)$ du point en de nouvelles coordonnées $(ax + by\,;\,cx + dy)$, où $a$, $b$, $c$, $d$ sont quatre réglages choisis par le technicien.

Un jour, le technicien constate un phénomène étrange : avec certains réglages, deux motifs différents donnent \emph{exactement la même impression}. Le réglage a « écrasé » le plan : toute une partie de l'information a disparu, impossible de revenir en arrière. Avec d'autres réglages, au contraire, on peut toujours retrouver le motif original à partir de l'impression.

Toute cette leçon naît de trois questions naturelles :
\begin{itemize}
\item Quels réglages \emph{conservent} toute l'information (on dira : quelles applications sont \cle{bijectives}) ?
\item Quels vecteurs sont envoyés sur le vecteur nul (ce sera le \cle{noyau}) ?
\item Quelles images peut-on atteindre (ce sera l'\cle{image}) ?
\end{itemize}

Pour y répondre, il faut d'abord comprendre la propriété remarquable que possèdent toutes les règles de la forme $(x\,;\,y) \mapsto (ax + by\,;\,cx + dy)$ : elles \emph{respectent} l'addition des vecteurs et la multiplication par un réel. C'est précisément ce que nous allons appeler la \cle{linéarité}.

\begin{savaistu}[Les applications linéaires sont partout]
Les logiciels de dessin (Photoshop, Illustrator) et les moteurs de jeux vidéo déplacent, tournent et déforment les personnages à l'écran grâce à des applications linéaires : chaque point de l'image subit une transformation du type $(x\,;\,y) \mapsto (ax + by\,;\,cx + dy)$, calculée des millions de fois par seconde. Les rotations de la \emph{carte} de ton téléphone quand tu tournes l'écran, c'est aussi une application linéaire !
\end{savaistu}

\courssub{Définition d'une application linéaire}

\textbf{Intuition.} Dans le plan vectoriel $E$, on sait faire deux choses : \emph{additionner} des vecteurs et \emph{multiplier} un vecteur par un réel. Une application linéaire est une application qui « commute » avec ces deux opérations : transformer d'abord puis additionner, ou additionner d'abord puis transformer, donne le même résultat. Géométriquement, une application linéaire envoie un parallélogramme sur un parallélogramme (éventuellement aplati).

\begin{definition}[Application linéaire du plan vectoriel]
Soit $E$ le plan vectoriel. On dit qu'une application $f : E \to E$ est une \cle{application linéaire} si elle vérifie les deux conditions suivantes :
\begin{enumerate}
\item \textbf{Additivité} : pour tous vecteurs $\vec{u}$ et $\vec{v}$ de $E$,
\[ f(\vec{u} + \vec{v}) = f(\vec{u}) + f(\vec{v}) \,; \]
\item \textbf{Homogénéité} : pour tout vecteur $\vec{u}$ de $E$ et tout réel $k$,
\[ f(k\,\vec{u}) = k\,f(\vec{u}). \]
\end{enumerate}
Une application linéaire de $E$ dans lui-même est aussi appelée \cle{endomorphisme} de $E$.
\end{definition}

\begin{attention}[Les deux conditions sont indispensables]
Il ne suffit pas de vérifier l'additivité seule, ni l'homogénéité seule : une application peut vérifier l'une sans vérifier l'autre. Par exemple, $f(x\,;\,y) = (x^2\,;\,y)$ vérifie $f(k\vec{u}) = k f(\vec{u})$ pour $k \geq 0$ seulement, et n'est pas additive. Au Baccalauréat, rédige toujours la vérification des \emph{deux} conditions, ou utilise la forme analytique (voir la méthode plus bas).
\end{attention}

La figure suivante illustre la condition d'additivité : l'image d'une somme est la somme des images, autrement dit le parallélogramme est conservé.

\begin{popfigure}
\begin{tikzpicture}[scale=0.85,>=Stealth]
% Plan de départ
\begin{scope}
\draw[->,popmuted] (-0.5,0) -- (4.5,0) node[below left] {};
\draw[->,popmuted] (0,-0.5) -- (0,3.5);
\draw[->,very thick,popblue] (0,0) -- (2.5,0.6) node[midway,below] {$\vec{u}$};
\draw[->,very thick,poporange] (0,0) -- (0.8,2.2) node[midway,left] {$\vec{v}$};
\draw[->,very thick,popink] (0,0) -- (3.3,2.8) node[midway,above right] {$\vec{u}+\vec{v}$};
\draw[dashed,popmuted] (2.5,0.6) -- (3.3,2.8) -- (0.8,2.2);
\fill (0,0) circle (1.5pt) node[below left] {$O$};
\node[popdark] at (2,-1) {Plan de départ};
\end{scope}
% Flèche f
\draw[->,very thick,popdark] (5,1.5) -- (6.6,1.5) node[midway,above] {$f$};
% Plan d'arrivée : f(u)=(2.2,0.2), f(v)=(-0.3,1.8)
\begin{scope}[xshift=7.2cm]
\draw[->,popmuted] (-1.2,0) -- (4,0);
\draw[->,popmuted] (0,-0.5) -- (0,3.5);
\draw[->,very thick,popblue] (0,0) -- (2.2,0.2) node[midway,below] {$f(\vec{u})$};
\draw[->,very thick,poporange] (0,0) -- (-0.3,1.8) node[midway,left] {$f(\vec{v})$};
\draw[->,very thick,popink] (0,0) -- (1.9,2) node[midway,above right] {$f(\vec{u}+\vec{v})$};
\draw[dashed,popmuted] (2.2,0.2) -- (1.9,2) -- (-0.3,1.8);
\fill (0,0) circle (1.5pt) node[below left] {$O$};
\node[popdark] at (1.4,-1) {Plan d'arrivée};
\end{scope}
\end{tikzpicture}
\legende{Additivité d'une application linéaire : l'image du parallélogramme construit sur $\vec{u}$ et $\vec{v}$ est le parallélogramme construit sur $f(\vec{u})$ et $f(\vec{v})$.}
\end{popfigure}

\courssub{Premières conséquences : $f(\vec{0}) = \vec{0}$ et $f(-\vec{u}) = -f(\vec{u})$}

Avant même de donner des exemples, tirons de la définition deux conséquences immédiates mais capitales : une application linéaire envoie \emph{toujours} le vecteur nul sur le vecteur nul, et elle transforme les opposés en opposés.

\begin{propriete}[Conséquences immédiates de la linéarité]
Si $f$ est une application linéaire du plan vectoriel $E$, alors :
\[ f(\vec{0}) = \vec{0}, \qquad f(-\vec{u}) = -f(\vec{u}) \ \text{pour tout } \vec{u} \in E, \qquad f(\vec{u} - \vec{v}) = f(\vec{u}) - f(\vec{v}). \]
\end{propriete}

\begin{experience}[Démonstration guidée (exigible)]
Démontrons ces trois affirmations en n'utilisant que la définition.
\begin{enumerate}
\item \textbf{Montrons que $f(\vec{0}) = \vec{0}$.} L'astuce consiste à écrire $\vec{0} = 0 \cdot \vec{u}$ pour un vecteur $\vec{u}$ quelconque. Par homogénéité avec $k = 0$ :
\[ f(\vec{0}) = f(0 \cdot \vec{u}) = 0 \cdot f(\vec{u}) = \vec{0}. \]
\item \textbf{Montrons que $f(-\vec{u}) = -f(\vec{u})$.} On écrit $-\vec{u} = (-1)\,\vec{u}$. Par homogénéité avec $k = -1$ :
\[ f(-\vec{u}) = f\big((-1)\,\vec{u}\big) = (-1)\,f(\vec{u}) = -f(\vec{u}). \]
\item \textbf{Montrons que $f(\vec{u} - \vec{v}) = f(\vec{u}) - f(\vec{v})$.} On écrit $\vec{u} - \vec{v} = \vec{u} + (-\vec{v})$, puis on applique l'additivité et le résultat précédent :
\[ f(\vec{u} - \vec{v}) = f(\vec{u}) + f(-\vec{v}) = f(\vec{u}) - f(\vec{v}). \]
\end{enumerate}
Retiens l'astuce de départ : \emph{écrire le vecteur nul comme $0 \cdot \vec{u}$}. C'est le point de départ de presque toutes les démonstrations sur les applications linéaires.
\end{experience}

\begin{aretenir}[Un test rapide de non-linéarité]
Si $f(\vec{0}) \neq \vec{0}$, alors $f$ \textbf{n'est pas linéaire}. C'est un critère d'élimination très rapide : dès qu'une application envoie l'origine ailleurs qu'à l'origine, inutile d'aller plus loin. Attention à la réciproque : $f(\vec{0}) = \vec{0}$ ne suffit \emph{pas} à prouver la linéarité (par exemple $f(x\,;\,y) = (x^2\,;\,y)$ vérifie $f(\vec{0}) = \vec{0}$ sans être linéaire).
\end{aretenir}

\courssub{Exemples fondamentaux d'applications linéaires}

Les transformations vectorielles que tu as rencontrées dans les classes précédentes sont presque toutes linéaires. Vérifions-le sur quatre exemples à connaître par cœur.

\textbf{1. L'homothétie vectorielle de rapport $k$.} C'est l'application $h_k : \vec{u} \mapsto k\,\vec{u}$. Pour tous $\vec{u}, \vec{v}$ et tout réel $\lambda$ :
\[ h_k(\vec{u} + \vec{v}) = k(\vec{u} + \vec{v}) = k\vec{u} + k\vec{v} = h_k(\vec{u}) + h_k(\vec{v}), \]
\[ h_k(\lambda \vec{u}) = k(\lambda \vec{u}) = \lambda (k \vec{u}) = \lambda\, h_k(\vec{u}). \]
Donc $h_k$ est linéaire. C'est le réglage « agrandir / réduire » de la machine de sérigraphie. Le cas $k = 1$ donne l'\emph{identité} $\mathrm{id}_E$, et $k = 0$ donne l'\emph{application nulle}, toutes deux linéaires.

\textbf{2. La symétrie vectorielle.} Étant donné une droite vectorielle $D$ et une droite $\Delta$ supplémentaire ($E = D \oplus \Delta$), tout vecteur $\vec{u}$ s'écrit de manière unique $\vec{u} = \vec{u}_1 + \vec{u}_2$ avec $\vec{u}_1 \in D$ et $\vec{u}_2 \in \Delta$. La symétrie par rapport à $D$ parallèlement à $\Delta$ est $s(\vec{u}) = \vec{u}_1 - \vec{u}_2$. Sa linéarité se vérifie en décomposant $\vec{u} + \vec{v}$ et $k\vec{u}$ : elle est linéaire.

\textbf{3. La projection vectorielle sur une droite.} Avec les mêmes notations, la projection sur $D$ parallèlement à $\Delta$ est $p(\vec{u}) = \vec{u}_1$. Si $\vec{u} = \vec{u}_1 + \vec{u}_2$ et $\vec{v} = \vec{v}_1 + \vec{v}_2$, alors $\vec{u} + \vec{v} = (\vec{u}_1 + \vec{v}_1) + (\vec{u}_2 + \vec{v}_2)$, donc
\[ p(\vec{u} + \vec{v}) = \vec{u}_1 + \vec{v}_1 = p(\vec{u}) + p(\vec{v}), \]
et de même $p(k\vec{u}) = k\,\vec{u}_1 = k\,p(\vec{u})$. La projection est linéaire. Remarque au passage : tous les vecteurs de $\Delta$ sont envoyés sur $\vec{0}$ : c'est le réglage qui « écrase » le plan de notre situation d'accroche !

\begin{popfigure}
\begin{tikzpicture}[scale=1,>=Stealth]
% Droite D (horizontale) et Delta (oblique)
\draw[very thick,popblue] (-0.5,0) -- (6.5,0) node[right] {$D$};
\draw[very thick,poporange] (2.2,-0.8) -- (4.6,3.6) node[above] {$\Delta$};
\fill (3,1) circle (0pt);
\fill (0,0) circle (1.5pt) node[below left] {$O$};
% Vecteur u
\coordinate (U) at (5,2.4);
\draw[->,very thick,popink] (0,0) -- (U) node[above right] {$\vec{u}$};
% Projection parallèle à Delta : direction (0.6,0.9) normalisée ~ pente 3/2
\draw[dashed,popmuted] (U) -- (3.4,0);
\draw[->,very thick,popgreen] (0,0) -- (3.4,0) node[midway,below] {$p(\vec{u})$};
% Vecteur v
\coordinate (V) at (1.6,1.5);
\draw[->,very thick,popink!60] (0,0) -- (V) node[above left] {$\vec{v}$};
\draw[dashed,popmuted] (V) -- (0.6,0);
\draw[->,very thick,popgreen!70!black] (0,0) -- (0.6,0) node[midway,below left] {$p(\vec{v})$};
\end{tikzpicture}
\legende{Projection vectorielle sur la droite $D$ parallèlement à la droite $\Delta$ : on « glisse » le long des droites parallèles à $\Delta$ jusqu'à $D$.}
\end{popfigure}

\textbf{4. La rotation vectorielle d'angle $\theta$.} La rotation $r_\theta$ fait tourner chaque vecteur d'un angle $\theta$ autour de l'origine, en conservant les normes. Comme elle conserve les parallélogrammes (elle envoie le parallélogramme construit sur $\vec{u}, \vec{v}$ sur celui construit sur $r_\theta(\vec{u}), r_\theta(\vec{v})$) et que $r_\theta(k\vec{u}) = k\,r_\theta(\vec{u})$, elle est linéaire. Nous verrons dans la section 3 son écriture analytique :
\[ r_\theta(x\,;\,y) = \big(x\cos\theta - y\sin\theta\,;\ x\sin\theta + y\cos\theta\big). \]

\courssub{Contre-exemples : ce qui n'est PAS linéaire}

Savoir reconnaître ce qui n'est pas linéaire est aussi important que l'inverse. Voici trois pièges classiques.

\textbf{1. Les translations.} Une translation de vecteur $\vec{a} \neq \vec{0}$ envoie $\vec{0}$ sur $\vec{a} \neq \vec{0}$. Or une application linéaire vérifie toujours $f(\vec{0}) = \vec{0}$. Donc \emph{aucune translation non triviale n'est linéaire}.

\textbf{2. Les applications affines avec terme constant.} Considérons $f(x\,;\,y) = (2x + 1\,;\,y - 3)$. On a $f(0\,;\,0) = (1\,;\,-3) \neq (0\,;\,0)$, donc $f$ n'est pas linéaire. Les termes constants $+1$ et $-3$ correspondent en fait à une translation.

\textbf{3. Les applications non homogènes en degré.} Considérons $f(x\,;\,y) = (x^2\,;\,y)$. On a bien $f(\vec{0}) = \vec{0}$, mais testons l'homogénéité avec $k = 2$ et $\vec{u} = (1\,;\,0)$ :
\[ f(2\vec{u}) = f(2\,;\,0) = (4\,;\,0) \quad \text{alors que} \quad 2f(\vec{u}) = 2(1\,;\,0) = (2\,;\,0). \]
Comme $(4\,;\,0) \neq (2\,;\,0)$, l'homogénéité échoue : $f$ n'est pas linéaire.

\begin{attention}[L'erreur la plus fréquente au Bac]
Beaucoup d'élèves affirment que $(x\,;\,y) \mapsto (x + 1\,;\,2y)$ est linéaire « parce qu'il n'y a que des $x$ et des $y$ ». C'est faux ! Le terme constant $+1$ l'interdit : $f(0\,;\,0) = (1\,;\,0) \neq \vec{0}$. \textbf{Réflexe à acquérir} : avant toute chose, calculer l'image de $(0\,;\,0)$. Si ce n'est pas $(0\,;\,0)$, l'application n'est pas linéaire, point final.
\end{attention}

\courssub{Méthode : reconnaître la linéarité à partir de l'écriture analytique}

Revenons à la machine de sérigraphie : ses réglages produisent des applications de la forme $(x\,;\,y) \mapsto (ax + by\,;\,cx + dy)$. La bonne nouvelle, c'est que \emph{toutes} ces applications sont linéaires, et ce sont (nous le verrons en section 3) les seules.

\begin{propriete}[Forme analytique des applications linéaires du plan]
Toute application $f : \mathbb{R}^2 \to \mathbb{R}^2$ de la forme
\[ f(x\,;\,y) = (ax + by\,;\ cx + dy), \quad \text{où } a, b, c, d \in \mathbb{R}, \]
est une application linéaire. Réciproquement, toute application linéaire du plan admet une écriture de cette forme dans une base donnée.
\end{propriete}

\begin{experience}[Démonstration guidée : la forme analytique implique la linéarité]
Soient $\vec{u} = (x\,;\,y)$, $\vec{v} = (x'\,;\,y')$ et $k \in \mathbb{R}$.
\begin{enumerate}
\item \textbf{Additivité.} On a $\vec{u} + \vec{v} = (x + x'\,;\,y + y')$, donc
\begin{align*}
f(\vec{u} + \vec{v}) &= \big(a(x + x') + b(y + y')\,;\ c(x + x') + d(y + y')\big) \\
&= \big(ax + by + ax' + by'\,;\ cx + dy + cx' + dy'\big) \\
&= (ax + by\,;\ cx + dy) + (ax' + by'\,;\ cx' + dy') \\
&= f(\vec{u}) + f(\vec{v}).
\end{align*}
\item \textbf{Homogénéité.} On a $k\vec{u} = (kx\,;\,ky)$, donc
\[ f(k\vec{u}) = \big(a(kx) + b(ky)\,;\ c(kx) + d(ky)\big) = k\,(ax + by\,;\ cx + dy) = k\,f(\vec{u}). \]
\end{enumerate}
Les deux conditions sont vérifiées : $f$ est linéaire. $\square$
\end{experience}

\begin{methode}[Prouver qu'une application est linéaire]
Deux stratégies au choix :
\begin{enumerate}
\item \textbf{Méthode analytique (la plus rapide)} : si l'application est donnée par une formule, vérifier qu'elle s'écrit $f(x\,;\,y) = (ax + by\,;\,cx + dy)$, c'est-à-dire que chaque coordonnée est une combinaison linéaire de $x$ et $y$ \emph{sans terme constant, sans carré, sans produit $xy$, sans valeur absolue}. Si oui, conclure directement : « $f$ est linéaire d'après la propriété du cours ».
\item \textbf{Méthode par la définition} : si l'application est définie géométriquement (projection, symétrie, rotation...), poser $\vec{u} = (x\,;\,y)$, $\vec{v} = (x'\,;\,y')$, $k \in \mathbb{R}$, et vérifier les deux conditions $f(\vec{u} + \vec{v}) = f(\vec{u}) + f(\vec{v})$ et $f(k\vec{u}) = k f(\vec{u})$.
\item \textbf{Pour prouver la NON-linéarité} : calculer $f(0\,;\,0)$ ; si ce n'est pas $(0\,;\,0)$, c'est terminé. Sinon, exhiber un contre-exemple à l'homogénéité ou à l'additivité avec des vecteurs bien choisis.
\end{enumerate}
\end{methode}

\begin{exoresolu}[Étude complète d'un premier exemple]
On considère l'application $f : \mathbb{R}^2 \to \mathbb{R}^2$ définie par
\[ f(x\,;\,y) = (2x - y\,;\ x + 3y). \]
\begin{enumerate}
\item Montrer que $f$ est une application linéaire.
\item Calculer l'image du vecteur $\vec{u} = (1\,;\,-2)$.
\item Calculer $f(3\vec{u})$ de deux façons différentes.
\end{enumerate}
\tcblower
\textbf{Solution détaillée.}
\begin{enumerate}
\item Chaque coordonnée de $f(x\,;\,y)$ est une combinaison linéaire de $x$ et $y$ sans terme constant : $2x - y$ et $x + 3y$. L'application est bien de la forme $(ax + by\,;\,cx + dy)$ avec $a = 2$, $b = -1$, $c = 1$, $d = 3$. Donc $f$ est \textbf{linéaire}.

\emph{Vérification par la définition (pour s'entraîner)} : soient $\vec{u} = (x\,;\,y)$, $\vec{v} = (x'\,;\,y')$.
\begin{align*}
f(\vec{u} + \vec{v}) &= \big(2(x + x') - (y + y')\,;\ (x + x') + 3(y + y')\big) \\
&= (2x - y\,;\ x + 3y) + (2x' - y'\,;\ x' + 3y') = f(\vec{u}) + f(\vec{v}),
\end{align*}
et $f(k\vec{u}) = (2kx - ky\,;\ kx + 3ky) = k\,(2x - y\,;\ x + 3y) = k f(\vec{u})$. Les deux conditions sont confirmées.

\item On remplace $x$ par $1$ et $y$ par $-2$ :
\[ f(1\,;\,-2) = \big(2(1) - (-2)\,;\ 1 + 3(-2)\big) = (4\,;\,-5). \]

\item \textbf{Première façon} : $3\vec{u} = (3\,;\,-6)$, donc
\[ f(3\vec{u}) = f(3\,;\,-6) = \big(2(3) - (-6)\,;\ 3 + 3(-6)\big) = (12\,;\,-15). \]
\textbf{Deuxième façon} : par homogénéité,
\[ f(3\vec{u}) = 3 f(\vec{u}) = 3\,(4\,;\,-5) = (12\,;\,-15). \]
Les deux méthodes donnent le même résultat, ce qui illustre concrètement la propriété $f(k\vec{u}) = k f(\vec{u})$.
\end{enumerate}
\end{exoresolu}

\begin{exercice}[À toi de jouer]
Pour chacune des applications suivantes, dire si elle est linéaire en justifiant :
\begin{enumerate}
\item $f_1(x\,;\,y) = (3x + 2y\,;\ x - y)$ ;
\item $f_2(x\,;\,y) = (x + 2\,;\,3y)$ ;
\item $f_3(x\,;\,y) = (xy\,;\,x + y)$ ;
\item $f_4(x\,;\,y) = (x - y\,;\,0)$.
\end{enumerate}
\end{exercice}

\begin{corrige}[Exercice : À toi de jouer]
\begin{enumerate}
\item $f_1$ est de la forme $(ax + by\,;\,cx + dy)$ avec $a = 3$, $b = 2$, $c = 1$, $d = -1$ : elle est \textbf{linéaire}.
\item $f_2(0\,;\,0) = (2\,;\,0) \neq (0\,;\,0)$ : $f_2$ \textbf{n'est pas linéaire} (terme constant).
\item $f_3(0\,;\,0) = (0\,;\,0)$, mais testons l'homogénéité : $f_3(2 \cdot (1\,;\,1)) = f_3(2\,;\,2) = (4\,;\,4)$ alors que $2 f_3(1\,;\,1) = 2(1\,;\,2) = (2\,;\,4)$. Comme $(4\,;\,4) \neq (2\,;\,4)$, $f_3$ \textbf{n'est pas linéaire} (à cause du produit $xy$).
\item $f_4$ est de la forme $(ax + by\,;\,cx + dy)$ avec $a = 1$, $b = -1$, $c = 0$, $d = 0$ : elle est \textbf{linéaire}. C'est d'ailleurs (à un changement de repère près) une projection : tous les vecteurs de la forme $(t\,;\,t)$ sont envoyés sur $\vec{0}$, ce qui annonce la notion de noyau.
\end{enumerate}
\end{corrige}

\begin{aretenir}[L'essentiel de cette section]
\begin{itemize}
\item Une application $f : E \to E$ est \cle{linéaire} si $f(\vec{u} + \vec{v}) = f(\vec{u}) + f(\vec{v})$ et $f(k\vec{u}) = k f(\vec{u})$ pour tous $\vec{u}, \vec{v} \in E$ et tout $k \in \mathbb{R}$.
\item Toute application linéaire vérifie $f(\vec{0}) = \vec{0}$ et $f(-\vec{u}) = -f(\vec{u})$.
\item Sont linéaires : homothéties, symétries, projections, rotations vectorielles. Ne sont \emph{pas} linéaires : translations, applications avec terme constant, applications avec des carrés ou des produits de coordonnées.
\item Critère analytique : $f$ est linéaire si et seulement si $f(x\,;\,y) = (ax + by\,;\,cx + dy)$ avec $a, b, c, d \in \mathbb{R}$.
\end{itemize}
Dans la prochaine section, nous distinguerons parmi ces applications celles qui sont bijectives : les \emph{automorphismes}, seuls réglages de la machine qui ne perdent aucune information.
\end{aretenir}

\coursec{Endomorphismes, isomorphismes, automorphismes}

Dans la section précédente, nous avons défini les applications linéaires du plan vectoriel et découvert nos premiers exemples : homothéties, projections, symétries, rotations vectorielles. Toutes ces applications partagent une propriété remarquable : elles \emph{respectent} l'addition des vecteurs et la multiplication par un réel. Mais elles ne se ressemblent pas toutes ! Une homothétie de rapport $2$ \og étire \fg{} le plan sans rien perdre, tandis qu'une projection \og écrase \fg{} tout le plan sur une droite : une infinité de vecteurs différents ont alors la même image. Cette différence de comportement est fondamentale, et elle motive le vocabulaire que nous allons construire dans cette section : \cle{endomorphisme}, \cle{isomorphisme}, \cle{automorphisme}. Ces trois mots, qui effraient parfois au premier abord, ne font que traduire précisément deux questions très simples : \emph{où arrive-t-on ?} et \emph{perd-on de l'information en chemin ?}

\courssub{Endomorphismes : des applications linéaires du plan dans lui-même}

Commençons par une image concrète. Un opérateur de CAMTEL qui traite un signal ne change pas la \emph{nature} du signal : il prend un signal en entrée et produit un signal en sortie. De même, toutes les applications linéaires que nous avons rencontrées prennent un vecteur du plan et produisent \emph{un vecteur du même plan}. On dit qu'elles agissent \og en interne \fg{}.

\begin{definition}[Endomorphisme]
On appelle \cle{endomorphisme} du plan vectoriel $\mathcal{P}$ toute \textbf{application linéaire de $\mathcal{P}$ dans lui-même}, c'est-à-dire toute application $f : \mathcal{P} \to \mathcal{P}$ vérifiant, pour tous vecteurs $\vec{u}, \vec{v}$ de $\mathcal{P}$ et tout réel $\lambda$ :
\[ f(\vec{u} + \vec{v}) = f(\vec{u}) + f(\vec{v}) \qquad \text{et} \qquad f(\lambda \vec{u}) = \lambda\, f(\vec{u}). \]
Le préfixe \emph{endo-} vient du grec \emph{endon}, \og à l'intérieur \fg{} : l'application reste à l'intérieur du même espace.
\end{definition}

\begin{aretenir}[Tous nos exemples sont des endomorphismes]
Puisque nous travaillons dans cette leçon avec des applications linéaires \emph{du plan vectoriel dans lui-même}, \textbf{toutes} les applications linéaires étudiées ici sont des endomorphismes du plan : l'application nulle, l'identité, les homothéties vectorielles, les projections, les symétries, les rotations vectorielles. Le mot \og endomorphisme \fg{} ne dit donc rien de plus que \og application linéaire du plan dans le plan \fg{}. La vraie distinction viendra de la bijectivité.
\end{aretenir}

\courssub{Rappel : injectivité, surjectivité, bijectivité}

Avant de définir les isomorphismes, rappelons soigneusement trois notions vues en classe de Première, car elles sont au cœur de toute la section.

\begin{definition}[Injectivité, surjectivité, bijectivité]
Soit $f : E \to F$ une application.
\begin{itemize}
\item $f$ est \cle{injective} si deux éléments distincts de $E$ ont toujours des images distinctes, autrement dit :
\[ f(\vec{u}) = f(\vec{v}) \implies \vec{u} = \vec{v}. \]
Aucun vecteur de $F$ n'est atteint \og deux fois \fg{}.
\item $f$ est \cle{surjective} si tout élément de $F$ possède au moins un antécédent dans $E$ :
\[ \forall \vec{w} \in F,\ \exists\, \vec{u} \in E,\ f(\vec{u}) = \vec{w}. \]
Tout l'espace d'arrivée est atteint.
\item $f$ est \cle{bijective} si elle est à la fois injective et surjective : chaque élément de $F$ possède alors \textbf{exactement un} antécédent. On peut alors définir la \cle{bijection réciproque} $f^{-1} : F \to E$.
\end{itemize}
\end{definition}

Pour une application linéaire, ces notions prennent un sens géométrique très parlant :
\begin{itemize}
\item \textbf{Injective} signifie : $f$ n'\emph{écrase} rien. Deux vecteurs différents ne sont jamais \og collés \fg{} l'un sur l'autre.
\item \textbf{Surjective} signifie : $f$ \emph{atteint tout} le plan d'arrivée. Aucune zone n'est oubliée.
\item \textbf{Bijective} signifie : $f$ est un \emph{relooking parfaitement réversible} du plan. On peut toujours revenir en arrière de façon unique.
\end{itemize}

\begin{exemplebox}[Testons nos exemples de la section 1]
\begin{itemize}
\item L'\textbf{homothétie} $h$ de rapport $k \neq 0$ : si $h(\vec{u}) = h(\vec{v})$, alors $k\vec{u} = k\vec{v}$, donc $\vec{u} = \vec{v}$ (on divise par $k \neq 0$) : elle est injective. Tout vecteur $\vec{w}$ a pour antécédent $\dfrac{1}{k}\vec{w}$ : elle est surjective. C'est une application \textbf{bijective}.
\item La \textbf{projection} $p$ sur l'axe des abscisses, parallèlement à l'axe des ordonnées : les vecteurs $(1\,;\,0)$ et $(1\,;\,5)$ ont la même image $(1\,;\,0)$, donc $p$ n'est \textbf{pas injective}. De plus, aucun vecteur d'ordonnée non nulle n'est atteint : $p$ n'est \textbf{pas surjective}.
\item L'\textbf{application nulle} envoie tout le plan sur $\vec{0}$ : ni injective, ni surjective.
\end{itemize}
\end{exemplebox}

\courssub{Isomorphismes et automorphismes}

\begin{definition}[Isomorphisme]
On appelle \cle{isomorphisme} toute application linéaire \textbf{bijective} entre deux espaces vectoriels. Le préfixe \emph{iso-} signifie \og égal, identique \fg{} : un isomorphisme établit une correspondance parfaite entre les deux espaces, qui ont alors exactement la \emph{même structure}.
\end{definition}

\begin{definition}[Automorphisme]
On appelle \cle{automorphisme} du plan vectoriel tout \textbf{endomorphisme bijectif}, c'est-à-dire toute application linéaire bijective du plan \emph{dans lui-même}. Le préfixe \emph{auto-} signifie \og soi-même \fg{}.
\end{definition}

\begin{aretenir}[La hiérarchie à retenir]
\[ \text{Automorphisme} = \text{endomorphisme} + \text{bijectivité} = \text{isomorphisme d'un espace sur lui-même}. \]
En particulier : tout automorphisme est un isomorphisme, mais la notion d'automorphisme exige en plus que l'espace de départ et l'espace d'arrivée soient \textbf{le même}.
\end{aretenir}

\begin{attention}[Ne pas confondre isomorphisme et automorphisme]
L'erreur la plus fréquente au Baccalauréat est d'employer les deux mots indifféremment. Retiens bien :
\begin{itemize}
\item Un \textbf{isomorphisme} peut relier deux espaces \emph{différents} (par exemple le plan vectoriel et $\mathbb{R}^2$).
\item Un \textbf{automorphisme} va toujours d'un espace \emph{sur lui-même}.
\end{itemize}
Dans le plan, toute application linéaire bijective $f : \mathcal{P} \to \mathcal{P}$ est à la fois un isomorphisme \emph{et} un automorphisme ; mais si l'énoncé précise \og de $\mathcal{P}$ dans $\mathcal{P}$ \fg{}, le mot attendu est \emph{automorphisme}. Autre piège : une homothétie de rapport $k = 0$ est l'application nulle, qui n'est \textbf{pas} bijective. La condition $k \neq 0$ est indispensable pour parler d'automorphisme !
\end{attention}

\begin{exemplebox}[Deux situations typiques dans le plan]
\begin{itemize}
\item \textbf{Automorphisme :} l'homothétie vectorielle de rapport $k \neq 0$, la symétrie centrale, toute rotation vectorielle, toute symétrie axiale vectorielle. Chacune transforme le plan en lui-même sans perte d'information.
\item \textbf{Endomorphisme non bijectif :} toute projection sur une droite. Elle envoie tout le plan sur une droite : une infinité de vecteurs ont la même image, et presque tout le plan n'est pas atteint. L'information est partiellement détruite.
\end{itemize}
\end{exemplebox}

\begin{popfigure}
\begin{tikzpicture}[scale=0.85, >=Stealth]
% ---- Colonne gauche : automorphisme ----
\begin{scope}
\draw[popline, ->] (-2.3,0) -- (2.3,0) node[below left, font=\small] {$x$};
\draw[popline, ->] (0,-1.6) -- (0,2.2) node[below left, font=\small] {$y$};
\foreach \p/\c in {(1,0.5)/popblue, (-1,1)/poporange, (0.5,-0.8)/popgreen, (-0.8,-0.6)/poppurple, (1.4,1.4)/popgold}{
  \draw[\c, very thick, ->] (0,0) -- \p;
}
\node[popdark, font=\small\bfseries] at (0,-2.1) {Automorphisme};
\node[popmuted, font=\small, text width=4.2cm, align=center] at (0,-2.9) {tout le plan est atteint,\\ aucun vecteur n'est écrasé};
\end{scope}
% ---- Colonne droite : projection ----
\begin{scope}[xshift=8cm]
\draw[popline, ->] (-2.3,0) -- (2.3,0) node[below left, font=\small] {$x$};
\draw[popline, ->] (0,-1.6) -- (0,2.2) node[below left, font=\small] {$y$};
\draw[popink, very thick] (-2.3,0) -- (2.3,0);
\foreach \x/\y in {1/1.5, -1.2/0.9, 0.6/-1.1, 1.7/0.6}{
  \draw[poporange, thick, dashed, ->] (\x,\y) -- (\x,0);
  \fill[popblue] (\x,\y) circle (1.6pt);
  \fill[popink] (\x,0) circle (1.6pt);
}
\node[popdark, font=\small\bfseries] at (0,-2.1) {Projection (endomorphisme non bijectif)};
\node[popmuted, font=\small, text width=4.2cm, align=center] at (0,-2.9) {tout le plan est écrasé\\ sur une seule droite};
\end{scope}
\end{tikzpicture}
\legende{À gauche, un automorphisme : les vecteurs du plan sont transformés mais couvrent encore tout le plan, et deux vecteurs distincts restent distincts. À droite, une projection sur l'axe des abscisses : des vecteurs différents (points bleus) tombent sur la même image (points foncés) ; l'application n'est ni injective ni surjective.}
\end{popfigure}

\courssub{Exemple chiffré complet : un automorphisme et sa réciproque}

Pour montrer qu'une application linéaire du plan est bijective, la méthode la plus directe à ce stade consiste à \textbf{résoudre l'équation $f(x\,;\,y) = (x'\,;\,y')$} et à vérifier qu'elle admet, pour tout $(x'\,;\,y')$, une solution unique.

\begin{exoresolu}[$f(x\,;\,y) = (x - y\,;\,x + y)$ est un automorphisme]
On considère l'application $f$ du plan vectoriel dans lui-même définie par
\[ f(x\,;\,y) = (x - y\,;\,x + y). \]
\begin{enumerate}
\item Vérifier que $f$ est un endomorphisme du plan.
\item Montrer que $f$ est bijective et déterminer sa bijection réciproque.
\item Conclure.
\end{enumerate}
\tcblower
\textbf{1. Linéarité.} Soient $\vec{u} = (x\,;\,y)$, $\vec{v} = (x'\,;\,y')$ et $\lambda \in \mathbb{R}$.
\begin{align*}
f(\vec{u} + \vec{v}) &= f(x + x'\,;\,y + y') = \big((x+x') - (y+y')\,;\,(x+x') + (y+y')\big) \\
&= (x - y\,;\,x + y) + (x' - y'\,;\,x' + y') = f(\vec{u}) + f(\vec{v}).
\end{align*}
\begin{align*}
f(\lambda \vec{u}) &= f(\lambda x\,;\,\lambda y) = (\lambda x - \lambda y\,;\,\lambda x + \lambda y) = \lambda\,(x - y\,;\,x + y) = \lambda f(\vec{u}).
\end{align*}
$f$ est donc linéaire, de $\mathcal{P}$ dans $\mathcal{P}$ : c'est un \textbf{endomorphisme}.

\textbf{2. Bijectivité.} Soit $(x'\,;\,y')$ un vecteur quelconque du plan. Résolvons $f(x\,;\,y) = (x'\,;\,y')$ :
\[ \begin{cases} x - y = x' \\ x + y = y' \end{cases} \iff \begin{cases} 2x = x' + y' \\ 2y = y' - x' \end{cases} \iff \begin{cases} x = \dfrac{x' + y'}{2} \\[6pt] y = \dfrac{y' - x'}{2} \end{cases} \]
Pour \emph{tout} vecteur $(x'\,;\,y')$, le système admet une solution, et cette solution est \emph{unique}. Donc $f$ est à la fois surjective et injective : $f$ est \textbf{bijective}, et sa réciproque est
\[ f^{-1}(x'\,;\,y') = \left(\frac{x' + y'}{2}\,;\,\frac{y' - x'}{2}\right). \]
\textbf{Vérification numérique :} $f(3\,;\,1) = (3-1\,;\,3+1) = (2\,;\,4)$, et en appliquant la réciproque : $f^{-1}(2\,;\,4) = \left(\dfrac{2+4}{2}\,;\,\dfrac{4-2}{2}\right) = (3\,;\,1)$. On retombe bien sur le vecteur de départ.

\textbf{3. Conclusion.} $f$ est un endomorphisme bijectif du plan : c'est un \textbf{automorphisme} du plan vectoriel.
\end{exoresolu}

\begin{methode}[Prouver qu'une application linéaire du plan est un automorphisme]
Pour montrer que $f(x\,;\,y) = (ax + by\,;\,cx + dy)$ est un automorphisme :
\begin{enumerate}
\item Justifier la linéarité (souvent admise par la forme analytique : coordonnées de la forme $ax + by$).
\item Poser l'équation $f(x\,;\,y) = (x'\,;\,y')$, c'est-à-dire le système $\begin{cases} ax + by = x' \\ cx + dy = y' \end{cases}$.
\item Montrer que ce système admet, pour tout $(x'\,;\,y')$, une \textbf{unique} solution (c'est le cas lorsque $ad - bc \neq 0$, fait que nous retrouverons avec l'écriture analytique).
\item Conclure : $f$ est un endomorphisme bijectif, donc un automorphisme, et donner éventuellement $f^{-1}$.
\end{enumerate}
\end{methode}

\courssub{Composée de deux applications linéaires}

Terminons cette section par un résultat qui sera un outil précieux dans toute la suite de la leçon, notamment pour l'écriture analytique des applications linéaires.

\begin{propriete}[Composée de deux applications linéaires (admise)]
Soient $f$ et $g$ deux applications linéaires du plan vectoriel dans lui-même. Alors l'application composée $g \circ f$, définie par
\[ (g \circ f)(\vec{u}) = g\big(f(\vec{u})\big), \]
est encore une \textbf{application linéaire}. Autrement dit, la composée de deux endomorphismes du plan est un endomorphisme du plan.
\end{propriete}

Ce résultat est admis au programme de Terminale TI, mais sa démonstration est très courte et formatrice : pour tous vecteurs $\vec{u}, \vec{v}$ et tout réel $\lambda$, en utilisant d'abord la linéarité de $f$ puis celle de $g$ :
\begin{align*}
(g \circ f)(\vec{u} + \vec{v}) &= g\big(f(\vec{u} + \vec{v})\big) = g\big(f(\vec{u}) + f(\vec{v})\big) = g(f(\vec{u})) + g(f(\vec{v})), \\
(g \circ f)(\lambda \vec{u}) &= g\big(f(\lambda \vec{u})\big) = g\big(\lambda f(\vec{u})\big) = \lambda\, g(f(\vec{u})).
\end{align*}
La linéarité se \og transmet \fg{} donc de $f$ et $g$ à leur composée.

\begin{exemplebox}[Composer une homothétie et une projection]
Soient $h(x\,;\,y) = (2x\,;\,2y)$ (homothétie de rapport $2$) et $p(x\,;\,y) = (x\,;\,0)$ (projection sur l'axe des abscisses). Alors :
\[ (p \circ h)(x\,;\,y) = p(2x\,;\,2y) = (2x\,;\,0), \qquad (h \circ p)(x\,;\,y) = h(x\,;\,0) = (2x\,;\,0). \]
Les deux composées sont bien linéaires. Remarque : ici $p \circ h = h \circ p$, mais ce n'est \textbf{pas} une règle générale ! En général, $g \circ f \neq f \circ g$.
\end{exemplebox}

\begin{attention}[La composée n'est pas commutative en général]
Contrairement à la multiplication des nombres réels, l'ordre de composition compte : $g \circ f$ signifie \og on applique d'abord $f$, puis $g$ \fg{}. Prends l'habitude de lire les composées \emph{de droite à gauche}. De plus, la composée de deux automorphismes est un automorphisme, mais la composée d'un automorphisme et d'une projection n'est jamais bijective : dès qu'une application \og perd de l'information \fg{} dans la chaîne, la perte est irrémédiable.
\end{attention}

\begin{savaistu}[Les automorphismes au cœur de la cryptographie]
Pourquoi les mathématiciens s'intéressent-ils tant aux automorphismes ? Parce qu'ils sont \textbf{réversibles} ! En cryptographie — la science des messages secrets — on chiffre un message en lui appliquant une transformation que seul le destinataire peut inverser. Un automorphisme est parfait pour cela : comme il est bijectif, sa réciproque existe et permet de déchiffrer sans ambiguïté. À l'inverse, une projection serait un désastre : plusieurs messages différents donneraient le même message chiffré, impossible à décoder ! Le célèbre chiffre de Hill (1929) chiffre les lettres par blocs à l'aide d'automorphismes : chaque paire de lettres devient un vecteur, transformé par une application linéaire bijective. Aujourd'hui, ces principes protègent tes transactions Mobile Money et tes messages : derrière chaque transfert MTN MoMo sécurisé se cachent des transformations réversibles... donc des automorphismes !
\end{savaistu}

\begin{exercice}[À toi de jouer]
On considère l'application $g$ du plan vectoriel dans lui-même définie par $g(x\,;\,y) = (x + 2y\,;\,2x + 4y)$.
\begin{enumerate}
\item Justifier que $g$ est un endomorphisme du plan.
\item Trouver deux vecteurs distincts ayant la même image par $g$. Que peut-on en conclure ?
\item Le vecteur $(0\,;\,1)$ possède-t-il un antécédent par $g$ ? Conclure sur la nature de $g$.
\end{enumerate}
\end{exercice}

\begin{corrige}[Exercice : l'application $g$]
\textbf{1.} Les coordonnées de $g(x\,;\,y)$ sont de la forme $ax + by$ : $g$ est une application linéaire du plan dans lui-même, donc un endomorphisme (on peut aussi vérifier directement les deux identités de linéarité).

\textbf{2.} $g(2\,;\,-1) = (2 - 2\,;\,4 - 4) = (0\,;\,0)$ et $g(0\,;\,0) = (0\,;\,0)$. Deux vecteurs distincts ont la même image : $g$ n'est \textbf{pas injective}.

\textbf{3.} On cherche $(x\,;\,y)$ tel que $g(x\,;\,y) = (0\,;\,1)$, c'est-à-dire $\begin{cases} x + 2y = 0 \\ 2x + 4y = 1 \end{cases}$. Or la deuxième équation s'écrit $2(x + 2y) = 1$, soit $0 = 1$ : impossible. Le vecteur $(0\,;\,1)$ n'a pas d'antécédent : $g$ n'est \textbf{pas surjective}.

\textbf{Conclusion :} $g$ est un endomorphisme \textbf{non bijectif} du plan (toutes ses images sont de la forme $(t\,;\,2t)$ : le plan est écrasé sur la droite d'équation $y = 2x$). Ce n'est donc ni un isomorphisme, ni un automorphisme.
\end{corrige}

En résumé, nous avons appris à classer nos applications linéaires : un \cle{endomorphisme} reste dans le plan, un \cle{isomorphisme} est linéaire et bijectif, un \cle{automorphisme} cumule les deux qualités. Il reste une question pratique : comment décider \emph{rapidement} si un endomorphisme est bijectif, sans résoudre un système à chaque fois ? C'est toute l'utilité de l'écriture analytique dans une base, puis du noyau et de l'image, que nous étudierons dans les sections suivantes.

\coursec{Expression d'une application linéaire dans une base}

Dans la section précédente, nous avons appris à reconnaître les endomorphismes, les isomorphismes et les automorphismes du plan vectoriel, et nous avons découvert que la bijectivité joue un rôle central. Mais une question pratique reste en suspens : comment \emph{décrire concrètement} une application linéaire ? Dire « $f$ est linéaire » ne suffit pas pour calculer : il nous faut une \cle{formule}, une \cle{recette de calcul} qui donne l'image de n'importe quel vecteur. C'est exactement l'objet de cette section, et la bonne nouvelle est spectaculaire : pour connaître entièrement une application linéaire du plan, il suffit de connaître les images de \textbf{deux} vecteurs bien choisis.

\courssub{Le théorème fondamental : deux images suffisent}

Commençons par une intuition. Imagine un réseau de téléphonie mobile qui couvre le Cameroun : pour connaître la couverture partout, l'opérateur n'a pas besoin de tester chaque mètre carré du territoire ; il lui suffit de connaître la position et la portée de quelques antennes relais, le reste se déduit par calcul. Il en va de même pour les applications linéaires : le plan tout entier se fabrique à partir d'une base $(\vec{i}, \vec{j})$, donc l'application tout entière se fabrique à partir de $f(\vec{i})$ et $f(\vec{j})$.

Rappelons un fait essentiel sur les bases : si $(\vec{i}, \vec{j})$ est une base du plan vectoriel, alors \textbf{tout} vecteur $\vec{u}$ du plan s'écrit de manière \emph{unique} sous la forme
\[ \vec{u} = x\,\vec{i} + y\,\vec{j}, \]
où $x$ et $y$ sont les coordonnées de $\vec{u}$ dans cette base. Autrement dit, chaque vecteur est une « recette » : $x$ cuillères de $\vec{i}$, $y$ cuillères de $\vec{j}$. Or la linéarité de $f$ signifie précisément que $f$ \emph{respecte les recettes} : elle transforme une combinaison en la combinaison des images, avec les mêmes doses.

\begin{propriete}[Théorème fondamental : une application linéaire est déterminée par les images d'une base]
Soit $f$ une application linéaire du plan vectoriel dans lui-même et $(\vec{i}, \vec{j})$ une base du plan. Pour tout vecteur $\vec{u}$ de coordonnées $(x\,; y)$ dans cette base, c'est-à-dire $\vec{u} = x\,\vec{i} + y\,\vec{j}$, on a :
\[ f(\vec{u}) = x\,f(\vec{i}) + y\,f(\vec{j}). \]
Par conséquent, la donnée des deux vecteurs $f(\vec{i})$ et $f(\vec{j})$ détermine \textbf{entièrement} l'application $f$ : deux applications linéaires qui coïncident sur $\vec{i}$ et sur $\vec{j}$ sont égales partout.
\end{propriete}

\begin{experience}[Démonstration guidée du théorème fondamental]
Cette démonstration est courte, formatrice, et tu dois savoir la refaire : elle n'utilise que la définition de la linéarité. Suivons-la pas à pas.

\textbf{Hypothèses.} $f$ est linéaire, $(\vec{i}, \vec{j})$ est une base, et $\vec{u}$ est un vecteur quelconque du plan.

\textbf{Étape 1 : décomposer $\vec{u}$ dans la base.} Puisque $(\vec{i}, \vec{j})$ est une base, il existe deux réels uniques $x$ et $y$ tels que
\[ \vec{u} = x\,\vec{i} + y\,\vec{j}. \]

\textbf{Étape 2 : appliquer $f$ des deux côtés.}
\[ f(\vec{u}) = f\bigl(x\,\vec{i} + y\,\vec{j}\bigr). \]

\textbf{Étape 3 : utiliser l'additivité de $f$.} Par définition d'une application linéaire, l'image d'une somme est la somme des images :
\[ f\bigl(x\,\vec{i} + y\,\vec{j}\bigr) = f(x\,\vec{i}) + f(y\,\vec{j}). \]

\textbf{Étape 4 : utiliser l'homogénéité de $f$.} L'image d'un produit par un réel est le produit de l'image par ce réel :
\[ f(x\,\vec{i}) = x\,f(\vec{i}) \qquad \text{et} \qquad f(y\,\vec{j}) = y\,f(\vec{j}). \]

\textbf{Étape 5 : conclure.} En rassemblant les étapes 3 et 4 :
\[ f(\vec{u}) = x\,f(\vec{i}) + y\,f(\vec{j}). \qquad \blacksquare \]

\textbf{Lecture du résultat.} Remarque ce que dit cette formule : les coordonnées $x$ et $y$ de $\vec{u}$ « traversent » l'application sans changer. Seuls les vecteurs de base sont transformés. Connaître $f(\vec{i})$ et $f(\vec{j})$, c'est donc connaître $f$ tout entière.
\end{experience}

\begin{aretenir}[L'idée à retenir]
Une application linéaire du plan est un « moule » à deux ingrédients : dès que l'on sait ce qu'elle fait à $\vec{i}$ et à $\vec{j}$, on sait ce qu'elle fait à \textbf{tous} les vecteurs du plan, par la formule
\[ f(x\,\vec{i} + y\,\vec{j}) = x\,f(\vec{i}) + y\,f(\vec{j}). \]
C'est la propriété la plus importante de cette leçon : tout le reste en découle.
\end{aretenir}

\begin{popfigure}
\begin{center}
\begin{tikzpicture}[scale=1.05, >=Stealth]
% Grille légère
\draw[popline, very thin, step=1] (-2.5,-1.5) grid (5.5,4.5);
% Axes
\draw[->, popmuted] (-2.5,0) -- (5.7,0) node[below left] {\small $x$};
\draw[->, popmuted] (0,-1.5) -- (0,4.7) node[below left] {\small $y$};
% Base canonique
\draw[->, very thick, popblue] (0,0) -- (1,0) node[midway, below] {$\vec{i}$};
\draw[->, very thick, popblue] (0,0) -- (0,1) node[midway, left] {$\vec{j}$};
% Images de la base
\draw[->, very thick, poporange] (0,0) -- (2,1) node[midway, above left=-1pt] {$f(\vec{i})$};
\draw[->, very thick, popgreen] (0,0) -- (-1,2) node[midway, above right=-1pt] {$f(\vec{j})$};
% Vecteur u = 1.5 i + j -> (1.5, 1)
\draw[->, very thick, popink] (0,0) -- (1.5,1) node[above right=-2pt] {$\vec{u}$};
% Construction de f(u) = 1.5 f(i) + f(j)
\coordinate (FU) at (2,3.5); % 1.5*(2,1) + (-1,2) = (3-1, 1.5+2) = (2, 3.5)
\draw[->, very thick, poppink] (0,0) -- (FU) node[above] {$f(\vec{u})$};
% Parallélogramme de construction pour f(u)
\draw[->, thick, poporange!70, dashed] (0,0) -- (3,1.5) node[midway, below right=-2pt] {\small $1{,}5\,f(\vec{i})$};
\draw[->, thick, popgreen!80!black, dashed] (0,0) -- (-1,2) node[midway, above left=-2pt] {\small $f(\vec{j})$};
\draw[thick, popmuted, dashed] (3,1.5) -- (FU);
\draw[thick, popmuted, dashed] (-1,2) -- (FU);
% Points
\fill[popink] (FU) circle (1.6pt);
\fill[popink] (0,0) circle (1.6pt) node[below left=-2pt] {\small $O$};
\end{tikzpicture}
\end{center}
\legende{Construction de $f(\vec{u})$ par la règle du parallélogramme : ici $\vec{u} = 1{,}5\,\vec{i} + \vec{j}$, donc $f(\vec{u}) = 1{,}5\,f(\vec{i}) + f(\vec{j})$. Les coordonnées de $\vec{u}$ deviennent les coefficients de la combinaison des images.}
\end{popfigure}

\courssub{Écriture analytique dans une base}

Passons maintenant du dessin au calcul. Puisque $f(\vec{i})$ et $f(\vec{j})$ sont des vecteurs du plan, ils possèdent eux-mêmes des coordonnées dans la base $(\vec{i}, \vec{j})$. Notons :
\[ f(\vec{i}) = a\,\vec{i} + c\,\vec{j} \qquad \text{et} \qquad f(\vec{j}) = b\,\vec{i} + d\,\vec{j}, \]
c'est-à-dire : $f(\vec{i})$ a pour coordonnées $(a\,; c)$ et $f(\vec{j})$ a pour coordonnées $(b\,; d)$. Appliquons le théorème fondamental à un vecteur $\vec{u}$ de coordonnées $(x\,; y)$ :
\begin{align*}
f(\vec{u}) &= x\,f(\vec{i}) + y\,f(\vec{j}) \\
&= x\bigl(a\,\vec{i} + c\,\vec{j}\bigr) + y\bigl(b\,\vec{i} + d\,\vec{j}\bigr) \\
&= (ax + by)\,\vec{i} + (cx + dy)\,\vec{j}.
\end{align*}
Les coordonnées de $f(\vec{u})$ sont donc $(ax + by\,; cx + dy)$. Nous venons d'obtenir l'écriture analytique de $f$.

\begin{definition}[Écriture analytique d'une application linéaire dans une base]
Soit $f$ une application linéaire du plan et $(\vec{i}, \vec{j})$ une base. Si $f(\vec{i})$ a pour coordonnées $(a\,; c)$ et $f(\vec{j})$ a pour coordonnées $(b\,; d)$ dans cette base, alors pour tous réels $x$ et $y$ :
\[ f(x\,; y) = (ax + by\,; cx + dy). \]
Cette formule s'appelle l'\cle{écriture analytique} (ou \cle{expression analytique}) de $f$ dans la base $(\vec{i}, \vec{j})$.
\end{definition}

\begin{attention}[Qui donne les coefficients de $x$ ? Qui donne ceux de $y$ ?]
C'est \textbf{l'erreur la plus fréquente} au Baccalauréat, alors lis attentivement. Dans l'écriture
\[ f(x\,; y) = (ax + by\,; cx + dy), \]
les coefficients de $x$ (c'est-à-dire $a$ en abscisse et $c$ en ordonnée) proviennent de $f(\vec{i})$, et les coefficients de $y$ (c'est-à-dire $b$ et $d$) proviennent de $f(\vec{j})$. Beaucoup d'élèves inversent et écrivent les coordonnées de $f(\vec{i})$ en facteur de $y$. Pour ne jamais te tromper, retiens la logique de la démonstration :
\[ f(x\,; y) = x \times \underbrace{f(\vec{i})}_{(a\,; c)} \;+\; y \times \underbrace{f(\vec{j})}_{(b\,; d)}. \]
Le $x$ reste collé à $f(\vec{i})$, le $y$ reste collé à $f(\vec{j})$. Vérifie systématiquement ta formule en calculant $f(1\,; 0)$ : tu dois retrouver exactement les coordonnées de $f(\vec{i})$, et $f(0\,; 1)$ doit redonner celles de $f(\vec{j})$.
\end{attention}

\begin{methode}[Obtenir l'écriture analytique de $f$ dans la base $(\vec{i}, \vec{j})$]
\begin{enumerate}
\item \textbf{Calculer} les coordonnées de $f(\vec{i})$ dans la base $(\vec{i}, \vec{j})$ : notons-les $(a\,; c)$.
\item \textbf{Calculer} les coordonnées de $f(\vec{j})$ dans la base $(\vec{i}, \vec{j})$ : notons-les $(b\,; d)$.
\item \textbf{Écrire} la formule générale : $f(x\,; y) = x\,f(\vec{i}) + y\,f(\vec{j})$, c'est-à-dire
\[ f(x\,; y) = (ax + by\,; cx + dy). \]
\item \textbf{Contrôler} : vérifier que $f(1\,; 0) = (a\,; c)$ et $f(0\,; 1) = (b\,; d)$.
\end{enumerate}
\end{methode}

\begin{exemplebox}[Un exemple chiffré complet]
Dans une base $(\vec{i}, \vec{j})$ du plan, on considère l'application linéaire $f$ définie par
\[ f(\vec{i}) = 2\,\vec{i} + \vec{j} \qquad \text{et} \qquad f(\vec{j}) = -\vec{i} + 3\,\vec{j}. \]
\textbf{1. Écriture analytique de $f$.} Les coordonnées de $f(\vec{i})$ sont $(2\,; 1)$ et celles de $f(\vec{j})$ sont $(-1\,; 3)$. Pour un vecteur $\vec{u}$ de coordonnées $(x\,; y)$ :
\begin{align*}
f(\vec{u}) &= x\bigl(2\,\vec{i} + \vec{j}\bigr) + y\bigl(-\vec{i} + 3\,\vec{j}\bigr) \\
&= (2x - y)\,\vec{i} + (x + 3y)\,\vec{j}.
\end{align*}
Donc
\[ \boxed{\,f(x\,; y) = (2x - y\,; x + 3y)\,} \]
\textbf{Contrôle} : $f(1\,; 0) = (2\,; 1)$, qui est bien $f(\vec{i})$ ; $f(0\,; 1) = (-1\,; 3)$, qui est bien $f(\vec{j})$. La formule est cohérente.

\textbf{2. Image du vecteur de coordonnées $(4\,; -1)$.} On remplace $x$ par $4$ et $y$ par $-1$ :
\begin{align*}
f(4\,; -1) &= \bigl(2\times 4 - (-1)\,; 4 + 3\times(-1)\bigr) \\
&= (8 + 1\,; 4 - 3) \\
&= (9\,; 1).
\end{align*}
Autre présentation, directement issue du théorème : $f(4\,\vec{i} - \vec{j}) = 4\,f(\vec{i}) - f(\vec{j}) = 4(2\,\vec{i} + \vec{j}) - (-\vec{i} + 3\,\vec{j}) = 9\,\vec{i} + \vec{j}$. Même résultat, évidemment.
\end{exemplebox}

\courssub{Lecture inverse : retrouver $f(\vec{i})$ et $f(\vec{j})$ à partir de l'écriture analytique}

Au Baccalauréat, l'énoncé donne le plus souvent l'écriture analytique et te demande de travailler avec. Il faut donc savoir faire le chemin en sens inverse : lire $f(\vec{i})$ et $f(\vec{j})$ directement sur la formule.

\begin{propriete}[Lecture directe sur l'écriture analytique]
Si l'application linéaire $f$ est définie dans la base $(\vec{i}, \vec{j})$ par
\[ f(x\,; y) = (ax + by\,; cx + dy), \]
alors
\[ f(\vec{i}) = f(1\,; 0) = (a\,; c) = a\,\vec{i} + c\,\vec{j} \quad \text{et} \quad f(\vec{j}) = f(0\,; 1) = (b\,; d) = b\,\vec{i} + d\,\vec{j}. \]
Autrement dit : $f(\vec{i})$ se lit dans la « colonne des $x$ » et $f(\vec{j})$ dans la « colonne des $y$ ».
\end{propriete}

\begin{exemplebox}[Lecture inverse sur un exemple]
Soit $g$ l'application du plan définie dans une base $(\vec{i}, \vec{j})$ par
\[ g(x\,; y) = (3x - 5y\,; -2x + 4y). \]
On lit immédiatement :
\[ g(\vec{i}) = 3\,\vec{i} - 2\,\vec{j} \qquad \text{et} \qquad g(\vec{j}) = -5\,\vec{i} + 4\,\vec{j}. \]
Remarque au passage que cette écriture prouve aussi que $g$ est linéaire : ses deux coordonnées sont des expressions de la forme $ax + by$, sans constante ni terme en $x^2$, $xy$, etc.
\end{exemplebox}

\courssub{Présentation en tableau : vers la notion de matrice}

Les quatre coefficients $a, b, c, d$ contiennent toute l'information sur $f$. Il est commode de les ranger dans un tableau à deux lignes et deux colonnes, en plaçant \textbf{en colonnes} les coordonnées de $f(\vec{i})$ et de $f(\vec{j})$ :
\[ f(x\,; y) = (ax + by\,; cx + dy) \qquad \longleftrightarrow \qquad \begin{pmatrix} a & b \\ c & d \end{pmatrix} \]
La première colonne contient les coordonnées de $f(\vec{i})$, la deuxième celles de $f(\vec{j})$ ; la première ligne correspond à l'abscisse de l'image, la deuxième à son ordonnée.

\begin{savaistu}[Complément Bac : ce tableau s'appelle une matrice]
Ce tableau de nombres porte un nom : c'est la \cle{matrice} de l'application linéaire $f$ dans la base $(\vec{i}, \vec{j})$. Les matrices ne figurent pas au programme strict de Terminale TI, mais tu les rencontreras dès la première année de l'enseignement supérieur (classes préparatoires, universités, écoles d'ingénieurs comme l'ENSP de Yaoundé ou l'IUT de Douala). Elles permettent de calculer des images, de composer des applications linéaires et de décider de la bijectivité par un simple calcul de déterminant : $f$ est bijective si et seulement si $ad - bc \neq 0$. Retiens simplement pour l'instant : \emph{une application linéaire du plan, c'est quatre nombres rangés dans le bon ordre}.
\end{savaistu}

\courssub{Et si la base n'est pas orthonormée ?}

Un dernier point de vigilance. Tout ce que nous avons fait dans cette section repose sur un seul ingrédient : le fait que $(\vec{i}, \vec{j})$ soit une \cle{base}, c'est-à-dire deux vecteurs non colinéaires, qui permettent d'écrire chaque vecteur du plan de façon unique sous la forme $x\,\vec{i} + y\,\vec{j}$. Nulle part nous n'avons utilisé le fait que $\vec{i}$ et $\vec{j}$ soient perpendiculaires ou de norme $1$.

\begin{aretenir}[La méthode est universelle]
Le théorème fondamental, l'écriture analytique $f(x\,; y) = (ax + by\,; cx + dy)$ et la lecture inverse restent \textbf{exactement identiques} dans une base quelconque, orthonormée ou non. Seule l'interprétation géométrique des coordonnées change (axes obliques, unités différentes), jamais la méthode algébrique.
\end{aretenir}

\begin{exoresolu}[Application linéaire dans une base non orthonormée]
Le plan est rapporté à une base $(\vec{i}, \vec{j})$ non orthonormée. Soit $f$ l'application linéaire définie par
\[ f(\vec{i}) = \vec{i} + 2\,\vec{j} \qquad \text{et} \qquad f(\vec{j}) = 3\,\vec{i} - \vec{j}. \]
\begin{enumerate}
\item Déterminer l'écriture analytique de $f$ dans la base $(\vec{i}, \vec{j})$.
\item Déterminer l'image du vecteur $\vec{u} = 2\,\vec{i} - 5\,\vec{j}$.
\item Le vecteur $\vec{v} = -\vec{i} + 7\,\vec{j}$ est-il l'image d'un vecteur de coordonnées $(1\,; y)$ pour un certain réel $y$ ?
\end{enumerate}
\tcblower
\textbf{1. Écriture analytique.} Les coordonnées de $f(\vec{i})$ sont $(1\,; 2)$ et celles de $f(\vec{j})$ sont $(3\,; -1)$. Le théorème fondamental s'applique sans modification, même si la base n'est pas orthonormée :
\[ f(x\,; y) = x\,f(\vec{i}) + y\,f(\vec{j}) = (x + 3y\,; 2x - y). \]
Contrôle : $f(1\,; 0) = (1\,; 2) = f(\vec{i})$ et $f(0\,; 1) = (3\,; -1) = f(\vec{j})$. C'est cohérent.

\textbf{2. Image de $\vec{u}$.} Le vecteur $\vec{u}$ a pour coordonnées $(2\,; -5)$. On applique la formule :
\[ f(2\,; -5) = \bigl(2 + 3\times(-5)\,; 2\times 2 - (-5)\bigr) = (2 - 15\,; 4 + 5) = (-13\,; 9). \]
Donc $f(\vec{u}) = -13\,\vec{i} + 9\,\vec{j}$.

\textbf{3. Recherche de $y$.} On cherche un réel $y$ tel que $f(1\,; y) = (-1\,; 7)$. Or
\[ f(1\,; y) = (1 + 3y\,; 2 - y). \]
On obtient le système
\[ \begin{cases} 1 + 3y = -1 \\ 2 - y = 7 \end{cases} \qquad \text{c'est-à-dire} \qquad \begin{cases} y = -\dfrac{2}{3} \\[4pt] y = -5 \end{cases} \]
Les deux équations donnent des valeurs différentes de $y$ : le système est incompatible. Il n'existe donc \textbf{aucun} réel $y$ tel que $f(1\,; y) = \vec{v}$. Ce petit calcul annonce la section suivante : nous verrons que l'ensemble des vecteurs atteints par $f$ (l'\emph{image} de $f$) et l'ensemble des vecteurs envoyés sur $\vec{0}$ (le \emph{noyau} de $f$) sont les clés pour décider si $f$ est bijective.
\end{exoresolu}

\begin{exercice}[À toi de jouer]
Dans une base $(\vec{i}, \vec{j})$ du plan, on considère l'application linéaire $h$ définie par $h(\vec{i}) = -\vec{i} + 4\,\vec{j}$ et $h(\vec{j}) = 2\,\vec{i} + \vec{j}$.
\begin{enumerate}
\item Donner l'écriture analytique de $h$ dans la base $(\vec{i}, \vec{j})$.
\item Calculer les images des vecteurs de coordonnées $(3\,; 2)$ et $(-1\,; 5)$.
\item Écrire le tableau (matrice) des coefficients de $h$ dans la base $(\vec{i}, \vec{j})$.
\end{enumerate}
\end{exercice}

\begin{corrige}[Exercice : À toi de jouer]
\textbf{1.} Les coordonnées de $h(\vec{i})$ sont $(-1\,; 4)$, celles de $h(\vec{j})$ sont $(2\,; 1)$. Donc
\[ h(x\,; y) = x\,h(\vec{i}) + y\,h(\vec{j}) = (-x + 2y\,; 4x + y). \]
Contrôle : $h(1\,; 0) = (-1\,; 4)$ et $h(0\,; 1) = (2\,; 1)$ : conforme.

\textbf{2.} $h(3\,; 2) = (-3 + 4\,; 12 + 2) = (1\,; 14)$ et $h(-1\,; 5) = (1 + 10\,; -4 + 5) = (11\,; 1)$.

\textbf{3.} En plaçant $h(\vec{i})$ et $h(\vec{j})$ en colonnes :
\[ \begin{pmatrix} -1 & 2 \\ 4 & 1 \end{pmatrix}. \]
\end{corrige}

\coursec{Noyau d'une application linéaire}

Dans la section précédente, nous avons appris à \emph{représenter} une application linéaire $f$ du plan par son écriture analytique dans une base $(\vec{i}, \vec{j})$ : à tout vecteur $\vec{u}(x ; y)$, on associe $f(\vec{u})(x' ; y')$ avec
\[
\begin{cases} x' = ax + by \\ y' = cx + dy \end{cases}
\]
Cette écriture nous donne un outil de calcul puissant. Mais elle ne répond pas encore à une question fondamentale : \textbf{que fait exactement $f$ du plan ?} Certains vecteurs sont-ils « écrasés » sur le vecteur nul ? L'application écrase-t-elle tout le plan, ou seulement une partie ? La réponse à la première de ces questions s'appelle le \cle{noyau} de $f$. C'est l'objet de cette section.

\courssub{Intuition : les vecteurs que l'application « écrase »}

Imagine une projection orthogonale sur une droite $(D)$ du plan, par exemple la projection sur l'axe des abscisses : $p(x ; y) = (x ; 0)$. Tous les vecteurs \emph{verticaux}, c'est-à-dire de la forme $(0 ; y)$, sont envoyés sur le vecteur nul $(0 ; 0)$. On dit qu'ils sont \emph{écrasés} par $p$. L'ensemble de tous ces vecteurs écrasés forme l'axe des ordonnées : c'est le noyau de $p$.

Cette idée est très générale : pour n'importe quelle application linéaire $f$ du plan, on peut regarder l'ensemble des vecteurs que $f$ envoie sur $\vec{0}$. Cet ensemble mesure « ce que $f$ perd » comme information.

\begin{definition}[Noyau d'une application linéaire]
Soit $f$ une application linéaire du plan vectoriel $\mathcal{V}$ dans lui-même. On appelle \cle{noyau} de $f$, et on note $\ker(f)$ (ou $\mathrm{Ker}(f)$), l'ensemble des vecteurs $\vec{u}$ du plan dont l'image par $f$ est le vecteur nul :
\[
\ker(f) = \left\{ \vec{u} \in \mathcal{V} \; \middle| \; f(\vec{u}) = \vec{0} \right\}.
\]
Autrement dit, $\ker(f)$ est l'ensemble des \emph{antécédents} du vecteur nul par $f$.
\end{definition}

\begin{attention}[Le noyau n'est jamais vide]
Erreur fréquente : croire qu'une application linéaire peut avoir un noyau vide. C'est impossible ! Comme $f$ est linéaire, on a toujours $f(\vec{0}) = \vec{0}$. Donc $\vec{0} \in \ker(f)$ dans \textbf{tous} les cas. Le noyau contient \emph{au minimum} le vecteur nul. La question intéressante est donc : contient-il \emph{autre chose} que $\vec{0}$ ?
\end{attention}

\courssub{Détermination pratique du noyau : le système homogène}

Grâce à l'écriture analytique de $f$ dans une base $(\vec{i}, \vec{j})$, la recherche du noyau devient un simple calcul. Dire que $\vec{u}(x ; y)$ appartient à $\ker(f)$, c'est dire que $f(\vec{u}) = \vec{0}$, c'est-à-dire que ses coordonnées-images sont nulles :
\[
\begin{cases} ax + by = 0 \\ cx + dy = 0 \end{cases}
\]
On obtient un \emph{système linéaire homogène} (les seconds membres sont nuls). Résoudre ce système, c'est déterminer entièrement le noyau.

\begin{methode}[Déterminer le noyau d'une application linéaire du plan]
Soit $f$ l'application linéaire définie dans une base $(\vec{i}, \vec{j})$ par $f(x ; y) = (ax + by \,;\, cx + dy)$.
\begin{enumerate}
\item \textbf{Écrire le système homogène} associé : $\begin{cases} ax + by = 0 \\ cx + dy = 0 \end{cases}$.
\item \textbf{Résoudre le système} en exprimant une inconnue en fonction de l'autre (ou en montrant que la seule solution est $(0 ; 0)$).
\item \textbf{Conclure sur la nature du noyau} :
\begin{itemize}
\item si la seule solution est $(0 ; 0)$, alors $\ker(f) = \{\vec{0}\}$ ;
\item si les solutions sont de la forme $(x ; y) = t \cdot (x_0 ; y_0)$ avec $t \in \mathbb{R}$, alors $\ker(f)$ est la \emph{droite vectorielle} de vecteur directeur $(x_0 ; y_0)$ ;
\item si tout couple $(x ; y)$ est solution, alors $\ker(f)$ est le plan tout entier (cas où $f$ est l'application nulle).
\end{itemize}
\item \textbf{Donner une base du noyau} : quand $\ker(f)$ est une droite, une base est formée d'un seul vecteur directeur non nul ; quand $\ker(f) = \{\vec{0}\}$, le noyau est \emph{nul} (sa base est la famille vide, sa dimension est $0$).
\end{enumerate}
\end{methode}

\begin{exemplebox}[Un noyau qui est une droite]
Soit $f$ l'application linéaire définie dans la base $(\vec{i}, \vec{j})$ par
\[
f(x ; y) = (2x - y \,;\, -4x + 2y).
\]
Un vecteur $(x ; y)$ est dans $\ker(f)$ si et seulement si
\[
\begin{cases} 2x - y = 0 \\ -4x + 2y = 0 \end{cases}
\]
La deuxième équation est exactement la première multipliée par $-2$ : le système se réduit à l'unique équation $2x - y = 0$, c'est-à-dire $y = 2x$. Les solutions sont donc les couples $(x ; 2x) = x(1 ; 2)$, avec $x \in \mathbb{R}$. Conclusion :
\[
\ker(f) = \left\{ x(1 ; 2) \; ; \; x \in \mathbb{R} \right\} = \mathrm{Vect}\big( (1 ; 2) \big).
\]
Le noyau est la \textbf{droite vectorielle d'équation $2x - y = 0$}, et une \textbf{base} de $\ker(f)$ est le vecteur $(1 ; 2)$.
\end{exemplebox}

\begin{definition}[Équation caractéristique du noyau]
On appelle \cle{équation caractéristique} du noyau l'équation (ou le système d'équations) cartésienne dont l'ensemble des solutions est exactement $\ker(f)$. Dans l'exemple précédent, l'équation caractéristique du noyau est $2x - y = 0$. Lorsque $\ker(f) = \{\vec{0}\}$, le noyau est caractérisé par le système $x = 0$ et $y = 0$.
\end{definition}

\courssub{Nature géométrique du noyau dans le plan}

Dans le plan, seules trois situations sont possibles pour le noyau d'une application linéaire :

\begin{center}
\begin{tabularx}{\linewidth}{|l|X|X|}
\hline
\rowcolor{popblueL} \textbf{Noyau} & \textbf{Signification} & \textbf{Exemple} \\
\hline
$\ker(f) = \{\vec{0}\}$ & Seul le vecteur nul est écrasé : $f$ ne perd aucune information. & $f(x;y) = (x + y \,;\, x - y)$ \\
\hline
$\ker(f)$ est une droite vectorielle & Tous les vecteurs d'une direction sont écrasés sur $\vec{0}$. & $f(x;y) = (2x - y \,;\, -4x + 2y)$ \\
\hline
$\ker(f) = \mathcal{V}$ (le plan entier) & Tout est écrasé : $f$ est l'application nulle. & $f(x;y) = (0 ; 0)$ \\
\hline
\end{tabularx}
\end{center}

\begin{propriete}[Stabilité du noyau — Sous-espace vectoriel]
Le noyau d'une application linéaire $f$ est un \textbf{sous-espace vectoriel} du plan de départ. Cela signifie qu'il est stable par addition et par multiplication par un réel :
\begin{itemize}
\item Si $\vec{u} \in \ker(f)$ et $\vec{v} \in \ker(f)$, alors $\vec{u} + \vec{v} \in \ker(f)$.
\item Si $\vec{u} \in \ker(f)$ et $\lambda \in \mathbb{R}$, alors $\lambda \vec{u} \in \ker(f)$.
\end{itemize}
\emph{Justification :} $f(\vec{u}+\vec{v}) = f(\vec{u})+f(\vec{v}) = \vec{0}+\vec{0} = \vec{0}$ et $f(\lambda \vec{u}) = \lambda f(\vec{u}) = \lambda \vec{0} = \vec{0}$. Cette propriété explique pourquoi le noyau ne peut être que $\{\vec{0}\}$, une droite vectorielle, ou le plan tout entier (les trois seuls sous-espaces vectoriels du plan).
\end{propriete}

La figure suivante illustre le cas intermédiaire, le plus riche : le noyau est une droite, et tous ses vecteurs sont envoyés sur l'origine.

\begin{popfigure}
\begin{tikzpicture}[scale=1.1, >=Stealth]
% Axes
\draw[->, popmuted, thick] (-3.2,0) -- (3.2,0) node[below right] {$x$};
\draw[->, popmuted, thick] (0,-2.6) -- (0,2.6) node[above left] {$y$};
% Droite Ker(f) : y = 2x
\draw[popblue, very thick] (-1.35,-2.7) -- (1.35,2.7) node[above right, popblue] {$\ker(f) : 2x - y = 0$};
% Vecteurs du noyau
\draw[->, poporange, very thick] (0,0) -- (1,2) node[midway, right, poporange] {$(1 ; 2)$};
\draw[->, poporange, thick] (0,0) -- (0.6,1.2);
\draw[->, poporange, thick] (0,0) -- (-0.8,-1.6);
% Flèches vers l'origine (vecteurs écrasés)
\draw[->, poppink, thick, dashed] (1.15,2.3) .. controls (0.7,1.5) .. (0.12,0.12);
\draw[->, poppink, thick, dashed] (-1.0,-2.0) .. controls (-0.6,-1.2) .. (-0.12,-0.12);
\draw[->, poppink, thick, dashed] (0.75,1.5) .. controls (0.5,0.9) .. (0.1,0.1);
% Origine
\fill[popdark] (0,0) circle (2.2pt) node[below left, popdark] {$\vec{0}$};
\node[poppink, align=center] at (2.6,-1.9) {\small tous envoyés\\ \small sur $\vec{0}$ par $f$};
\end{tikzpicture}
\legende{Le noyau de $f(x;y) = (2x - y \,;\, -4x + 2y)$ est la droite d'équation $2x - y = 0$. Tous les vecteurs de cette droite (en orange) sont « écrasés » sur le vecteur nul.}
\end{popfigure}

\courssub{Noyau et injectivité : le théorème fondamental}

Rappelons qu'une application est \emph{injective} lorsque deux vecteurs distincts ont toujours des images distinctes. Le noyau donne un critère remarquablement simple pour tester l'injectivité d'une application linéaire : il suffit de regarder si $\vec{0}$ est le seul vecteur écrasé.

\begin{propriete}[Caractérisation de l'injectivité par le noyau]
Soit $f$ une application linéaire du plan dans lui-même. Alors :
\[
f \text{ est injective} \quad \Longleftrightarrow \quad \ker(f) = \{\vec{0}\}.
\]
\end{propriete}

\begin{experience}[Démonstration guidée du théorème]
Nous démontrons les deux implications. La clé est la linéarité : $f(\vec{u}) - f(\vec{v}) = f(\vec{u} - \vec{v})$.

\textbf{Sens direct} : supposons $f$ injective et montrons que $\ker(f) = \{\vec{0}\}$.
\begin{enumerate}
\item Soit $\vec{u} \in \ker(f)$. Que vaut $f(\vec{u})$ ? Par définition du noyau, $f(\vec{u}) = \vec{0}$.
\item Or $f(\vec{0}) = \vec{0}$ aussi (linéarité). Donc $f(\vec{u}) = f(\vec{0})$.
\item Comme $f$ est injective, deux vecteurs ayant la même image sont égaux : $\vec{u} = \vec{0}$.
\item Conclusion : le seul élément du noyau est $\vec{0}$, donc $\ker(f) = \{\vec{0}\}$.
\end{enumerate}

\textbf{Sens réciproque} : supposons $\ker(f) = \{\vec{0}\}$ et montrons que $f$ est injective.
\begin{enumerate}
\item Soient $\vec{u}$ et $\vec{v}$ deux vecteurs tels que $f(\vec{u}) = f(\vec{v})$. Nous devons prouver que $\vec{u} = \vec{v}$.
\item Par linéarité : $f(\vec{u} - \vec{v}) = f(\vec{u}) - f(\vec{v}) = \vec{0}$.
\item Donc $\vec{u} - \vec{v} \in \ker(f)$. Mais par hypothèse $\ker(f) = \{\vec{0}\}$, donc $\vec{u} - \vec{v} = \vec{0}$.
\item Conclusion : $\vec{u} = \vec{v}$, donc $f$ est injective. $\blacksquare$
\end{enumerate}
\end{experience}

\begin{aretenir}[L'idée à retenir]
Pour une application linéaire, l'injectivité se lit \emph{uniquement} sur le noyau : si $f$ n'écrase que le vecteur nul, elle est injective ; si elle écrase toute une droite de vecteurs, alors chaque image a une infinité d'antécédents (tous obtenus en ajoutant un vecteur du noyau), et $f$ n'est pas injective. C'est un critère purement algébrique : une simple résolution de système homogène suffit.
\end{aretenir}

\begin{exoresolu}[Noyau, base et injectivité]
\textbf{Énoncé.} On considère les applications linéaires définies dans une base $(\vec{i}, \vec{j})$ par :
\[
f(x ; y) = (2x - y \,;\, -4x + 2y) \qquad \text{et} \qquad g(x ; y) = (x + 3y \,;\, 2x - y).
\]
\begin{enumerate}
\item Déterminer $\ker(f)$, en donner une équation caractéristique et une base. L'application $f$ est-elle injective ?
\item Déterminer $\ker(g)$. L'application $g$ est-elle injective ?
\end{enumerate}
\tcblower
\textbf{Solution détaillée.}

\textbf{1.} Un vecteur $(x ; y)$ appartient à $\ker(f)$ si et seulement si
\[
\begin{cases} 2x - y = 0 \\ -4x + 2y = 0 \end{cases}
\]
La seconde équation vaut $-2$ fois la première : le système équivaut à $y = 2x$. Les solutions sont les $(x ; 2x) = x(1 ; 2)$, $x \in \mathbb{R}$. Donc
\[
\ker(f) = \mathrm{Vect}\big((1 ; 2)\big), \quad \text{équation caractéristique : } 2x - y = 0,
\]
et une base de $\ker(f)$ est le vecteur $(1 ; 2)$. Comme $\ker(f) \neq \{\vec{0}\}$, l'application $f$ \textbf{n'est pas injective}. Vérification : $f(1 ; 2) = (2 - 2 \,;\, -4 + 4) = (0 ; 0) = f(0 ; 0)$, avec $(1 ; 2) \neq (0 ; 0)$.

\textbf{2.} Un vecteur $(x ; y)$ appartient à $\ker(g)$ si et seulement si
\[
\begin{cases} x + 3y = 0 \\ 2x - y = 0 \end{cases}
\]
De la deuxième équation, $y = 2x$. En reportant dans la première : $x + 3(2x) = 7x = 0$, donc $x = 0$, puis $y = 0$. La seule solution est $(0 ; 0)$, donc $\ker(g) = \{\vec{0}\}$. Par le théorème, $g$ \textbf{est injective}.
\end{exoresolu}

\begin{attention}[Pièges de rédaction autour du noyau]
\begin{itemize}
\item Ne jamais écrire $\ker(f) = \emptyset$ : le noyau contient toujours $\vec{0}$. Le « plus petit » noyau possible est $\{\vec{0}\}$.
\item Ne pas confondre $\ker(f)$ (ensemble de \emph{vecteurs du plan de départ}) avec une condition sur les images : on résout $f(\vec{u}) = \vec{0}$, pas $f(\vec{u}) = \vec{u}$.
\item Quand le système homogène se réduit à une seule équation, ne pas conclure « le noyau est $\mathbb{R}$ » : c'est une \emph{droite de vecteurs}, qu'il faut écrire sous la forme $\mathrm{Vect}((x_0 ; y_0))$ en exhibant un vecteur directeur.
\item Pour justifier la non-injectivité, il ne suffit pas de dire « le système a plusieurs solutions » : il faut exhiber (ou citer) un vecteur non nul du noyau, par exemple $(1 ; 2)$ dans l'exemple de $f$.
\end{itemize}
\end{attention}

\begin{savaistu}[Le noyau dans ta poche : la compression d'images]
Quand tu envoies une photo par WhatsApp ou que tu la publies sur les réseaux sociaux, elle est \emph{compressée} (format JPEG) pour passer plus vite, même avec un réseau MTN ou Orange capricieux. Le principe mathématique ? L'image est découpée en blocs, chaque bloc est transformé par une application linéaire (la « transformée en cosinus »), puis on \emph{supprime volontairement} les composantes qui portent peu d'information visible : on projette sur un sous-espace, et tout ce qui vit dans le \textbf{noyau} de cette projection est définitivement perdu. C'est exactement l'idée de cette leçon : ce que le noyau contient ne peut plus être récupéré. Les mathématiciens du signal dosent la taille du noyau pour trouver le meilleur compromis entre légèreté du fichier et qualité de l'image !
\end{savaistu}

\begin{exercice}[À toi de jouer]
Dans une base $(\vec{i}, \vec{j})$, on considère l'application linéaire $h(x ; y) = (3x + y \,;\, 6x + 2y)$.
\begin{enumerate}
\item Déterminer $\ker(h)$ et en donner une base.
\item Donner une équation caractéristique de $\ker(h)$.
\item L'application $h$ est-elle injective ? Justifier.
\end{enumerate}
\end{exercice}

\begin{corrige}[Exercice : noyau de $h$]
\textbf{1.} $(x ; y) \in \ker(h)$ équivaut à $\begin{cases} 3x + y = 0 \\ 6x + 2y = 0 \end{cases}$. La deuxième équation est le double de la première : le système équivaut à $y = -3x$. Les solutions sont $(x ; -3x) = x(1 ; -3)$, $x \in \mathbb{R}$. Donc $\ker(h) = \mathrm{Vect}\big((1 ; -3)\big)$, de base le vecteur $(1 ; -3)$.

\textbf{2.} Une équation caractéristique du noyau est $3x + y = 0$.

\textbf{3.} Comme $\ker(h) \neq \{\vec{0}\}$ (il contient par exemple $(1 ; -3) \neq \vec{0}$), l'application $h$ n'est pas injective, d'après le théorème de caractérisation de l'injectivité par le noyau.
\end{corrige}

Le noyau nous renseigne donc sur ce que $f$ « perd ». Symétriquement, il est naturel de s'intéresser à ce que $f$ « atteint » : l'ensemble de toutes les images possibles. C'est l'objet de la section suivante, consacrée à l'\emph{image} d'une application linéaire — et nous verrons que noyau et image sont deux faces d'une même médaille.

\coursec{Image d'une application linéaire et stabilité}

Dans la section précédente, nous avons étudié le \cle{noyau} d'une application linéaire $f$, c'est-à-dire l'ensemble des vecteurs que $f$ \og écrase \fg{} sur le vecteur nul : $\mathrm{Ker}(f)$ répond à la question \emph{d'où partent les vecteurs qui disparaissent ?}. Nous allons maintenant regarder l'autre extrémité du voyage : \emph{où arrivent les vecteurs ?} Autrement dit, quelle portion du plan d'arrivée est réellement atteinte par $f$ ? C'est exactement la notion d'\cle{image} d'une application linéaire. Noyau et image sont les deux jumeaux qui, ensemble, décrivent entièrement le comportement d'une application linéaire.

\courssub{Définition de l'image d'une application linéaire}

\begin{experience}[Une intuition avant la définition]
Imagine un photographe qui prend en photo la place du marché de Mokolo. Le plan de départ, c'est la réalité ; le plan d'arrivée, c'est la pellicule. Certains points de la réalité peuvent se superposer sur la photo (plusieurs objets alignés dans l'axe de l'appareil donnent le même point-image) : c'est le rôle du noyau. Mais surtout, la photo ne montre pas \emph{tout} l'univers : elle ne montre que ce que l'objectif a capté. L'ensemble des points réellement visibles sur la photo, c'est l'\emph{image} du photographe. Pour une application linéaire $f$ du plan dans lui-même, l'image est l'ensemble des vecteurs que $f$ peut effectivement produire.
\end{experience}

\begin{definition}[Image d'une application linéaire]
Soit $f$ une application linéaire du plan vectoriel $\mathcal{P}$ dans lui-même. On appelle \cle{image} de $f$, et on note $\mathrm{Im}(f)$, l'ensemble des vecteurs de la forme $f(\vec{u})$ lorsque $\vec{u}$ décrit tout le plan :
\[ \mathrm{Im}(f) = \left\{\, f(\vec{u}) \; ; \; \vec{u} \in \mathcal{P} \,\right\}. \]
Autrement dit, un vecteur $\vec{v}$ appartient à $\mathrm{Im}(f)$ si et seulement s'il existe (au moins) un vecteur $\vec{u}$ tel que $f(\vec{u}) = \vec{v}$ : on dit que $\vec{v}$ \emph{a un antécédent} par $f$.
\end{definition}

\begin{attention}[Ne pas confondre image et ensemble d'arrivée]
L'ensemble d'arrivée de $f$ est toujours le plan $\mathcal{P}$ tout entier. Mais $\mathrm{Im}(f)$ peut être strictement plus petit : ce sont les vecteurs \emph{effectivement atteints}. Dire $\vec{v} \in \mathcal{P}$ ne signifie pas $\vec{v} \in \mathrm{Im}(f)$. Par exemple, si $f$ est l'application nulle, $\mathrm{Im}(f) = \{\vec{0}\}$ alors que l'ensemble d'arrivée est le plan entier.
\end{attention}

\courssub{Propriété fondamentale : l'image est engendrée par $f(\vec{\imath})$ et $f(\vec{\jmath})$}

Voici la propriété qui rend le calcul de $\mathrm{Im}(f)$ presque mécanique. Son intuition est simple : une application linéaire est entièrement déterminée par ce qu'elle fait aux deux vecteurs de base ; donc tout ce qu'elle produit doit pouvoir se fabriquer à partir de $f(\vec{\imath})$ et $f(\vec{\jmath})$.

\begin{propriete}[$\mathrm{Im}(f)$ est engendrée par les images des vecteurs de base]
Soit $f$ une application linéaire du plan et $(\vec{\imath}, \vec{\jmath})$ une base du plan. Alors
\[ \mathrm{Im}(f) = \mathrm{Vect}\bigl(f(\vec{\imath}),\, f(\vec{\jmath})\bigr), \]
c'est-à-dire que $\mathrm{Im}(f)$ est l'ensemble de toutes les combinaisons linéaires $\alpha\, f(\vec{\imath}) + \beta\, f(\vec{\jmath})$ avec $\alpha, \beta \in \mathbb{R}$.
\end{propriete}

\begin{experience}[Démonstration (exigible)]
Montrons la double inclusion.
\begin{enumerate}
\item \textbf{Tout élément de $\mathrm{Im}(f)$ est combinaison linéaire de $f(\vec{\imath})$ et $f(\vec{\jmath})$.} Soit $\vec{v} \in \mathrm{Im}(f)$. Il existe $\vec{u} \in \mathcal{P}$ tel que $\vec{v} = f(\vec{u})$. Comme $(\vec{\imath}, \vec{\jmath})$ est une base, il existe deux réels $x$ et $y$ tels que $\vec{u} = x\,\vec{\imath} + y\,\vec{\jmath}$. Par linéarité de $f$ :
\[ \vec{v} = f(\vec{u}) = f\bigl(x\,\vec{\imath} + y\,\vec{\jmath}\bigr) = x\,f(\vec{\imath}) + y\,f(\vec{\jmath}). \]
Donc $\vec{v}$ est bien combinaison linéaire de $f(\vec{\imath})$ et $f(\vec{\jmath})$.
\item \textbf{Toute combinaison linéaire de $f(\vec{\imath})$ et $f(\vec{\jmath})$ est dans $\mathrm{Im}(f)$.} Soient $\alpha, \beta \in \mathbb{R}$. Posons $\vec{u} = \alpha\,\vec{\imath} + \beta\,\vec{\jmath}$. Alors, toujours par linéarité :
\[ f(\vec{u}) = \alpha\,f(\vec{\imath}) + \beta\,f(\vec{\jmath}), \]
donc $\alpha\,f(\vec{\imath}) + \beta\,f(\vec{\jmath})$ est l'image de $\vec{u}$ : il appartient à $\mathrm{Im}(f)$.
\end{enumerate}
Les deux inclusions étant établies, $\mathrm{Im}(f) = \mathrm{Vect}\bigl(f(\vec{\imath}), f(\vec{\jmath})\bigr)$. \hfill$\blacksquare$
\end{experience}

\begin{attention}[L'erreur classique du Bac]
Il est fréquent de lire : \og $\mathrm{Im}(f)$ est engendrée par $\vec{\imath}$ et $\vec{\jmath}$ \fg{}. C'est \textbf{faux} ! Ce sont les \emph{images} $f(\vec{\imath})$ et $f(\vec{\jmath})$ qui engendrent $\mathrm{Im}(f)$, pas les vecteurs de base eux-mêmes. Les vecteurs $\vec{\imath}$ et $\vec{\jmath}$ vivent dans le plan de \emph{départ} ; l'image vit dans le plan d'\emph{arrivée}. Écrire $\mathrm{Im}(f) = \mathrm{Vect}(\vec{\imath}, \vec{\jmath})$ reviendrait à dire que toute application linéaire est surjective, ce qui est évidemment faux (pense à l'application nulle).
\end{attention}

\courssub{Nature géométrique de l'image dans le plan}

Puisque $\mathrm{Im}(f) = \mathrm{Vect}\bigl(f(\vec{\imath}), f(\vec{\jmath})\bigr)$, tout dépend de la famille $\bigl(f(\vec{\imath}), f(\vec{\jmath})\bigr)$. Dans le plan, trois cas seulement sont possibles.

\begin{propriete}[Les trois visages de l'image]
Soit $f$ une application linéaire du plan et $(\vec{\imath}, \vec{\jmath})$ une base. Alors $\mathrm{Im}(f)$ est :
\begin{enumerate}
\item \textbf{réduite au vecteur nul} $\{\vec{0}\}$, si $f(\vec{\imath}) = f(\vec{\jmath}) = \vec{0}$ (c'est l'application nulle) ;
\item \textbf{une droite vectorielle}, si $f(\vec{\imath})$ et $f(\vec{\jmath})$ sont colinéaires et non tous deux nuls ;
\item \textbf{le plan tout entier} $\mathcal{P}$, si $f(\vec{\imath})$ et $f(\vec{\jmath})$ ne sont pas colinéaires (ils forment alors une base du plan).
\end{enumerate}
\end{propriete}

\begin{aretenir}[Base de l'image]
Pour obtenir une \cle{base} de $\mathrm{Im}(f)$, on \emph{extrait une famille libre} de la famille génératrice $\bigl(f(\vec{\imath}), f(\vec{\jmath})\bigr)$ :
\begin{itemize}
\item si les deux vecteurs sont non colinéaires, $\bigl(f(\vec{\imath}), f(\vec{\jmath})\bigr)$ est déjà une base de $\mathrm{Im}(f) = \mathcal{P}$ ;
\item s'ils sont colinéaires, on garde l'un des deux (un vecteur non nul) : il forme une base de la droite $\mathrm{Im}(f)$ ;
\item s'ils sont tous deux nuls, $\mathrm{Im}(f) = \{\vec{0}\}$ n'a pas de base.
\end{itemize}
Le nombre de vecteurs d'une base de $\mathrm{Im}(f)$ s'appelle le \cle{rang} de $f$ : il vaut $0$, $1$ ou $2$.
\end{aretenir}

\begin{methode}[Déterminer l'image d'une application linéaire]
Pour déterminer $\mathrm{Im}(f)$ lorsque $f$ est donnée par son écriture analytique :
\begin{enumerate}
\item Calculer $f(\vec{\imath})$ et $f(\vec{\jmath})$ en prenant $(x\,;y) = (1\,;0)$ puis $(x\,;y) = (0\,;1)$. Ce sont les \emph{colonnes} de l'écriture de $f$.
\item Tester la colinéarité de $f(\vec{\imath})$ et $f(\vec{\jmath})$ à l'aide du déterminant : $\det\bigl(f(\vec{\imath}), f(\vec{\jmath})\bigr)$.
\item Conclure :
\begin{itemize}
\item déterminant non nul $\Rightarrow$ $\mathrm{Im}(f) = \mathcal{P}$, base $\bigl(f(\vec{\imath}), f(\vec{\jmath})\bigr)$ ;
\item déterminant nul avec au moins un vecteur non nul $\Rightarrow$ $\mathrm{Im}(f)$ est la droite vectorielle de base le vecteur non nul retenu ;
\item les deux vecteurs nuls $\Rightarrow$ $\mathrm{Im}(f) = \{\vec{0}\}$.
\end{itemize}
\end{enumerate}
\end{methode}

\begin{exemplebox}[Une image qui est une droite]
Soit $f$ l'application linéaire définie par $f(x\,;y) = (x + 2y\,;\; 3x + 6y)$.
\begin{itemize}
\item $f(\vec{\imath}) = f(1\,;0) = (1\,;3)$ et $f(\vec{\jmath}) = f(0\,;1) = (2\,;6)$.
\item On remarque que $f(\vec{\jmath}) = 2\,f(\vec{\imath})$ : les deux vecteurs sont colinéaires (vérification par le déterminant : $1 \times 6 - 3 \times 2 = 0$).
\item Donc $\mathrm{Im}(f) = \mathrm{Vect}\bigl((1\,;3)\bigr)$ : c'est la droite vectorielle passant par l'origine et dirigée par le vecteur $(1\,;3)$, c'est-à-dire la droite d'équation $3x - y = 0$. Une base de $\mathrm{Im}(f)$ est $\bigl((1\,;3)\bigr)$.
\end{itemize}
Toute l'information du plan de départ s'est \og écrasée \fg{} sur une seule droite : c'est le phénomène de projection que l'on observe ici.
\end{exemplebox}

\begin{popfigure}
\begin{tikzpicture}[scale=1.05, >=Stealth]
% Plan de départ
\begin{scope}[shift={(0,0)}]
\draw[popline, ->] (-0.4,0) -- (2.6,0) ;
\draw[popline, ->] (0,-0.4) -- (0,2.6) ;
\draw[->, very thick, popblue] (0,0) -- (1,0) node[below] {$\vec{\imath}$};
\draw[->, very thick, popblue] (0,0) -- (0,1) node[left] {$\vec{\jmath}$};
\draw[->, thick, poporange] (0,0) -- (1.6,1.2) node[above right] {$\vec{u}=x\vec{\imath}+y\vec{\jmath}$};
\node[popdark] at (1.2,-0.9) {\textbf{Plan de départ}};
\end{scope}
% Flèche f
\draw[->, very thick, poppurple] (3.1,1.1) -- (4.4,1.1) node[midway, above] {$f$};
% Plan d'arrivée
\begin{scope}[shift={(4.9,0)}]
\draw[popline, ->] (-0.4,0) -- (3.1,0) ;
\draw[popline, ->] (0,-0.4) -- (0,2.6) ;
% droite Im(f) : direction (1;3) réduite
\draw[very thick, popgreen, dashed] (-0.35,-1.05) -- (0.85,2.55) ;
\draw[->, very thick, popgreen] (0,0) -- (0.5,1.5) node[right] {$f(\vec{\imath})$};
\draw[->, very thick, popgreen!70!popdark] (0,0) -- (1,3) ;
\node[popgreen!70!popdark] at (1.35,2.9) {$f(\vec{\jmath})=2f(\vec{\imath})$};
\draw[->, thick, poporange] (0,0) -- (0.8,2.4) node[above left] {$f(\vec{u})$};
\node[popgreen] at (2.35,1.9) {$\mathrm{Im}(f)$ : droite};
\node[popdark] at (1.3,-0.9) {\textbf{Plan d'arrivée}};
\end{scope}
\end{tikzpicture}
\legende{L'application $f(x\,;y) = (x+2y\,;3x+6y)$ envoie tout le plan sur une seule droite : $f(\vec{\imath})$ et $f(\vec{\jmath})$ sont colinéaires, donc ils n'engendrent qu'une droite vectorielle, qui est $\mathrm{Im}(f)$.}
\end{popfigure}

\courssub{Stabilité du noyau et de l'image}

Revenons un instant sur le noyau pour établir une propriété commune aux deux ensembles : la \cle{stabilité}. Dire qu'un ensemble est stable pour une opération signifie que cette opération, appliquée à des éléments de l'ensemble, ne fait jamais \og sortir \fg{} de l'ensemble. C'est la propriété qui donne au noyau et à l'image leur statut de \emph{sous-espaces vectoriels}.

\begin{propriete}[Stabilité du noyau]
Soit $f$ une application linéaire du plan. Le noyau $\mathrm{Ker}(f)$ est stable pour l'addition et pour la multiplication par un réel :
\begin{enumerate}
\item si $\vec{u} \in \mathrm{Ker}(f)$ et $\vec{v} \in \mathrm{Ker}(f)$, alors $\vec{u} + \vec{v} \in \mathrm{Ker}(f)$ ;
\item si $\vec{u} \in \mathrm{Ker}(f)$ et $\lambda \in \mathbb{R}$, alors $\lambda\,\vec{u} \in \mathrm{Ker}(f)$.
\end{enumerate}
De plus, $\vec{0} \in \mathrm{Ker}(f)$ car $f(\vec{0}) = \vec{0}$.
\end{propriete}

\begin{experience}[Démonstration (courte et exigible)]
\begin{enumerate}
\item Soient $\vec{u}, \vec{v} \in \mathrm{Ker}(f)$ : on a $f(\vec{u}) = \vec{0}$ et $f(\vec{v}) = \vec{0}$. Par linéarité :
\[ f(\vec{u} + \vec{v}) = f(\vec{u}) + f(\vec{v}) = \vec{0} + \vec{0} = \vec{0}, \]
donc $\vec{u} + \vec{v} \in \mathrm{Ker}(f)$.
\item Soient $\vec{u} \in \mathrm{Ker}(f)$ et $\lambda \in \mathbb{R}$. Par linéarité :
\[ f(\lambda\,\vec{u}) = \lambda\,f(\vec{u}) = \lambda\,\vec{0} = \vec{0}, \]
donc $\lambda\,\vec{u} \in \mathrm{Ker}(f)$. \hfill$\blacksquare$
\end{enumerate}
\end{experience}

\begin{propriete}[Stabilité de l'image]
Soit $f$ une application linéaire du plan. L'image $\mathrm{Im}(f)$ est stable pour l'addition et pour la multiplication par un réel :
\begin{enumerate}
\item si $\vec{a} \in \mathrm{Im}(f)$ et $\vec{b} \in \mathrm{Im}(f)$, alors $\vec{a} + \vec{b} \in \mathrm{Im}(f)$ ;
\item si $\vec{a} \in \mathrm{Im}(f)$ et $\lambda \in \mathbb{R}$, alors $\lambda\,\vec{a} \in \mathrm{Im}(f)$.
\end{enumerate}
De plus, $\vec{0} \in \mathrm{Im}(f)$ car $\vec{0} = f(\vec{0})$.
\end{propriete}

\begin{experience}[Démonstration (courte et exigible)]
\begin{enumerate}
\item Soient $\vec{a}, \vec{b} \in \mathrm{Im}(f)$. Par définition de l'image, il existe $\vec{u}$ et $\vec{v}$ tels que $\vec{a} = f(\vec{u})$ et $\vec{b} = f(\vec{v})$. Alors, par linéarité :
\[ \vec{a} + \vec{b} = f(\vec{u}) + f(\vec{v}) = f(\vec{u} + \vec{v}), \]
donc $\vec{a} + \vec{b}$ est l'image de $\vec{u} + \vec{v}$ : il appartient à $\mathrm{Im}(f)$.
\item Soient $\vec{a} \in \mathrm{Im}(f)$ et $\lambda \in \mathbb{R}$. Il existe $\vec{u}$ tel que $\vec{a} = f(\vec{u})$. Alors :
\[ \lambda\,\vec{a} = \lambda\,f(\vec{u}) = f(\lambda\,\vec{u}), \]
donc $\lambda\,\vec{a}$ est l'image de $\lambda\,\vec{u}$ : il appartient à $\mathrm{Im}(f)$. \hfill$\blacksquare$
\end{enumerate}
\end{experience}

\begin{aretenir}[Noyau et image sont des sous-espaces vectoriels]
Un sous-ensemble du plan contenant $\vec{0}$ et stable pour l'addition et la multiplication par un réel s'appelle un \cle{sous-espace vectoriel} du plan (vocabulaire admis). Ainsi :
\[ \mathrm{Ker}(f) \text{ et } \mathrm{Im}(f) \text{ sont des sous-espaces vectoriels du plan.} \]
Dans le plan, les seuls sous-espaces vectoriels sont : $\{\vec{0}\}$, les droites vectorielles, et le plan entier. C'est pourquoi noyau et image ne peuvent prendre que ces trois formes.
\end{aretenir}

\begin{attention}[La stabilité se démontre, elle ne s'affirme pas]
Au Bac, lorsqu'on te demande de justifier qu'une somme ou qu'un produit par un réel reste dans le noyau ou dans l'image, il faut rédiger la démonstration à partir de la linéarité de $f$, comme ci-dessus. Écrire simplement \og c'est stable par propriété du cours \fg{} sans citer $f(\vec{u}+\vec{v}) = f(\vec{u}) + f(\vec{v})$ fait perdre des points de rigueur.
\end{attention}

\courssub{Image et surjectivité}

Rappelons qu'une application est \cle{surjective} lorsque tout vecteur du plan d'arrivée possède au moins un antécédent. Or \og posséder un antécédent \fg{}, c'est précisément appartenir à $\mathrm{Im}(f)$. Le lien est donc immédiat.

\begin{propriete}[Caractérisation de la surjectivité par l'image]
Soit $f$ une application linéaire du plan dans lui-même. Alors :
\[ f \text{ est surjective} \quad \Longleftrightarrow \quad \mathrm{Im}(f) = \mathcal{P} \quad \Longleftrightarrow \quad \bigl(f(\vec{\imath}), f(\vec{\jmath})\bigr) \text{ est une base du plan.} \]
Autrement dit, $f$ est surjective si et seulement si son rang vaut $2$.
\end{propriete}

\begin{exoresolu}[Image, stabilité et surjectivité]
Soit $f$ l'application linéaire définie par $f(x\,;y) = (2x - y\,;\; -4x + 2y)$.
\begin{enumerate}
\item Déterminer $\mathrm{Im}(f)$ et en donner une base. L'application $f$ est-elle surjective ?
\item Le vecteur $\vec{a} = (1\,;-2)$ appartient-il à $\mathrm{Im}(f)$ ? Le vecteur $\vec{b} = (1\,;1)$ ?
\item Montrer que si $\vec{a} \in \mathrm{Im}(f)$, alors $-5\,\vec{a} \in \mathrm{Im}(f)$.
\end{enumerate}
\tcblower
\textbf{Solution détaillée.}
\begin{enumerate}
\item On calcule les images des vecteurs de base :
\[ f(\vec{\imath}) = f(1\,;0) = (2\,;-4), \qquad f(\vec{\jmath}) = f(0\,;1) = (-1\,;2). \]
Déterminant : $2 \times 2 - (-4) \times (-1) = 4 - 4 = 0$ : les vecteurs sont colinéaires (on voit d'ailleurs que $f(\vec{\imath}) = -2\,f(\vec{\jmath})$). Donc
\[ \mathrm{Im}(f) = \mathrm{Vect}\bigl((-1\,;2)\bigr), \]
c'est la droite vectorielle de base $\bigl((-1\,;2)\bigr)$, d'équation $2x + y = 0$. Comme $\mathrm{Im}(f) \neq \mathcal{P}$, l'application $f$ \textbf{n'est pas surjective}.
\item $\vec{a} = (1\,;-2) = -1 \times (-1\,;2)$ : c'est un multiple du vecteur de base, donc $\vec{a} \in \mathrm{Im}(f)$. (Vérification : $f(0\,;-1) = (1\,;-2) = \vec{a}$.) En revanche, $\vec{b} = (1\,;1)$ ne vérifie pas $2x + y = 0$ (car $2 + 1 = 3 \neq 0$), donc $\vec{b} \notin \mathrm{Im}(f)$ : $\vec{b}$ n'a aucun antécédent par $f$.
\item Comme $\vec{a} \in \mathrm{Im}(f)$, il existe $\vec{u}$ tel que $\vec{a} = f(\vec{u})$ (ici $\vec{u} = (0\,;-1)$). Par linéarité :
\[ -5\,\vec{a} = -5\,f(\vec{u}) = f(-5\,\vec{u}), \]
donc $-5\,\vec{a}$ est l'image de $-5\,\vec{u}$ : il appartient à $\mathrm{Im}(f)$. C'est exactement la stabilité de l'image par multiplication par un réel.
\end{enumerate}
\end{exoresolu}

\begin{popfigure}
\begin{tikzpicture}[node distance=0.35cm, >=Stealth,
  cas/.style={draw=popdark, thick, rounded corners=3pt, fill=popcream, text width=3.1cm, align=center, minimum height=0.9cm, font=\small},
  kerim/.style={draw, thick, rounded corners=3pt, text width=3.1cm, align=center, font=\small}]
\node[cas, align=center] (c0) {\textbf{Rang 0}\\ $f$ application nulle};
\node[cas, below=0.7cm of c0, align=center] (c1) {\textbf{Rang 1}\\ $\det\bigl(f(\vec{\imath}),f(\vec{\jmath})\bigr)=0$};
\node[cas, below=0.7cm of c1, align=center] (c2) {\textbf{Rang 2}\\ $\det\bigl(f(\vec{\imath}),f(\vec{\jmath})\bigr)\neq 0$};
\node[kerim, right=1.6cm of c0, draw=popblue, fill=popblueL, align=center] (k0) {$\mathrm{Ker}(f) = \mathcal{P}$\\ (plan entier)};
\node[kerim, right=0.5cm of k0, draw=popgreen, fill=popgreenL] (i0) {$\mathrm{Im}(f) = \{\vec{0}\}$};
\node[kerim, right=1.6cm of c1, draw=popblue, fill=popblueL, align=center] (k1) {$\mathrm{Ker}(f)$ :\\ droite vectorielle};
\node[kerim, right=0.5cm of k1, draw=popgreen, fill=popgreenL, align=center] (i1) {$\mathrm{Im}(f)$ :\\ droite vectorielle};
\node[kerim, right=1.6cm of c2, draw=popblue, fill=popblueL, align=center] (k2) {$\mathrm{Ker}(f) = \{\vec{0}\}$\\ ($f$ injective)};
\node[kerim, right=0.5cm of k2, draw=popgreen, fill=popgreenL, align=center] (i2) {$\mathrm{Im}(f) = \mathcal{P}$\\ ($f$ surjective)};
\draw[->, popmuted, thick] (c0) -- (k0); \draw[->, popmuted, thick] (k0) -- (i0);
\draw[->, popmuted, thick] (c1) -- (k1); \draw[->, popmuted, thick] (k1) -- (i1);
\draw[->, popmuted, thick] (c2) -- (k2); \draw[->, popmuted, thick] (k2) -- (i2);
\node[poporange, font=\small\bfseries, below=0.25cm of i2, text width=3.1cm, align=center] {Cas bijectif};
\end{tikzpicture}
\legende{Les trois configurations possibles du couple $\bigl(\mathrm{Ker}(f), \mathrm{Im}(f)\bigr)$ selon le rang de $f$. Remarque le phénomène de balance : plus le noyau est gros, plus l'image est petite, et réciproquement.}
\end{popfigure}

\begin{savaistu}[La balance noyau--image]
Observe la figure récapitulative : quand le noyau \og grossit \fg{}, l'image \og maigrit \fg{}, et inversement. Ce n'est pas un hasard : une célèbre formule, le \emph{théorème du rang} (que tu rencontreras à l'université), affirme que la dimension du noyau plus la dimension de l'image est toujours égale à la dimension de l'espace de départ, ici $2$. C'est une loi de conservation, comme en physique : ce que l'application \og perd \fg{} dans le noyau, elle le retrouve exactement dans son image. Les ingénieurs de CAMTEL ou de MTN qui compressent des signaux numériques utilisent précisément cette idée : comprendre ce qu'une transformation perd (le noyau) et ce qu'elle conserve (l'image).
\end{savaistu}

\begin{exercice}[À toi de jouer]
Soit $g$ l'application linéaire définie par $g(x\,;y) = (x - 3y\,;\; 2x + y)$.
\begin{enumerate}
\item Calculer $g(\vec{\imath})$ et $g(\vec{\jmath})$.
\item Montrer que $\mathrm{Im}(g) = \mathcal{P}$. Que peut-on en déduire pour $g$ ?
\item Montrer que la somme de deux vecteurs de $\mathrm{Im}(g)$ appartient encore à $\mathrm{Im}(g)$, en rédigeant la démonstration complète.
\end{enumerate}
\end{exercice}

\begin{corrige}[Exercice : image de $g$]
\begin{enumerate}
\item $g(\vec{\imath}) = g(1\,;0) = (1\,;2)$ et $g(\vec{\jmath}) = g(0\,;1) = (-3\,;1)$.
\item Déterminant : $1 \times 1 - 2 \times (-3) = 1 + 6 = 7 \neq 0$. Les vecteurs $g(\vec{\imath})$ et $g(\vec{\jmath})$ ne sont pas colinéaires : ils forment une base du plan, donc $\mathrm{Im}(g) = \mathrm{Vect}\bigl(g(\vec{\imath}), g(\vec{\jmath})\bigr) = \mathcal{P}$. On en déduit que $g$ est surjective.
\item Soient $\vec{a}, \vec{b} \in \mathrm{Im}(g)$. Il existe $\vec{u}$ et $\vec{v}$ tels que $\vec{a} = g(\vec{u})$ et $\vec{b} = g(\vec{v})$. Par linéarité de $g$ : $\vec{a} + \vec{b} = g(\vec{u}) + g(\vec{v}) = g(\vec{u} + \vec{v})$, donc $\vec{a} + \vec{b} \in \mathrm{Im}(g)$. (Ici, c'était d'ailleurs évident puisque $\mathrm{Im}(g) = \mathcal{P}$, mais la démonstration reste valable pour toute application linéaire.)
\end{enumerate}
\end{corrige}

En résumé : l'image de $f$ est engendrée par les images des vecteurs de base, elle est stable pour l'addition et la multiplication par un réel (c'est un sous-espace vectoriel), et sa comparaison avec le plan entier caractérise la surjectivité. Il ne nous reste plus, dans la section suivante, qu'à assembler noyau, image et écriture analytique pour obtenir une caractérisation complète de la bijectivité.

\coursec{Caractérisation de la bijectivité et synthèse}

Dans la section précédente, tu as appris à déterminer l'\emph{image} d'une application linéaire du plan : c'est soit le plan tout entier, soit une droite vectorielle, soit le singleton $\{\vec{0}\}$. Tu as aussi vu, dans la section 4, que le \emph{noyau} connaît exactement les trois mêmes destins. Une question naturelle s'impose alors : \emph{existe-t-il un lien entre le noyau et l'image ? Et comment savoir, au premier coup d'œil, si une application linéaire est bijective ?} C'est toute l'ambition de cette section : découvrir que, pour un endomorphisme du plan, \textbf{tous les critères de bijectivité se rejoignent} et se résument en un seul nombre : $ad - bc$.

\courssub{Le théorème central : six façons de dire « bijective »}

\textbf{Intuition.} Imagine une machine à trier les cartes SIM dans un centre MTN : chaque client (vecteur d'arrivée) doit recevoir exactement une carte (vecteur de départ). Pour que le service soit parfait, il faut deux choses : aucun client ne doit être oublié (\emph{surjectivité}) et deux clients distincts ne doivent jamais recevoir la même carte (\emph{injectivité}). Le miracle des applications linéaires du plan, c'est que ces deux conditions, en apparence indépendantes, sont en réalité \emph{la même condition}, et qu'elle se lit directement sur les coefficients de l'écriture analytique.

\begin{propriete}[Théorème de caractérisation de la bijectivité (exigible au Bac)]
Soit $f$ un endomorphisme du plan vectoriel $E$, d'écriture analytique
\[ f(x\,;\,y) = (ax + by\,;\,cx + dy) \]
dans une base $(\vec{i}\,;\,\vec{j})$ de $E$. Les six propriétés suivantes sont \cle{équivalentes} :
\begin{enumerate}
\item $f$ est \cle{bijective} ;
\item $f$ est \cle{injective} ;
\item $f$ est \cle{surjective} ;
\item $\ker(f) = \{\vec{0}\}$ ;
\item $\mathrm{Im}(f) = E$ ;
\item $ad - bc \neq 0$.
\end{enumerate}
Dans ce cas, $f$ est un \cle{automorphisme} de $E$, et sa réciproque $f^{-1}$ est aussi une application linéaire.
\end{propriete}

\begin{experience}[Démonstration guidée du théorème]
Nous allons établir les équivalences en trois temps. Complète mentalement chaque étape avant de lire la suite.

\textbf{Étape 1 : $f$ injective $\iff \ker(f) = \{\vec{0}\}$.}
\begin{itemize}
\item Supposons $f$ injective. Soit $\vec{u} \in \ker(f)$ : alors $f(\vec{u}) = \vec{0} = f(\vec{0})$, donc par injectivité $\vec{u} = \vec{0}$. Ainsi $\ker(f) = \{\vec{0}\}$.
\item Réciproquement, supposons $\ker(f) = \{\vec{0}\}$. Si $f(\vec{u}) = f(\vec{v})$, alors par linéarité $f(\vec{u} - \vec{v}) = \vec{0}$, donc $\vec{u} - \vec{v} \in \ker(f) = \{\vec{0}\}$, d'où $\vec{u} = \vec{v}$ : $f$ est injective.
\end{itemize}

\textbf{Étape 2 : $f$ surjective $\iff \mathrm{Im}(f) = E$.} C'est la définition même de la surjectivité : tout vecteur de $E$ possède un antécédent.

\textbf{Étape 3 : le lien avec $ad - bc$.} Notons $\vec{u} = f(\vec{i}) = (a\,;\,c)$ et $\vec{v} = f(\vec{j}) = (b\,;\,d)$. On sait (section 5) que $\mathrm{Im}(f) = \mathrm{Vect}(\vec{u}\,;\,\vec{v})$.
\begin{itemize}
\item Si $ad - bc \neq 0$, les vecteurs $\vec{u}$ et $\vec{v}$ ne sont pas colinéaires (le critère de colinéarité $ad - bc = 0$ n'est pas vérifié). Ils forment donc une base de $E$, et $\mathrm{Im}(f) = E$ : $f$ est surjective.
\item Si $ad - bc = 0$, alors $\vec{u}$ et $\vec{v}$ sont colinéaires, donc $\mathrm{Im}(f)$ est au plus une droite : $f$ n'est pas surjective.
\end{itemize}

\textbf{Étape 4 : dans le plan, injectivité $\iff$ surjectivité.} Si $f$ n'est pas injective, alors $\ker(f) \neq \{\vec{0}\}$ : il existe $\vec{w} \neq \vec{0}$ avec $f(\vec{w}) = \vec{0}$. Or $\vec{w}$ est combinaison de $\vec{i}$ et $\vec{j}$, donc $\vec{0} = f(\vec{w})$ est une combinaison non triviale de $\vec{u}$ et $\vec{v}$ : ceux-ci sont colinéaires, donc $f$ n'est pas surjective. En passant aux contraposées, surjectivité entraîne injectivité, et le raisonnement symétrique donne l'autre sens.

\textbf{Conclusion.} Les propriétés 2, 3, 4, 5, 6 sont toutes équivalentes entre elles ; comme bijective signifie « injective \emph{et} surjective », la propriété 1 leur est équivalente aussi. $\blacksquare$
\end{experience}

\begin{popfigure}
\begin{center}
\begin{tikzpicture}[
  node distance=1.1cm and 1.6cm,
  box/.style={draw=popblue, thick, rounded corners=4pt, fill=popblueL, text=popdark, minimum height=0.9cm, inner sep=6pt, font=\small\bfseries},
  eq/.style={<->, >=Stealth, thick, poporange}
]
\node[box] (bij) {$f$ bijective};
\node[box, right=of bij] (inj) {$f$ injective};
\node[box, right=of inj] (surj) {$f$ surjective};
\node[box, below=of bij] (ker) {$\ker(f) = \{\vec{0}\}$};
\node[box, below=of inj] (im) {$\mathrm{Im}(f) = E$};
\node[box, below=of surj] (det) {$ad - bc \neq 0$};
\draw[eq] (bij) -- (inj);
\draw[eq] (inj) -- (surj);
\draw[eq] (bij) -- (ker);
\draw[eq] (inj) -- (im);
\draw[eq] (surj) -- (det);
\draw[eq] (ker) -- (im);
\draw[eq] (im) -- (det);
\draw[eq, dashed] (ker.east) .. controls +(1.2,-0.8) and +(-1.2,-0.8) .. (det.west);
\end{tikzpicture}
\end{center}
\legende{Le graphe des équivalences : chaque double flèche signifie « si et seulement si ». Une seule propriété vérifiée suffit à les obtenir toutes.}
\end{popfigure}

\courssub{Le déterminant $ad - bc$ : un critère analytique éclair}

\begin{definition}[Déterminant d'un endomorphisme du plan]
Soit $f$ l'endomorphisme d'écriture analytique $f(x\,;\,y) = (ax + by\,;\,cx + dy)$. Le nombre réel
\[ ad - bc \]
est appelé \cle{déterminant} de $f$ (dans la base considérée). On le note souvent $\det(f)$.
\end{definition}

\begin{aretenir}[Le réflexe à acquérir]
Pour un endomorphisme du plan $f(x\,;\,y) = (ax + by\,;\,cx + dy)$ :
\[ f \text{ est un automorphisme} \iff ad - bc \neq 0. \]
Si $ad - bc = 0$, alors $f$ n'est ni injective ni surjective : son noyau et son image sont des droites vectorielles (sauf si $f$ est l'application nulle).
\end{aretenir}

\begin{exemplebox}[Un calcul qui décide de tout]
Soit $f$ l'application linéaire définie par $f(x\,;\,y) = (3x - y\,;\,2x + y)$. Ici $a = 3$, $b = -1$, $c = 2$, $d = 1$. On calcule :
\[ ad - bc = 3 \times 1 - (-1) \times 2 = 3 + 2 = 5 \neq 0. \]
Sans résoudre aucune équation, on conclut : $f$ est un \textbf{automorphisme} du plan. Son noyau est $\{\vec{0}\}$, son image est le plan entier, et elle admet une réciproque linéaire $f^{-1}$.
\end{exemplebox}

\begin{attention}[Le piège classique : « $f(\vec{0}) = \vec{0}$ donc $f$ est bijective »]
C'est une erreur très fréquente au Bac. \textbf{Toute} application linéaire vérifie $f(\vec{0}) = \vec{0}$, y compris les pires : l'application nulle $f(x\,;\,y) = (0\,;\,0)$ satisfait $f(\vec{0}) = \vec{0}$ et n'a pourtant rien d'injectif ! Vérifier $f(\vec{0}) = \vec{0}$ ne prouve \emph{rien}. Pour l'injectivité, il faut montrer que $\ker(f) = \{\vec{0}\}$, c'est-à-dire résoudre le système $f(x\,;\,y) = (0\,;\,0)$ et vérifier que $(0\,;\,0)$ en est l'\emph{unique} solution.
\end{attention}

\begin{savaistu}[Le déterminant mesure les aires]
Le nombre $ad - bc$ ne dit pas seulement si $f$ est bijective : sa valeur absolue $|ad - bc|$ est le \emph{facteur par lequel $f$ multiplie les aires}. Le carré unité construit sur $\vec{i}$ et $\vec{j}$ (aire $1$) est transformé en un parallélogramme d'aire $|ad - bc|$. Ainsi, l'application $f(x\,;\,y) = (3x - y\,;\,2x + y)$ multiplie toutes les aires par $5$ : un terrain de \SI{200}{m^2} à Yaoundé, vu à travers $f$, deviendrait un terrain de \SI{1000}{m^2}. Et si $ad - bc = 0$, le parallélogramme est aplati sur une droite : aire nulle, information perdue, bijectivité impossible. C'est la même idée qui fait fonctionner le codage des images et la cartographie numérique.
\end{savaistu}

\courssub{La réciproque d'un automorphisme}

\begin{propriete}[Réciproque d'un automorphisme]
Si $f$ est un automorphisme du plan, alors sa bijection réciproque $f^{-1}$ est également une application \cle{linéaire}. Pour obtenir son écriture analytique, on résout le système
\[ \begin{cases} x' = ax + by \\ y' = cx + dy \end{cases} \]
en exprimant $x$ et $y$ en fonction de $x'$ et $y'$ (ce qui est possible précisément parce que $ad - bc \neq 0$).
\end{propriete}

\begin{exemplebox}[Calcul d'une réciproque]
Reprenons $f(x\,;\,y) = (3x - y\,;\,2x + y)$. Posons $\begin{cases} x' = 3x - y \\ y' = 2x + y \end{cases}$. En additionnant : $x' + y' = 5x$, donc $x = \dfrac{x' + y'}{5}$. Puis $y = 3x - x' = \dfrac{3x' + 3y'}{5} - x' = \dfrac{-2x' + 3y'}{5}$. D'où :
\[ f^{-1}(x\,;\,y) = \left(\frac{x + y}{5}\,;\,\frac{-2x + 3y}{5}\right). \]
Vérification sur un vecteur : $f(1\,;\,2) = (1\,;\,4)$ et $f^{-1}(1\,;\,4) = \left(\dfrac{5}{5}\,;\,\dfrac{10}{5}\right) = (1\,;\,2)$. Parfait.
\end{exemplebox}

\courssub{Tableau de synthèse : les trois destins d'un endomorphisme du plan}

\begin{popfigure}
\begin{center}
\begin{tikzpicture}[scale=0.62, every node/.style={transform shape}]
% Cas 1 : bijectif
\begin{scope}
\draw[->, popmuted] (-2.2,0) -- (2.2,0); \draw[->, popmuted] (0,-2.2) -- (0,2.2);
\draw[popblue, thick, ->] (0,0) -- (1,0); \draw[popblue, thick, ->] (0,0) -- (0,1);
\node[popdark, font=\small\bfseries] at (0,-2.8) {$\det f \neq 0$ : bijective};
\node[popmuted, font=\scriptsize] at (0,2.6) {plan $\to$ plan};
\end{scope}
% Cas 2 : rang 1
\begin{scope}[xshift=7.5cm]
\draw[->, popmuted] (-2.2,0) -- (2.2,0); \draw[->, popmuted] (0,-2.2) -- (0,2.2);
\draw[poporange, very thick] (-1.8,-0.9) -- (1.8,0.9);
\fill[poporange] (1,0.5) circle (2pt); \fill[poporange] (-1,-0.5) circle (2pt);
\node[popdark, font=\small\bfseries] at (0,-2.8) {$\det f = 0$, $f \neq 0$ : rang $1$};
\node[popmuted, font=\scriptsize] at (0,2.6) {plan $\to$ droite};
\end{scope}
% Cas 3 : nulle
\begin{scope}[xshift=15cm]
\draw[->, popmuted] (-2.2,0) -- (2.2,0); \draw[->, popmuted] (0,-2.2) -- (0,2.2);
\fill[poppink] (0,0) circle (3pt);
\node[popdark, font=\small\bfseries] at (0,-2.8) {$f = 0$ : application nulle};
\node[popmuted, font=\scriptsize] at (0,2.6) {plan $\to$ point};
\end{scope}
\end{tikzpicture}
\end{center}
\legende{Les trois situations possibles pour l'image d'un endomorphisme du plan : le plan entier, une droite, ou un point.}
\end{popfigure}

\begin{center}
\begin{tabularx}{\linewidth}{|l|X|X|X|}
\hline
\rowcolor{popblueL} \textbf{Critère} & \textbf{Automorphisme} & \textbf{Rang 1} & \textbf{Application nulle} \\
\hline
Déterminant $ad - bc$ & $\neq 0$ & $= 0$ & $= 0$ \\
\hline
Noyau $\ker(f)$ & $\{\vec{0}\}$ & une droite vectorielle & le plan $E$ entier \\
\hline
Image $\mathrm{Im}(f)$ & le plan $E$ entier & une droite vectorielle & $\{\vec{0}\}$ \\
\hline
$f(\vec{i})$ et $f(\vec{j})$ & non colinéaires & colinéaires, non tous nuls & tous les deux nuls \\
\hline
Bijectivité & oui & non & non \\
\hline
\end{tabularx}
\end{center}

Remarque la belle symétrie : le noyau et l'image « échangent » leurs rôles. Quand l'un est réduit à $\{\vec{0}\}$, l'autre est le plan entier ; quand l'un est une droite, l'autre aussi. Ce n'est pas un hasard : c'est un cas particulier du \emph{théorème du rang}, que tu rencontreras dans l'enseignement supérieur.

\courssub{Stratégie d'examen : choisir le bon critère}

\begin{methode}[Prouver qu'une application linéaire est bijective : trois routes au choix]
Face à la question « Montrer que $f$ est bijective », trois stratégies sont possibles. Choisis la plus rapide selon les données :
\begin{enumerate}
\item \textbf{Route analytique (la plus rapide).} Si l'écriture $f(x\,;\,y) = (ax + by\,;\,cx + dy)$ est connue, calcule $ad - bc$. S'il est non nul, conclus immédiatement : « $ad - bc = \ldots \neq 0$, donc $f$ est un automorphisme. »
\item \textbf{Route du noyau.} Résous le système $f(x\,;\,y) = (0\,;\,0)$. Si l'unique solution est $(0\,;\,0)$, alors $\ker(f) = \{\vec{0}\}$, donc $f$ est injective, donc bijective (car on est dans le plan).
\item \textbf{Route géométrique.} Calcule $f(\vec{i})$ et $f(\vec{j})$. S'ils ne sont pas colinéaires, ils forment une base de $E$, donc $\mathrm{Im}(f) = E$ : $f$ est surjective, donc bijective.
\end{enumerate}
\emph{Conseil de rédaction :} cite toujours le théorème (« $f$ est un endomorphisme du plan et $\ker(f) = \{\vec{0}\}$, donc $f$ est bijective ») : c'est ce qui vaut les points de justification au Bac.
\end{methode}

\begin{attention}[Conditions d'application à ne pas oublier]
\begin{itemize}
\item L'équivalence « injective $\iff$ surjective » n'est vraie que parce que $f$ est un \textbf{endomorphisme du plan} (espace de départ = espace d'arrivée, de dimension finie). Ne l'utilise jamais pour une application quelconque.
\item Avant d'appliquer le critère $ad - bc \neq 0$, assure-toi d'avoir \textbf{d'abord justifié que $f$ est linéaire} si ce n'est pas donné dans l'énoncé.
\item Ne confonds pas : $ad - bc = 0$ ne signifie pas que $f$ est nulle ; cela signifie seulement que $f$ n'est pas bijective.
\end{itemize}
\end{attention}

\courssub{Entraînement type Bac : étude complète d'une application linéaire}

\begin{exoresolu}[Étude complète (sujet type Baccalauréat TI)]
On considère l'application $f$ du plan vectoriel $E$ dans lui-même définie par :
\[ f(x\,;\,y) = (2x + 3y\,;\,x - y). \]
\begin{enumerate}
\item Montrer que $f$ est une application linéaire.
\item Déterminer le noyau de $f$.
\item En déduire que $f$ est bijective, puis déterminer $\mathrm{Im}(f)$.
\item Retrouver la bijectivité par le critère analytique.
\item Déterminer l'expression de $f^{-1}$.
\end{enumerate}
\tcblower
\textbf{1. Linéarité.} Soient $\vec{u} = (x\,;\,y)$, $\vec{v} = (x'\,;\,y')$ et $\lambda \in \mathbb{R}$.
\begin{align*}
f(\vec{u} + \vec{v}) &= f(x + x'\,;\,y + y') = \big(2(x+x') + 3(y+y')\,;\,(x+x') - (y+y')\big) \\
&= (2x + 3y\,;\,x - y) + (2x' + 3y'\,;\,x' - y') = f(\vec{u}) + f(\vec{v}).
\end{align*}
De même, $f(\lambda \vec{u}) = (2\lambda x + 3\lambda y\,;\,\lambda x - \lambda y) = \lambda f(\vec{u})$. Donc $f$ est \textbf{linéaire}.

\textbf{2. Noyau.} On résout $f(x\,;\,y) = (0\,;\,0)$ :
\[ \begin{cases} 2x + 3y = 0 \\ x - y = 0 \end{cases} \iff \begin{cases} 2x + 3y = 0 \\ x = y \end{cases} \iff 5y = 0 \iff x = y = 0. \]
L'unique solution est $(0\,;\,0)$, donc $\ker(f) = \{\vec{0}\}$.

\textbf{3. Bijectivité et image.} $f$ est un endomorphisme du plan et $\ker(f) = \{\vec{0}\}$ : d'après le théorème de caractérisation, $f$ est \textbf{injective}, donc \textbf{bijective}. Par conséquent $\mathrm{Im}(f) = E$ : l'image est le plan tout entier.

\textbf{4. Critère analytique.} Avec $a = 2$, $b = 3$, $c = 1$, $d = -1$ :
\[ ad - bc = 2 \times (-1) - 3 \times 1 = -2 - 3 = -5 \neq 0, \]
ce qui confirme que $f$ est un automorphisme.

\textbf{5. Réciproque.} On résout $\begin{cases} x' = 2x + 3y \\ y' = x - y \end{cases}$. De la deuxième équation, $x = y' + y$ ; en substituant : $x' = 2y' + 5y$, donc $y = \dfrac{x' - 2y'}{5}$ et $x = y' + \dfrac{x' - 2y'}{5} = \dfrac{x' + 3y'}{5}$. Ainsi :
\[ f^{-1}(x\,;\,y) = \left(\frac{x + 3y}{5}\,;\,\frac{x - 2y}{5}\right). \]
Vérification : $f^{-1}(f(1\,;\,1)) = f^{-1}(5\,;\,0) = (1\,;\,1)$. $\checkmark$
\end{exoresolu}

\begin{exercice}[À ton tour]
Soit $g$ l'application du plan dans lui-même définie par $g(x\,;\,y) = (x + 2y\,;\,3x + 6y)$.
\begin{enumerate}
\item Justifier que $g$ est linéaire.
\item Calculer son déterminant. Que peut-on en conclure ?
\item Déterminer $\ker(g)$ et $\mathrm{Im}(g)$ : on donnera pour chacun une équation et un vecteur directeur.
\item $g$ est-elle bijective ? Justifier par le noyau.
\end{enumerate}
\end{exercice}

\begin{corrige}[Exercice : application $g$]
\textbf{1.} Pour $\vec{u} = (x\,;\,y)$, $\vec{v} = (x'\,;\,y')$, $\lambda \in \mathbb{R}$ : $g(\vec{u} + \vec{v}) = \big((x+x') + 2(y+y')\,;\,3(x+x') + 6(y+y')\big) = g(\vec{u}) + g(\vec{v})$ et $g(\lambda \vec{u}) = \lambda g(\vec{u})$. Donc $g$ est linéaire.

\textbf{2.} $\det(g) = 1 \times 6 - 2 \times 3 = 6 - 6 = 0$. Donc $g$ n'est \textbf{pas bijective} : ce n'est pas un automorphisme.

\textbf{3.} Noyau : $g(x\,;\,y) = (0\,;\,0) \iff x + 2y = 0$ (la deuxième équation $3x + 6y = 0$ est la même multipliée par $3$). Donc $\ker(g)$ est la droite d'équation $x + 2y = 0$, dirigée par $\vec{u} = (-2\,;\,1)$. Image : $g(\vec{i}) = (1\,;\,3)$ et $g(\vec{j}) = (2\,;\,6) = 2\,g(\vec{i})$, donc $\mathrm{Im}(g) = \mathrm{Vect}(1\,;\,3)$, la droite d'équation $3x - y = 0$.

\textbf{4.} $\ker(g) \neq \{\vec{0}\}$ (par exemple $g(-2\,;\,1) = (0\,;\,0)$ avec $(-2\,;\,1) \neq (0\,;\,0)$), donc $g$ n'est pas injective, donc pas bijective. On retrouve la situation « rang 1 » du tableau de synthèse : noyau et image sont tous deux des droites vectorielles.
\end{corrige}

\begin{aretenir}[L'essentiel de la section]
\begin{itemize}
\item Pour un endomorphisme du plan : \textbf{bijective $\iff$ injective $\iff$ surjective $\iff$ $\ker(f) = \{\vec{0}\}$ $\iff$ $\mathrm{Im}(f) = E$ $\iff$ $ad - bc \neq 0$}.
\item Le nombre $ad - bc$ est le \cle{déterminant} de $f$ : c'est le critère le plus rapide, et sa valeur absolue donne le facteur d'agrandissement des aires.
\item Si $f$ est bijective, c'est un \cle{automorphisme} et $f^{-1}$ est aussi linéaire ; on l'obtient en résolvant le système $\begin{cases} x' = ax + by \\ y' = cx + dy \end{cases}$.
\item Trois situations possibles : automorphisme (plan $\to$ plan), rang 1 (plan $\to$ droite, noyau = droite), application nulle (plan $\to$ point).
\item Piège mortel : $f(\vec{0}) = \vec{0}$ ne prouve jamais la bijectivité.
\end{itemize}
\end{aretenir}

Tu disposes maintenant de tous les outils pour étudier n'importe quelle application linéaire du plan, de sa définition jusqu'à sa bijectivité. Dans les leçons suivantes, tu verras comment ces outils s'appliquent aux transformations géométriques remarquables : symétries, rotations, homothéties et projections, qui sont toutes des applications linéaires particulières.

\newpage

\coursec{Activité d'intégration}

\coursec{Applications linéaires du plan}

\courssub{Synthèse des savoirs indispensables}

\begin{aretenir}[Boîte à outils de la leçon M19]
\begin{itemize}
\item \textbf{Définition.} Une application $f$ du plan vectoriel $\mathcal{P}$ dans lui-même est \cle{linéaire} si et seulement si, pour tous vecteurs $\vec{u}, \vec{v} \in \mathcal{P}$ et tout réel $\lambda \in \mathbb{R}$ :
\[ f(\vec{u}+\vec{v}) = f(\vec{u}) + f(\vec{v}) \qquad \text{et} \qquad f(\lambda\vec{u}) = \lambda f(\vec{u}). \]
Conséquence immédiate : $f(\vec{0}) = \vec{0}$. Toute application ne respectant pas cette condition (terme constant non nul) n'est pas linéaire (c'est une application affine).
\item \textbf{Endomorphisme, automorphisme, isomorphisme.} Une application linéaire d'un espace vectoriel dans lui-même est un \cle{endomorphisme}. Si elle est bijective, c'est un \cle{automorphisme}. Une application linéaire bijective entre deux espaces vectoriels (éventuellement différents) est un \cle{isomorphisme}.
\item \textbf{Matrice dans la base canonique.} Si $f(x\,;y) = (ax+by\,;\,cx+dy)$ dans la base canonique $(\vec{\imath},\vec{\jmath})$, alors $f(\vec{\imath}) = (a\,;c)$, $f(\vec{\jmath}) = (b\,;d)$ et la matrice de $f$ est $M = \begin{pmatrix} a & b \\ c & d \end{pmatrix}$.
\item \textbf{Noyau.} $\ker f = \{\vec{u} \in \mathcal{P} \mid f(\vec{u}) = \vec{0}\}$. On le détermine en résolvant le système homogène $\begin{cases} ax+by = 0 \\ cx+dy = 0 \end{cases}$. C'est un sous-espace vectoriel (stable par addition et multiplication par un réel).
\item \textbf{Image.} $\operatorname{Im} f = \{f(\vec{u}) \mid \vec{u} \in \mathcal{P}\}$. C'est le sous-espace vectoriel engendré par $f(\vec{\imath})$ et $f(\vec{\jmath})$ (les colonnes de la matrice), lui aussi stable par addition et multiplication par un réel.
\item \textbf{Caractérisation de la bijectivité (Théorème fondamental).} Pour un endomorphisme $f$ du plan, les propriétés suivantes sont \textbf{équivalentes} :
\[ f \text{ bijective} \iff \ker f = \{\vec{0}\} \iff \operatorname{Im} f = \mathcal{P} \iff \det M = ad - bc \neq 0. \]
\item \textbf{Bijection réciproque.} Si $f$ est bijective, $f^{-1}$ s'obtient en résolvant par rapport à $(x\,;y)$ le système $\begin{cases} x' = ax+by \\ y' = cx+dy \end{cases}$. Sa matrice est $M^{-1} = \frac{1}{ad-bc}\begin{pmatrix} d & -b \\ -c & a \end{pmatrix}$.
\item \textbf{Théorème du rang (dimension finie).} $\dim(\ker f) + \dim(\operatorname{Im} f) = \dim(\mathcal{P}) = 2$.
\end{itemize}
\end{aretenir}

\begin{attention}[Pièges à éviter au Bac]
\begin{itemize}
\item \textbf{Linéarité :} On ne prouve pas la linéarité sur un exemple. Il faut vérifier les deux propriétés pour des vecteurs \emph{quelconques} $(x_1,y_1), (x_2,y_2)$ et $\lambda \in \mathbb{R}$. En revanche, \textbf{un seul contre-exemple} (souvent $f(\vec{0}) \neq \vec{0}$) suffit pour dire « non linéaire ».
\item \textbf{Noyau :} $\ker f$ n'est \textbf{jamais} vide : il contient toujours $\vec{0}$. On écrit $\ker f = \{\vec{0}\}$ (et pas $\varnothing$).
\item \textbf{Image :} Ne pas écrire « l'image est le plan » ou « l'image est une droite » sans justification. Donner une \textbf{base} (famille génératrice libre) ou une \textbf{équation caractéristique} (si c'est une droite).
\item \textbf{Bijectivité :} Au Bac, on exige souvent de justifier par \textbf{deux méthodes différentes} (ex: noyau trivial ET déterminant non nul).
\item \textbf{Notations :} Bien distinguer le vecteur $\vec{u}$ de ses coordonnées $(x,y)$ dans une base. La matrice s'applique sur la matrice-colonne des coordonnées.
\end{itemize}
\end{attention}

\courssub{Activité d'intégration : Le réglage de la machine de sérigraphie}

\courssub{Objectifs de la séance}

À l'issue de cette activité, tu devras être capable de mobiliser \emph{ensemble} tous les savoir-faire de la leçon sur une situation concrète de modélisation :

\begin{itemize}
\item \cle{reconnaître} qu'une application du plan dans lui-même est linéaire (ou ne l'est pas) ;
\item \cle{déterminer} l'écriture analytique et la matrice d'une application linéaire dans la base canonique ;
\item \cle{déterminer} le noyau et l'image d'une application linéaire du plan, et vérifier leur stabilité ;
\item \cle{justifier} qu'une application linéaire du plan est (ou n'est pas) bijective, en utilisant au choix le noyau, l'image ou l'écriture analytique (déterminant) ;
\item \cle{déterminer} l'expression analytique de la bijection réciproque d'un automorphisme du plan ;
\item rédiger une argumentation mathématique complète sous forme de rapport technique.
\end{itemize}

\begin{aretenir}[Compétence visée]
Famille de situations : \emph{Représentations et transformations des configurations planes dans l'environnement}. Compétence : utiliser les applications linéaires du plan pour modéliser, analyser et contrôler une transformation de motifs.
\end{aretenir}

\courssub{Situation}

\textbf{Le réglage de la machine de sérigraphie.}

\medskip

L'atelier \og Pagnes et Couleurs \fg{} de Mokolo, à Yaoundé, vient d'acquérir une machine de sérigraphie programmable pour imprimer des motifs sur des pagnes destinés à la fête de la jeunesse. La machine travaille sur une plaque carrée rapportée à un repère orthonormé $(O\,;\vec{\imath},\vec{\jmath})$ (unité : le décimètre). Chaque point $M(x\,;y)$ du motif numérique est transformé par la machine en un point imprimé $M'(x'\,;y')$ sur le tissu.

Le technicien de maintenance, M. Etoundi, constate que la machine propose trois programmes de réglage, définis par les applications $f_A$, $f_B$ et $f_C$ du plan vectoriel $\mathcal{P}$ dans lui-même :

\begin{center}
\begin{tabularx}{\linewidth}{|l|X|}
\hline
\rowcolor{popblueL}
\textbf{Réglage} & \textbf{Programme de transformation} \\
\hline
Réglage A & $f_A : (x\,;y) \longmapsto (x'\,;y')$ avec $\begin{cases} x' = 2x + y \\ y' = x - y \end{cases}$ \\
\hline
Réglage B & $f_B : (x\,;y) \longmapsto (x'\,;y')$ avec $\begin{cases} x' = x - 2y \\ y' = -2x + 4y \end{cases}$ \\
\hline
Réglage C & $f_C : (x\,;y) \longmapsto (x'\,;y')$ avec $\begin{cases} x' = x + 1 \\ y' = 2y \end{cases}$ \\
\hline
\end{tabularx}
\end{center}

M. Etoundi sait qu'un réglage est \emph{fiable} seulement si l'application correspondante est \cle{linéaire} : cela garantit que la machine respecte les proportions et les combinaisons de vecteurs du motif. De plus :

\begin{itemize}
\item un réglage qui \og écrase \fg{} le motif sur une droite provoque une \textbf{perte d'information} irrémédiable : deux motifs différents peuvent donner la même impression ;
\item un réglage \textbf{réversible} (bijectif) permet, à partir de l'impression obtenue, de retrouver exactement le motif original : c'est indispensable pour le contrôle qualité.
\end{itemize}

Le chef d'atelier te missionne, toi et ton groupe, pour produire un \textbf{rapport technique} indiquant quel réglage utiliser, lequel bannir, et comment programmer la \og marche arrière \fg{} de la machine.

\courssub{Consignes}

Travail en groupes de quatre élèves. Durée : 55 minutes, dont 10 minutes de restitution au tableau. Chaque groupe rédige son rapport technique en répondant, dans l'ordre, aux questions suivantes.

\begin{enumerate}
\item \textbf{Test de conformité.} Pour chacun des trois réglages, vérifier si l'application est linéaire. Pour les réglages linéaires, on vérifiera les deux conditions $f(\vec{u}+\vec{v}) = f(\vec{u}) + f(\vec{v})$ et $f(\lambda \vec{u}) = \lambda f(\vec{u})$ ; pour le réglage défectueux, on exhibera un contre-exemple (par exemple en calculant l'image du vecteur nul). Conclure : quel réglage doit être éliminé d'office ?
\item \textbf{Écriture analytique et matricielle.} Pour chaque réglage linéaire, calculer les images des vecteurs de base $f(\vec{\imath})$ et $f(\vec{\jmath})$, puis écrire la matrice de l'application dans la base $(\vec{\imath},\vec{\jmath})$.
\item \textbf{Noyau.} Déterminer le noyau de chaque réglage linéaire (résoudre le système $f(x\,;y) = (0\,;0)$). Donner une interprétation géométrique : le noyau est-il réduit au vecteur nul, ou est-ce une droite vectorielle ? Que subissent les vecteurs du noyau lors de l'impression ?
\item \textbf{Image.} Déterminer l'image de chaque réglage linéaire (ensemble des vecteurs $(x'\,;y')$ atteignables). Donner une équation caractéristique de cette image lorsqu'elle est une droite.
\item \textbf{Diagnostic.} Identifier :
\begin{enumerate}
\item le réglage qui \og écrase \fg{} le motif sur une droite (perte d'information) — justifier par le noyau \emph{et} par l'image ;
\item le réglage réversible (automorphisme) — justifier la bijectivité par au moins deux méthodes différentes (noyau, image, déterminant de la matrice).
\end{enumerate}
\item \textbf{Marche arrière.} Pour le réglage bijectif, déterminer l'expression analytique de la bijection réciproque $f^{-1}$, c'est-à-dire exprimer $(x\,;y)$ en fonction de $(x'\,;y')$. Vérifier sur un exemple : le point imprimé $M'(5\,;1)$ provient-il bien d'un point $M$ que l'on calculera ?
\item \textbf{Rédaction du rapport.} Rédiger la conclusion technique destinée au chef d'atelier : réglage recommandé, réglage à bannir, et programme de la marche arrière.
\end{enumerate}

\courssub{Démarche et corrigé commenté}

\begin{exoresolu}[Rapport technique — corrigé complet]
On étudie les trois réglages $f_A$, $f_B$, $f_C$ définis dans la situation.

\tcblower

\textbf{1. Test de linéarité.}

\emph{Réglage A.} Soient $\vec{u}(x_1\,;y_1)$, $\vec{v}(x_2\,;y_2)$ deux vecteurs quelconques de $\mathcal{P}$ et $\lambda \in \mathbb{R}$.
\begin{itemize}
\item \textbf{Additivité :} $\vec{u}+\vec{v} = (x_1+x_2\,;\,y_1+y_2)$.
\begin{align*}
f_A(\vec{u}+\vec{v}) &= \bigl(2(x_1+x_2) + (y_1+y_2)\,;\,(x_1+x_2) - (y_1+y_2)\bigr) \\
&= (2x_1+y_1 + 2x_2+y_2\,;\,x_1-y_1 + x_2-y_2) \\
&= (2x_1+y_1\,;\,x_1-y_1) + (2x_2+y_2\,;\,x_2-y_2) \\
&= f_A(\vec{u}) + f_A(\vec{v}).
\end{align*}
\item \textbf{Homogénéité :} $\lambda\vec{u} = (\lambda x_1\,;\,\lambda y_1)$.
\[ f_A(\lambda\vec{u}) = (2\lambda x_1 + \lambda y_1\,;\,\lambda x_1 - \lambda y_1) = \lambda(2x_1+y_1\,;\,x_1-y_1) = \lambda f_A(\vec{u}). \]
\end{itemize}
Les deux propriétés sont vérifiées : $f_A$ est \textbf{linéaire}.

\emph{Réglage B.} Même démarche. Les coordonnées de $f_B$ sont des combinaisons linéaires de $x$ et $y$ sans terme constant. Pour $\vec{u}(x_1,y_1), \vec{v}(x_2,y_2)$ :
\begin{align*}
f_B(\vec{u}+\vec{v}) &= \bigl((x_1+x_2)-2(y_1+y_2)\,;\,-2(x_1+x_2)+4(y_1+y_2)\bigr) \\
&= (x_1-2y_1\,;\,-2x_1+4y_1) + (x_2-2y_2\,;\,-2x_2+4y_2) = f_B(\vec{u})+f_B(\vec{v}), \\
f_B(\lambda\vec{u}) &= (\lambda x_1 - 2\lambda y_1\,;\,-2\lambda x_1 + 4\lambda y_1) = \lambda(x_1-2y_1\,;\,-2x_1+4y_1) = \lambda f_B(\vec{u}).
\end{align*}
$f_B$ est \textbf{linéaire}.

\emph{Réglage C.} Testons la condition nécessaire $f(\vec{0}) = \vec{0}$ :
\[ f_C(0\,;0) = (0+1\,;\,2\times 0) = (1\,;0) \neq (0\,;0). \]
$f_C$ ne vérifie pas la condition nécessaire de linéarité. C'est une application \textbf{affine} (translation suivant $\vec{\imath}$ + homothétie de rapport 2 suivant $\vec{\jmath}$), mais \textbf{pas linéaire}.

\textbf{Conclusion : le réglage C est défectueux et éliminé d'office.}

\medskip

\textbf{2. Écriture analytique et matrices dans la base canonique $(\vec{\imath},\vec{\jmath})$.}

Pour une application linéaire $f(x,y) = (ax+by, cx+dy)$, la matrice est $M = \begin{pmatrix} a & b \\ c & d \end{pmatrix}$ avec $f(\vec{\imath}) = \begin{pmatrix} a \\ c \end{pmatrix}$, $f(\vec{\jmath}) = \begin{pmatrix} b \\ d \end{pmatrix}$.

\begin{itemize}
\item \textbf{Réglage A :} $f_A(x,y) = (2x+y, x-y)$.
$f_A(\vec{\imath}) = (2\,;1)$, $f_A(\vec{\jmath}) = (1\,;-1)$.
\[ M_A = \begin{pmatrix} 2 & 1 \\ 1 & -1 \end{pmatrix}. \]
\item \textbf{Réglage B :} $f_B(x,y) = (x-2y, -2x+4y)$.
$f_B(\vec{\imath}) = (1\,;-2)$, $f_B(\vec{\jmath}) = (-2\,;4)$.
\[ M_B = \begin{pmatrix} 1 & -2 \\ -2 & 4 \end{pmatrix}. \]
\end{itemize}

\medskip

\textbf{3. Noyaux.}

\emph{Noyau de $f_A$.} On résout $f_A(x,y) = (0,0)$ :
\[ \begin{cases} 2x + y = 0 \\ x - y = 0 \end{cases} \iff \begin{cases} y = -2x \\ y = x \end{cases} \iff x = 0,\ y = 0. \]
\[ \ker f_A = \{\vec{0}\}. \]
\textit{Interprétation :} Seul le point origine est envoyé sur l'origine. Aucun vecteur non nul n'est « effacé » à l'impression.

\emph{Noyau de $f_B$.} On résout $f_B(x,y) = (0,0)$ :
\[ \begin{cases} x - 2y = 0 \\ -2x + 4y = 0 \end{cases} \iff x = 2y \quad (\text{la 2\up{ème} équation est } -2 \times \text{la 1\up{ère}}). \]
\[ \ker f_B = \{(2y\,;\,y) \mid y \in \mathbb{R}\} = \operatorname{Vect}\bigl((2\,;1)\bigr). \]
C'est la droite vectorielle d'équation cartésienne $x - 2y = 0$.
\textit{Interprétation :} Tous les points du motif situés sur cette droite sont écrasés sur l'origine $(0,0)$ à l'impression : perte totale de l'information pour ces points.

\medskip

\textbf{4. Images.}

\emph{Image de $f_A$.} $\operatorname{Im} f_A = \operatorname{Vect}\bigl(f_A(\vec{\imath}), f_A(\vec{\jmath})\bigr) = \operatorname{Vect}\bigl((2\,;1), (1\,;-1)\bigr)$.
Le déterminant $\det\begin{pmatrix} 2 & 1 \\ 1 & -1 \end{pmatrix} = -2 - 1 = -3 \neq 0$, donc les deux vecteurs sont non colinéaires et forment une base du plan.
\[ \operatorname{Im} f_A = \mathcal{P}. \]
Tout point du tissu peut être atteint : l'impression couvre toute la surface.

\emph{Image de $f_B$.} $\operatorname{Im} f_B = \operatorname{Vect}\bigl((1\,;-2), (-2\,;4)\bigr)$.
On observe que $(-2\,;4) = -2(1\,;-2)$ : les deux vecteurs sont colinéaires.
\[ \operatorname{Im} f_B = \operatorname{Vect}\bigl((1\,;-2)\bigr). \]
C'est la droite vectorielle engendrée par $(1\,;-2)$. Une équation caractéristique : $2x' + y' = 0$ (car $2(1) + (-2) = 0$).
Vérification : pour tout $(x,y)$, $2x' + y' = 2(x-2y) + (-2x+4y) = 0$.
\textit{Interprétation :} L'impression entière est confinée sur cette droite. Le motif 2D est « aplati » sur un trait : perte d'information massive.

\medskip

\textbf{5. Diagnostic.}

\begin{enumerate}
\item \textbf{Réglage B : Écrasement (perte d'information).}
$\ker f_B$ est une droite (dimension 1) $\neq \{\vec{0}\}$ $\implies$ $f_B$ n'est \textbf{pas injective} (plusieurs antécédents pour le même point, notamment $\vec{0}$).
$\operatorname{Im} f_B$ est une droite (dimension 1) $\neq \mathcal{P}$ $\implies$ $f_B$ n'est \textbf{pas surjective} (des points du tissu ne sont jamais atteints).
Le théorème du rang est vérifié : $\dim(\ker f_B) + \dim(\operatorname{Im} f_B) = 1 + 1 = 2$.
\textbf{Verdict : Réglage B à bannir.}
\item \textbf{Réglage A : Réversibilité (automorphisme).} Trois justifications concordantes :
\begin{itemize}
\item $\ker f_A = \{\vec{0}\}$ $\implies$ $f_A$ injective.
\item $\operatorname{Im} f_A = \mathcal{P}$ $\implies$ $f_A$ surjective.
\item $\det M_A = -3 \neq 0$ $\implies$ $M_A$ inversible $\implies$ $f_A$ bijective.
\end{itemize}
$f_A$ est un \textbf{automorphisme} du plan.
\textbf{Verdict : Réglage A recommandé.}
\end{enumerate}

\medskip

\textbf{6. Marche arrière : calcul de $f_A^{-1}$.}

On résout le système par rapport à $(x,y)$ :
\[ \begin{cases} x' = 2x + y \\ y' = x - y \end{cases} \]
Méthode : addition membre à membre $\implies x' + y' = 3x \implies x = \dfrac{x'+y'}{3}$.
On substitue dans la deuxième équation : $y = x - y' = \dfrac{x'+y'}{3} - y' = \dfrac{x' - 2y'}{3}$.
D'où l'expression analytique de la bijection réciproque :
\[ f_A^{-1}(x'\,;y') = \left(\dfrac{x'+y'}{3}\,;\,\dfrac{x'-2y'}{3}\right). \]
Sa matrice dans la base canonique est $M_A^{-1} = \dfrac{1}{-3}\begin{pmatrix} -1 & -1 \\ -1 & 2 \end{pmatrix} = \begin{pmatrix} 1/3 & 1/3 \\ 1/3 & -2/3 \end{pmatrix}$.

\emph{Vérification sur l'exemple $M'(5\,;1)$.}
\[ x = \frac{5+1}{3} = 2, \qquad y = \frac{5-2}{3} = 1. \]
Le point original est $M(2\,;1)$. Contrôle : $f_A(2\,;1) = (2\times2+1\,;\,2-1) = (5\,;1)$. Le calcul est correct.

\medskip

\textbf{7. Conclusion du rapport technique.}

\begin{center}
\begin{tabular}{|p{4cm}|p{11cm}|}
\hline
\rowcolor{popgreenL} \textbf{Décision} & \textbf{Justification technique} \\
\hline
\textbf{Réglage A : ADOPTÉ} & Automorphisme (bijectif). Impression fidèle, proportions respectées, réversible. \\
\hline
\textbf{Réglage B : BANNI} & Noyau et image de dimension 1. Écrase le motif sur la droite $2x'+y'=0$. Perte d'information irrémédiable. \\
\hline
\textbf{Réglage C : BANNI} & Non linéaire ($f_C(\vec{0}) \neq \vec{0}$). Ne respecte pas la structure vectorielle (translation parasite). \\
\hline
\end{tabular}
\end{center}

\textbf{Programme de la marche arrière (pour réglage A) :}
\[ \boxed{ \begin{cases} x = \dfrac{x'+y'}{3} \\[6pt] y = \dfrac{x'-2y'}{3} \end{cases} } \]
Ce programme permet, à partir des coordonnées $(x',y')$ du point imprimé, de retrouver les coordonnées $(x,y)$ du point du motif original. Il sera implémenté dans le module de contrôle qualité de la machine.
\end{exoresolu}

\courssub{Bilan de l'activité}

Cette activité a permis de mettre en évidence, sur une situation authentique de l'artisanat yaoundéen, les points essentiels de la leçon M19 :

\begin{itemize}
\item \textbf{La linéarité n'est pas automatique.} Le réglage C, pourtant d'apparence simple, échoue au test fondamental $f(\vec{0}) = \vec{0}$ : la présence d'un terme constant (une translation) détruit la linéarité. Au Bac, vérifier $f(\vec{0})=\vec{0}$ est le réflexe rapide pour éliminer une application affine.
\item \textbf{Le noyau mesure la perte d'information (défaut d'injectivité).} Plus le noyau est \og gros \fg{} (dimension 1 au lieu de 0), plus l'application écrase de vecteurs distincts sur le même image ($\vec{0}$). Un noyau réduit à $\{\vec{0}\}$ garantit qu'aucun détail du motif n'est confondu avec un autre.
\item \textbf{L'image mesure la couverture (défaut de surjectivité).} Si l'image est une droite (dimension 1), tout le plan cible n'est pas atteint : des zones du tissu restent blanches quel que soit le motif envoyé.
\item \textbf{Noyau et image sont complémentaires (Théorème du rang).} Pour $f_B$, $\dim(\ker f_B) + \dim(\operatorname{Im} f_B) = 1 + 1 = 2$. Ce que l'application \og perd \fg{} d'un côté (noyau non trivial) se retrouve \og manquant \fg{} de l'autre (image non totale).
\item \textbf{La bijectivité se lit de plusieurs façons équivalentes.} Noyau trivial, image totale, déterminant non nul : trois critères qu'il faut savoir mobiliser selon les données de l'exercice — une exigence classique du Baccalauréat TI (souvent : « justifier par deux méthodes différentes »).
\item \textbf{La réciproque a un sens concret.} Calculer $f_A^{-1}$, c'est programmer la \og marche arrière \fg{} de la machine : les mathématiques deviennent un outil de contrôle qualité industriel.
\end{itemize}

\begin{savaistu}[Pour aller plus loin : le théorème du rang]
En dimension finie, le \emph{théorème du rang} (ou théorème noyau-image) affirme que pour toute application linéaire $f : E \to F$,
\[ \dim(\ker f) + \dim(\operatorname{Im} f) = \dim(E). \]
Tu l'as observé ici sans le nommer : pour $f_B$, $1 + 1 = 2$ ; pour $f_A$, $0 + 2 = 2$. Ce résultat, pierre angulaire de l'algèbre linéaire à l'université, gouverne aussi la compression d'images numériques (JPEG, envoi de photos sur WhatsApp via MTN/Orange) : on projette l'image sur un sous-espace de plus petite dimension (noyau non trivial = informations supprimées) pour réduire la taille du fichier.
\end{savaistu}

\newpage

\coursec{Méthodes et exercices résolus}

\courssub{Savoir-faire 1 : Reconnaître qu'une application est linéaire}

\begin{methode}[Prouver (ou réfuter) la linéarité d'une application du plan]
Pour montrer qu'une application $f : \mathbb{R}^2 \to \mathbb{R}^2$ est \cle{linéaire}, deux stratégies :
\begin{enumerate}
\item \textbf{Vérifier la définition} : pour tous vecteurs $\vec{u}, \vec{v}$ et tout réel $\lambda$, montrer que
\[ f(\vec{u}+\vec{v}) = f(\vec{u}) + f(\vec{v}) \quad \text{et} \quad f(\lambda \vec{u}) = \lambda f(\vec{u}). \]
On peut condenser en une seule vérification : $f(\lambda \vec{u} + \vec{v}) = \lambda f(\vec{u}) + f(\vec{v})$.
\item \textbf{Utiliser l'écriture analytique} : si $f(x,y) = (ax+by,\; cx+dy)$ avec $a,b,c,d$ réels (expressions \emph{du premier degré, sans terme constant}), alors $f$ est linéaire.
\end{enumerate}
Pour \textbf{réfuter} la linéarité, un contre-exemple suffit : tester $f(\vec{0}) = \vec{0}$ (si $f(\vec{0}) \neq \vec{0}$, $f$ n'est pas linéaire), ou exhiber $\vec{u}$ et $\lambda$ tels que $f(\lambda\vec{u}) \neq \lambda f(\vec{u})$.
\end{methode}

\begin{exoresolu}[Linéaire ou pas ?]
On considère les applications de $\mathbb{R}^2$ dans $\mathbb{R}^2$ définies par :
\[ f(x,y) = (2x - 3y,\; x + y) \;;\qquad g(x,y) = (x + 1,\; y) \;;\qquad h(x,y) = (x^2,\; y). \]
Pour chacune, dire si elle est linéaire, en justifiant.
\tcblower
\textbf{Étude de $f$.} Posons $\vec{u} = (x, y)$ et $\vec{u'} = (x', y')$, et soit $\lambda \in \mathbb{R}$.
\begin{align*}
f(\lambda \vec{u} + \vec{u'}) &= f(\lambda x + x',\; \lambda y + y') \\
&= \big(2(\lambda x + x') - 3(\lambda y + y'),\; (\lambda x + x') + (\lambda y + y')\big) \\
&= \lambda (2x - 3y,\; x + y) + (2x' - 3y',\; x' + y') \\
&= \lambda f(\vec{u}) + f(\vec{u'}).
\end{align*}
Donc $f$ est \textbf{linéaire}. On pouvait aussi le voir directement : les deux coordonnées sont des expressions du premier degré en $x$ et $y$, sans terme constant.

\textbf{Étude de $g$.} On a $g(0,0) = (1, 0) \neq (0,0)$. Or toute application linéaire envoie le vecteur nul sur le vecteur nul. Donc $g$ \textbf{n'est pas linéaire} (c'est une application \emph{affine}).

\textbf{Étude de $h$.} Testons avec $\vec{u} = (1, 0)$ et $\lambda = 2$ :
\[ h(2\vec{u}) = h(2, 0) = (4, 0) \quad \text{alors que} \quad 2h(\vec{u}) = 2(1, 0) = (2, 0). \]
Comme $h(2\vec{u}) \neq 2h(\vec{u})$, $h$ \textbf{n'est pas linéaire}.

\textbf{Conclusion.} Seule $f$ est une application linéaire.
\end{exoresolu}

\courssub{Endomorphismes, isomorphismes, automorphismes}

\begin{definition}[Endomorphisme, isomorphisme, automorphisme]
Soient $E$ et $F$ deux espaces vectoriels (ici $E=F=\mathbb{R}^2$).
\begin{itemize}
\item Une application linéaire $f : E \to E$ est appelée un \cle{endomorphisme} de $E$.
\item Une application linéaire $f : E \to F$ qui est bijective est appelée un \cle{isomorphisme} de $E$ dans $F$. Si $E=F$, on parle d'\cle{automorphisme} de $E$.
\item Deux espaces vectoriels sont \cle{isomorphes} s'il existe un isomorphisme de l'un dans l'autre. Tout espace vectoriel de dimension $2$ est isomorphe à $\mathbb{R}^2$.
\end{itemize}
\end{definition}

\begin{propriete}[Caractérisation des isomorphismes en dimension finie]
Soit $f : E \to F$ une application linéaire entre deux espaces vectoriels de même dimension finie $n$. Les propriétés suivantes sont équivalentes :
\begin{enumerate}
\item $f$ est un isomorphisme (bijective) ;
\item $f$ est injective ($\ker f = \{\vec{0}\}$) ;
\item $f$ est surjective ($\operatorname{Im} f = F$) ;
\item L'image par $f$ de toute base de $E$ est une base de $F$.
\end{enumerate}
\end{propriete}

\courssub{Savoir-faire 2 : Déterminer l'image d'un vecteur et l'expression dans une base}

\begin{methode}[Calculer $f(\vec{u})$ à partir de la matrice ou des images des vecteurs de base]
Soit $\mathcal{B} = (\vec{i}, \vec{j})$ une base du plan et $f$ une application linéaire.
\begin{enumerate}
\item \textbf{Si la matrice est connue} : $M = \begin{pmatrix} a & b \\ c & d \end{pmatrix}$ signifie $f(\vec{i}) = a\vec{i} + c\vec{j}$ (première colonne) et $f(\vec{j}) = b\vec{i} + d\vec{j}$ (deuxième colonne). Alors pour $\vec{u} = x\vec{i} + y\vec{j}$ :
\[ f(\vec{u}) = (ax + by)\vec{i} + (cx + dy)\vec{j}, \qquad \text{i.e.} \quad \begin{pmatrix} x' \\ y' \end{pmatrix} = \begin{pmatrix} a & b \\ c & d \end{pmatrix}\begin{pmatrix} x \\ y \end{pmatrix}. \]
\item \textbf{Si seules les images de la base sont connues} : décomposer $\vec{u} = x\vec{i} + y\vec{j}$, puis utiliser la linéarité : $f(\vec{u}) = x f(\vec{i}) + y f(\vec{j})$.
\end{enumerate}
\textbf{Réflexe} : les colonnes de la matrice sont les coordonnées des images des vecteurs de la base, dans cette base.
\end{methode}

\begin{exoresolu}[Image d'un vecteur par une application linéaire]
Le plan est muni d'une base $\mathcal{B} = (\vec{i}, \vec{j})$. Soit $f$ l'application linéaire définie par
\[ f(\vec{i}) = 3\vec{i} - \vec{j} \quad \text{et} \quad f(\vec{j}) = 2\vec{i} + \vec{j}. \]
\begin{enumerate}
\item Écrire la matrice de $f$ dans la base $\mathcal{B}$ et l'expression analytique de $f$.
\item Déterminer l'image du vecteur $\vec{u} = 4\vec{i} - 5\vec{j}$.
\item Un opérateur télécom modélise le déplacement d'un signal par $\vec{v} = -\vec{i} + 2\vec{j}$. Calculer $f(\vec{v})$.
\end{enumerate}
\tcblower
\textbf{1. Matrice et expression analytique.} Les colonnes de la matrice sont les coordonnées de $f(\vec{i})$ et $f(\vec{j})$ dans $\mathcal{B}$ :
\[ M = \begin{pmatrix} 3 & 2 \\ -1 & 1 \end{pmatrix}. \]
Pour $\vec{w} = x\vec{i} + y\vec{j}$, on a donc
\[ f(\vec{w}) = (3x + 2y)\vec{i} + (-x + y)\vec{j}, \quad \text{soit} \quad f(x, y) = (3x + 2y,\; -x + y). \]

\textbf{2. Image de $\vec{u} = 4\vec{i} - 5\vec{j}$.} Par linéarité :
\begin{align*}
f(\vec{u}) &= 4 f(\vec{i}) - 5 f(\vec{j}) \\
&= 4(3\vec{i} - \vec{j}) - 5(2\vec{i} + \vec{j}) \\
&= (12 - 10)\vec{i} + (-4 - 5)\vec{j} \\
&= 2\vec{i} - 9\vec{j}.
\end{align*}
Vérification par la formule analytique : $f(4, -5) = (3(4) + 2(-5),\; -4 + (-5)) = (2, -9)$. On retrouve bien $\vec{u'} = 2\vec{i} - 9\vec{j}$.

\textbf{3. Image de $\vec{v} = -\vec{i} + 2\vec{j}$.}
\[ f(\vec{v}) = -f(\vec{i}) + 2f(\vec{j}) = -(3\vec{i} - \vec{j}) + 2(2\vec{i} + \vec{j}) = \vec{i} + 3\vec{j}. \]

\textbf{Conclusion.} $M = \begin{pmatrix} 3 & 2 \\ -1 & 1 \end{pmatrix}$, $f(\vec{u}) = 2\vec{i} - 9\vec{j}$ et $f(\vec{v}) = \vec{i} + 3\vec{j}$.
\end{exoresolu}

\courssub{Savoir-faire 3 : Déterminer le noyau d'une application linéaire}

\begin{methode}[Déterminer $\ker f$ : équation caractéristique et base]
Le \cle{noyau} de $f$ est $\ker f = \{\vec{u} \in \mathbb{R}^2 \mid f(\vec{u}) = \vec{0}\}$.
\begin{enumerate}
\item Écrire l'équation $f(x, y) = (0, 0)$, c'est-à-dire le système $\begin{cases} ax + by = 0 \\ cx + dy = 0 \end{cases}$ appelé \cle{équation caractéristique} du noyau.
\item Résoudre le système (substitution ou combinaison).
\item Conclure selon le nombre de solutions :
\begin{itemize}
\item seule solution $(0,0)$ : $\ker f = \{\vec{0}\}$ ;
\item une infinité de solutions dépendant d'un paramètre $t$ : $\ker f$ est la droite vectorielle engendrée par un vecteur non nul $\vec{e}$, et $(\vec{e})$ en est une \textbf{base}.
\end{itemize}
\item \textbf{Vérifier} que le ou les vecteurs trouvés ont bien une image nulle.
\end{enumerate}
\end{methode}

\begin{exoresolu}[Noyau d'une application linéaire]
Soit $f$ l'application linéaire de $\mathbb{R}^2$ dans $\mathbb{R}^2$ définie par
\[ f(x, y) = (2x - y,\; -4x + 2y). \]
Déterminer le noyau de $f$ et en donner une base.
\tcblower
\textbf{Équation caractéristique.} $\vec{u} = (x, y) \in \ker f$ si et seulement si $f(\vec{u}) = \vec{0}$, c'est-à-dire :
\[ \begin{cases} 2x - y = 0 \quad (L_1) \\ -4x + 2y = 0 \quad (L_2) \end{cases} \]

\textbf{Résolution.} Remarquons que $(L_2) = -2 \times (L_1)$ : les deux équations sont équivalentes. Le système se réduit donc à $2x - y = 0$, soit $y = 2x$.

En posant $x = t$ ($t \in \mathbb{R}$), les solutions sont les vecteurs $(t, 2t) = t(1, 2)$.

\textbf{Conclusion.} Le noyau est la droite vectorielle
\[ \ker f = \{ t(1, 2) \mid t \in \mathbb{R} \} = \operatorname{Vect}\big(\vec{e}\big) \quad \text{avec} \quad \vec{e} = (1, 2). \]
Le vecteur $\vec{e}$ étant non nul, la famille $(\vec{e})$ est libre et engendre $\ker f$ : c'est une \textbf{base} de $\ker f$, qui est donc de dimension $1$.

\textbf{Vérification.} $f(1, 2) = (2(1) - 2,\; -4(1) + 2(2)) = (0, 0)$ : correct.
\end{exoresolu}

\courssub{Savoir-faire 4 : Déterminer l'image d'une application linéaire}

\begin{methode}[Déterminer $\operatorname{Im} f$ et en donner une base]
L'\cle{image} de $f$ est $\operatorname{Im} f = \{ f(\vec{u}) \mid \vec{u} \in \mathbb{R}^2 \}$.
\begin{enumerate}
\item Écrire $f(x, y) = x\, f(\vec{i}) + y\, f(\vec{j})$ : cela montre que
\[ \operatorname{Im} f = \operatorname{Vect}\big(f(\vec{i}), f(\vec{j})\big). \]
L'image est donc engendrée par les \textbf{colonnes de la matrice}.
\item Extraire une famille libre maximale :
\begin{itemize}
\item si $f(\vec{i})$ et $f(\vec{j})$ sont non colinéaires : $\operatorname{Im} f = \mathbb{R}^2$ tout entier ;
\item s'ils sont colinéaires (avec l'un non nul) : $\operatorname{Im} f$ est une droite vectorielle, dont une base est un des deux vecteurs non nul.
\end{itemize}
\item On peut aussi chercher une \textbf{équation} de l'image : $(x', y') \in \operatorname{Im} f$ s'il existe $(x, y)$ avec $x' = ax + by$ et $y' = cx + dy$.
\end{enumerate}
\textbf{Réflexe de contrôle} : $\dim \ker f + \dim \operatorname{Im} f = 2$ (dimension du plan).
\end{methode}

\begin{exoresolu}[Image d'une application linéaire]
On reprend l'application $f(x, y) = (2x - y,\; -4x + 2y)$ de l'exercice précédent.
\begin{enumerate}
\item Déterminer $\operatorname{Im} f$ et en donner une base.
\item Vérifier la relation $\dim \ker f + \dim \operatorname{Im} f = 2$.
\item Le point $A(3, -6)$ appartient-il à $\operatorname{Im} f$ ? Et le point $B(1, 1)$ ?
\end{enumerate}
\tcblower
\textbf{1. Détermination de l'image.} Pour tout $(x, y) \in \mathbb{R}^2$ :
\[ f(x, y) = x(2, -4) + y(-1, 2). \]
Donc $\operatorname{Im} f = \operatorname{Vect}\big((2, -4), (-1, 2)\big)$.

Or $(2, -4) = -2(-1, 2)$ : les deux vecteurs sont \textbf{colinéaires}. L'image est donc la droite vectorielle engendrée par $\vec{w} = (-1, 2)$ (ou par $(2,-4)$) :
\[ \operatorname{Im} f = \operatorname{Vect}\big(\vec{w}\big) = \{ k(-1, 2) \mid k \in \mathbb{R} \}, \]
et $(\vec{w})$ est une base de $\operatorname{Im} f$, qui est de dimension $1$.

Autre présentation (équation de l'image) : $(x', y') \in \operatorname{Im} f$ s'il existe $(x,y)$ tel que $x' = 2x - y$ et $y' = -4x + 2y = -2(2x - y) = -2x'$. Donc
\[ \operatorname{Im} f = \{ (x', y') \in \mathbb{R}^2 \mid y' = -2x' \}. \]

\textbf{2. Vérification.} $\dim \ker f + \dim \operatorname{Im} f = 1 + 1 = 2$ : conforme.

\textbf{3. Appartenance.} $A(3, -6)$ : on a $-6 = -2 \times 3$, donc $A \in \operatorname{Im} f$ (en effet $A = -3\vec{w}$, et par exemple $f(3, 3) = (3, -6)$). Pour $B(1, 1)$ : $1 \neq -2 \times 1$, donc $B \notin \operatorname{Im} f$ : l'équation $f(x,y) = (1,1)$ n'a pas de solution.

\textbf{Conclusion.} $\operatorname{Im} f$ est la droite vectorielle d'équation $y' = -2x'$, de base $\big((-1, 2)\big)$.
\end{exoresolu}

\courssub{Savoir-faire 5 : Stabilité du noyau et de l'image}

\begin{methode}[Exploiter la stabilité de $\ker f$ et $\operatorname{Im} f$]
Le noyau et l'image d'une application linéaire sont des \cle{sous-espaces vectoriels} : ils sont \textbf{stables} pour l'addition et la multiplication par un réel. Concrètement :
\begin{enumerate}
\item Si $\vec{u}, \vec{v} \in \ker f$, alors $\vec{u} + \vec{v} \in \ker f$ et $\lambda \vec{u} \in \ker f$ pour tout $\lambda \in \mathbb{R}$.
\item Si $\vec{a}, \vec{b} \in \operatorname{Im} f$, alors $\vec{a} + \vec{b} \in \operatorname{Im} f$ et $\lambda \vec{a} \in \operatorname{Im} f$.
\end{enumerate}
Pour le \textbf{prouver}, on utilise uniquement la linéarité de $f$ : par exemple, si $\vec{u}, \vec{v} \in \ker f$, alors $f(\vec{u} + \vec{v}) = f(\vec{u}) + f(\vec{v}) = \vec{0} + \vec{0} = \vec{0}$.
\end{methode}

\begin{exoresolu}[Stabilité du noyau]
Soit $f$ une application linéaire du plan. On sait que les vecteurs $\vec{u} = (1, 2)$ et $\vec{v} = (-3, -6)$ appartiennent à $\ker f$.
\begin{enumerate}
\item Démontrer que $\vec{w} = 2\vec{u} - \vec{v}$ appartient à $\ker f$.
\item Démontrer de manière générale que si $\vec{u}, \vec{v} \in \ker f$ et $\lambda \in \mathbb{R}$, alors $\vec{u} + \vec{v} \in \ker f$ et $\lambda \vec{u} \in \ker f$.
\end{enumerate}
\tcblower
\textbf{1.} Par linéarité de $f$ :
\[ f(\vec{w}) = f(2\vec{u} - \vec{v}) = 2 f(\vec{u}) - f(\vec{v}) = 2\vec{0} - \vec{0} = \vec{0}, \]
car $\vec{u}, \vec{v} \in \ker f$ signifie $f(\vec{u}) = f(\vec{v}) = \vec{0}$. Donc $\vec{w} \in \ker f$.

\textbf{2. Démonstration générale.} Soient $\vec{u}, \vec{v} \in \ker f$ et $\lambda \in \mathbb{R}$.
\begin{itemize}
\item \emph{Stabilité par addition} : $f(\vec{u} + \vec{v}) = f(\vec{u}) + f(\vec{v}) = \vec{0} + \vec{0} = \vec{0}$, donc $\vec{u} + \vec{v} \in \ker f$.
\item \emph{Stabilité par multiplication par un réel} : $f(\lambda \vec{u}) = \lambda f(\vec{u}) = \lambda \vec{0} = \vec{0}$, donc $\lambda \vec{u} \in \ker f$.
\end{itemize}
De plus $\vec{0} \in \ker f$ (car $f(\vec{0}) = \vec{0}$). Le noyau, non vide et stable par addition et multiplication par un réel, est donc un sous-espace vectoriel du plan. La démonstration est identique pour l'image : si $\vec{a} = f(\vec{u})$ et $\vec{b} = f(\vec{v})$, alors $\vec{a} + \vec{b} = f(\vec{u} + \vec{v}) \in \operatorname{Im} f$ et $\lambda \vec{a} = f(\lambda \vec{u}) \in \operatorname{Im} f$.

\textbf{Conclusion.} Noyau et image sont stables par addition et par multiplication par un réel : ce sont des sous-espaces vectoriels.
\end{exoresolu}

\courssub{Savoir-faire 6 : Étudier la bijectivité d'une application linéaire}

\begin{methode}[Caractériser la bijectivité : trois points de vue équivalents]
Pour une application linéaire $f$ du plan de matrice $M = \begin{pmatrix} a & b \\ c & d \end{pmatrix}$, les affirmations suivantes sont \textbf{équivalentes} :
\begin{enumerate}
\item $f$ est \textbf{bijective} (c'est un \cle{automorphisme}) ;
\item $\ker f = \{\vec{0}\}$ (le système $f(x,y) = (0,0)$ admet l'unique solution $(0,0)$) ;
\item $\operatorname{Im} f = \mathbb{R}^2$ (les vecteurs $f(\vec{i})$ et $f(\vec{j})$ sont non colinéaires) ;
\item le déterminant est non nul : $ad - bc \neq 0$.
\end{enumerate}
\textbf{Démarche au Bac} : choisir \emph{un seul} critère (le plus rapide, souvent le déterminant), le vérifier proprement, puis conclure en citant le théorème de caractérisation. Si $f$ est bijective, $f$ est un automorphisme et l'image de toute base est une base.
\end{methode}

\begin{exoresolu}[Bijectivité : deux applications à comparer]
Dans une base $\mathcal{B} = (\vec{i}, \vec{j})$, on considère les applications linéaires $f$ et $g$ de matrices
\[ M_f = \begin{pmatrix} 3 & 2 \\ -1 & 1 \end{pmatrix} \qquad \text{et} \qquad M_g = \begin{pmatrix} 2 & -1 \\ -4 & 2 \end{pmatrix}. \]
\begin{enumerate}
\item L'application $f$ est-elle un automorphisme du plan ? Justifier.
\item L'application $g$ est-elle bijective ? Justifier par le noyau, puis par l'image.
\end{enumerate}
\tcblower
\textbf{1. Étude de $f$.} Calculons le déterminant :
\[ \det M_f = 3 \times 1 - 2 \times (-1) = 3 + 2 = 5 \neq 0. \]
Le déterminant étant non nul, $f$ est bijective. Comme c'est un endomorphisme bijectif du plan, $f$ est un \textbf{automorphisme} du plan.

Vérification par le noyau : $f(x,y) = (0,0)$ donne $\begin{cases} 3x + 2y = 0 \\ -x + y = 0 \end{cases}$. La deuxième équation donne $y = x$, que l'on reporte : $3x + 2x = 5x = 0$, donc $x = y = 0$. Ainsi $\ker f = \{\vec{0}\}$, ce qui confirme la bijectivité.

\textbf{2. Étude de $g$.} On a $g(x, y) = (2x - y,\; -4x + 2y)$ : c'est l'application des exercices précédents.

\emph{Par le noyau} : nous avons montré que $\ker g = \operatorname{Vect}\big((1, 2)\big) \neq \{\vec{0}\}$. Le noyau n'étant pas réduit au vecteur nul, $g$ \textbf{n'est pas injective}, donc pas bijective.

\emph{Par l'image} : $\operatorname{Im} g = \operatorname{Vect}\big((-1, 2)\big)$ est une droite vectorielle, strictement incluse dans $\mathbb{R}^2$. Donc $g$ \textbf{n'est pas surjective}, donc pas bijective.

\emph{Contrôle par le déterminant} : $\det M_g = 2 \times 2 - (-1)(-4) = 4 - 4 = 0$ : cohérent.

\textbf{Conclusion.} $f$ est un automorphisme du plan ($\det M_f = 5 \neq 0$) ; $g$ n'est ni injective ni surjective ($\det M_g = 0$).
\end{exoresolu}

\begin{exoresolu}[Synthèse type Bac : étude complète d'une application linéaire]
Le plan est muni d'une base $\mathcal{B} = (\vec{i}, \vec{j})$. Soit $m$ un réel et $f_m$ l'application linéaire définie par
\[ f_m(x, y) = \big((m-1)x + 2y,\; 2x + (m+1)y\big). \]
\begin{enumerate}
\item Écrire la matrice de $f_m$ dans $\mathcal{B}$.
\item Déterminer les valeurs de $m$ pour lesquelles $f_m$ est un automorphisme.
\item Pour $m = -3$, déterminer le noyau et l'image de $f_{-3}$, avec une base de chacun.
\end{enumerate}
\tcblower
\textbf{1. Matrice.} Les colonnes sont les images de $\vec{i}$ et $\vec{j}$ :
\[ M_m = \begin{pmatrix} m-1 & 2 \\ 2 & m+1 \end{pmatrix}. \]

\textbf{2. Condition d'automorphisme.} $f_m$ est bijective si et seulement si $\det M_m \neq 0$ :
\begin{align*}
\det M_m &= (m-1)(m+1) - 2 \times 2 = m^2 - 1 - 4 = m^2 - 5.
\end{align*}
Or $m^2 - 5 = 0 \iff m = \sqrt{5}$ ou $m = -\sqrt{5}$. Donc
\[ f_m \text{ est un automorphisme} \iff m \notin \{-\sqrt{5},\; \sqrt{5}\}. \]

\textbf{3. Cas $m = -3$.} Alors $f_{-3}(x, y) = (-4x + 2y,\; 2x - 2y)$ et $\det M_{-3} = 9 - 5 = 4 \neq 0$ : en effet, $-3 \notin \{-\sqrt{5}, \sqrt{5}\}$, donc $f_{-3}$ est \textbf{bijective}. Par conséquent :
\[ \ker f_{-3} = \{\vec{0}\} \quad \text{et} \quad \operatorname{Im} f_{-3} = \mathbb{R}^2. \]
Le noyau, réduit au vecteur nul, a pour base la famille vide (dimension $0$) ; une base de l'image est la base canonique $(\vec{i}, \vec{j})$, ou toute base du plan, par exemple $\big(f_{-3}(\vec{i}), f_{-3}(\vec{j})\big) = \big((-4, 2), (2, -2)\big)$.

Vérifions le noyau directement : $\begin{cases} -4x + 2y = 0 \\ 2x - 2y = 0 \end{cases}$ donne $y = x$ (deuxième ligne) puis $-4x + 2x = -2x = 0$, d'où $x = y = 0$. Confirmé.

\textbf{Conclusion.} $f_m$ est un automorphisme pour tout $m \neq \pm\sqrt{5}$ ; en particulier $f_{-3}$ est bijective, de noyau $\{\vec{0}\}$ et d'image $\mathbb{R}^2$.
\end{exoresolu}

\begin{attention}[Les 5 erreurs qui coûtent des points au Bac]
\begin{enumerate}
\item \textbf{Oublier de vérifier $f(\vec{0}) = \vec{0}$ avant de conclure à la linéarité.} Une application comme $g(x,y) = (x+1, y)$ « ressemble » à une application linéaire mais ne l'est pas. Inversement, ne pas conclure « linéaire » à partir d'un seul exemple numérique : il faut une démonstration générale ou la forme analytique $(ax+by, cx+dy)$.
\item \textbf{Transposer la matrice.} Les \emph{colonnes} de la matrice sont les images des vecteurs de la base : $f(\vec{i})$ en première colonne, $f(\vec{j})$ en deuxième. Écrire $f(\vec{i})$ en ligne fait échouer tous les calculs d'images ensuite.
\item \textbf{Confondre noyau et image.} Le noyau se résout avec $f(x, y) = (0, 0)$ (on cherche les \emph{antécédents} de $\vec{0}$) ; l'image s'obtient avec les vecteurs $f(\vec{i}), f(\vec{j})$ (on cherche les \emph{valeurs} atteintes). Mélanger les deux systèmes est l'erreur la plus fréquente.
\item \textbf{Conclure « bijective » sans justification.} Il faut exhiber un critère vérifié : $\ker f = \{\vec{0}\}$, ou $\operatorname{Im} f = \mathbb{R}^2$, ou $ad - bc \neq 0$, puis citer le théorème d'équivalence. Écrire « $f$ est bijective car elle est linéaire » est faux : une application linéaire n'est pas automatiquement bijective.
\item \textbf{Erreurs de calcul sur le déterminant et les systèmes.} Le déterminant est $ad - bc$ (attention aux signes avec des coefficients négatifs : $(-1)(-4) = +4$). Dans la résolution des systèmes, reporter la solution dans \emph{les deux} équations pour vérifier, et ne pas oublier le paramètre $t$ quand il y a une infinité de solutions : écrire $\ker f = \{(t, 2t), t \in \mathbb{R}\}$, pas seulement « $y = 2x$ » sans conclusion ensembliste.
\end{enumerate}
\end{attention}

\newpage

\coursec{Exercices}

\courssub{Je vérifie mes connaissances}

\begin{exercice}[QCM : une seule bonne réponse]
Pour chaque question, entoure la seule réponse exacte.

\textbf{1.} Soit $f$ une application linéaire du plan vectoriel dans lui-même. Alors $f(\vec{0})$ est égal à :
\begin{itemize}
\item[a)] n'importe quel vecteur du plan ;
\item[b)] $\vec{0}$ ;
\item[c)] $\vec{i}+\vec{j}$ ;
\item[d)] on ne peut pas savoir.
\end{itemize}
\end{exercice}

\begin{exercice}[QCM : une seule bonne réponse]
\textbf{2.} Soit $f(x\,;y)=(3x-y\,;\,x+2y)$. L'image du vecteur $\vec{u}(1\,;-1)$ par $f$ est :
\begin{itemize}
\item[a)] $(4\,;-1)$ ; \qquad b) $(2\,;3)$ ; \qquad c) $(4\,;1)$ ; \qquad d) $(-4\,;-1)$.
\end{itemize}
\end{exercice}

\begin{exercice}[QCM : une seule bonne réponse]
\textbf{3.} Le noyau de l'application linéaire $f(x\,;y)=(x+y\,;\,2x+2y)$ est :
\begin{itemize}
\item[a)] réduit à $\{\vec{0}\}$ ;
\item[b)] la droite d'équation $x+y=0$ ;
\item[c)] le plan tout entier ;
\item[d)] la droite d'équation $y=2x$.
\end{itemize}
\end{exercice}

\begin{exercice}[QCM : une seule bonne réponse]
\textbf{4.} Une application linéaire $f$ du plan dans lui-même est un automorphisme si et seulement si :
\begin{itemize}
\item[a)] $\operatorname{Ker}(f)$ est une droite ;
\item[b)] $\operatorname{Im}(f)$ est une droite ;
\item[c)] $\operatorname{Ker}(f)=\{\vec{0}\}$ ;
\item[d)] $f(\vec{0})=\vec{0}$.
\end{itemize}
\end{exercice}

\begin{exercice}[Vrai ou faux ? Justifie ta réponse]
Réponds par VRAI ou FAUX en justifiant chaque réponse par une démonstration courte ou un contre-exemple.
\begin{itemize}
\item[a)] Toute application linéaire $f$ du plan vectoriel vérifie $f(\vec{0})=\vec{0}$.
\item[b)] Si le noyau d'une application linéaire $f$ du plan est une droite vectorielle, alors $f$ est surjective.
\item[c)] Si $f(x\,;y)=(ax+by\,;\,cx+dy)$ vérifie $ad-bc=0$, alors $f$ est l'application nulle.
\item[d)] La composée de deux automorphismes du plan est un automorphisme du plan.
\end{itemize}
\end{exercice}

\begin{exercice}[Reconnaître une application linéaire]
Parmi les applications suivantes, définies sur le plan vectoriel muni d'une base $(\vec{i},\vec{j})$, indique celles qui sont linéaires en justifiant ta réponse :
\begin{align*}
f_1(x\,;y)&=(2x-3y\,;\,x+y) ; & f_2(x\,;y)&=(x^2\,;\,y) ;\\
f_3(x\,;y)&=(x+1\,;\,y) ; & f_4(x\,;y)&=(x-y\,;\,0) ;\\
f_5(x\,;y)&=(y\,;\,x) ; & f_6(x\,;y)&=(3x\,;\,xy).
\end{align*}
\end{exercice}

\begin{exercice}[Calculs directs]
Soit $f$ l'application linéaire du plan définie par $f(x\,;y)=(x-2y\,;\,3x+y)$.
\begin{enumerate}
\item Calcule $f(\vec{i})$, $f(\vec{j})$ et $f(2\vec{i}-\vec{j})$.
\item Calcule $f(4\,;-2)$. Que remarques-tu ?
\item Détermine l'antécédent éventuel du vecteur $\vec{v}(5\,;1)$ par $f$.
\end{enumerate}
\end{exercice}

\courssub{Je m'entraîne}

\begin{exercice}[Écriture analytique et image d'un vecteur]
Dans la base $(\vec{i},\vec{j})$ du plan vectoriel, $f$ est l'application linéaire telle que :
\[ f(\vec{i})=\vec{i}-2\vec{j} \qquad \text{et} \qquad f(\vec{j})=3\vec{i}+\vec{j}. \]
\begin{enumerate}
\item Détermine l'écriture analytique de $f$, c'est-à-dire l'expression de $f(x\,;y)$ pour tout vecteur $\vec{u}=x\vec{i}+y\vec{j}$.
\item Détermine l'image du vecteur $\vec{u}=2\vec{i}-5\vec{j}$ par $f$.
\item L'application $f$ est-elle bijective ? Justifie ta réponse.
\end{enumerate}
\end{exercice}

\begin{exercice}[Noyau et image]
Soit $f$ l'application linéaire du plan définie par $f(x\,;y)=(2x-4y\,;\,-x+2y)$.
\begin{enumerate}
\item Détermine $\operatorname{Ker}(f)$ et donne-en une base.
\item Détermine $\operatorname{Im}(f)$ et donne-en une base.
\item Compare $\operatorname{Ker}(f)$ et $\operatorname{Im}(f)$. Interprète géométriquement le résultat obtenu.
\item L'application $f$ est-elle bijective ? Justifie.
\end{enumerate}
\end{exercice}

\begin{exercice}[Étude de la bijectivité de deux façons]
Soit $f$ l'application du plan vectoriel dans lui-même définie par :
\[ f(x\,;y)=(x+3y\,;\,2x-y). \]
\begin{enumerate}
\item Montre que $f$ est une application linéaire.
\item Détermine $\operatorname{Ker}(f)$ et $\operatorname{Im}(f)$.
\item L'application $f$ est-elle bijective ? Justifie ta réponse de deux façons différentes : d'une part en utilisant le noyau, d'autre part en utilisant l'écriture analytique (calcul de $ad-bc$).
\end{enumerate}
\end{exercice}

\begin{exercice}[Projection vectorielle]
Le plan est muni d'un repère $(O\,;\vec{i},\vec{j})$. On considère la projection vectorielle $p$ sur la droite $D$ d'équation $y=x$, parallèlement à l'axe des ordonnées.
\begin{enumerate}
\item Fais une figure, puis détermine géométriquement $p(\vec{i})$ et $p(\vec{j})$.
\item Déduis-en l'écriture analytique de $p$.
\item Détermine $\operatorname{Ker}(p)$ et $\operatorname{Im}(p)$. Que représentent-ils géométriquement ?
\item Vérifie par le calcul que $p\circ p=p$. Ce résultat était-il prévisible ? Explique.
\item $p$ est-elle un automorphisme du plan ? Justifie.
\end{enumerate}
\end{exercice}

\courssub{Je me prépare au Bac}

\begin{exercice}[Type Bac : la société de transport « TransLittoral »]
La société de transport « TransLittoral », basée à Douala, modélise le déplacement de ses camions dans le plan par des vecteurs. Un ingénieur de la société associe à chaque vecteur $\vec{u}(x\,;y)$ représentant un trajet, un vecteur corrigé $f(\vec{u})$ tenant compte des déviations, défini dans la base $(\vec{i},\vec{j})$ par :
\[ f(x\,;y)=(x+3y\,;\,2x-y). \]
\begin{enumerate}
\item Montre que $f$ est une application linéaire du plan vectoriel dans lui-même.
\item Un camion effectue le trajet représenté par le vecteur $\vec{u}(7\,;-1)$. Détermine le vecteur corrigé $f(\vec{u})$.
\item Détermine le noyau $\operatorname{Ker}(f)$ de $f$ et donne-en une base.
\item Détermine l'image $\operatorname{Im}(f)$ de $f$.
\item L'application $f$ est-elle un automorphisme du plan ? Justifie ta réponse de deux façons différentes :
\begin{enumerate}
\item en utilisant le noyau ;
\item en utilisant l'écriture analytique de $f$.
\end{enumerate}
\item Détermine l'antécédent par $f$ du vecteur $\vec{v}(10\,;6)$. Interprète le résultat pour l'ingénieur : le trajet corrigé $\vec{v}$ correspond-il à un trajet réel unique ?
\end{enumerate}
\end{exercice}

\begin{exercice}[Type Bac : étude d'une famille d'applications à paramètre]
Une entreprise de télécommunications utilise, pour coder les positions de ses antennes relais dans la région du Centre, une famille d'applications du plan dépendant d'un réel $m$, définies dans la base $(\vec{i},\vec{j})$ par :
\[ f_m(x\,;y)=(mx+y\,;\,x+my). \]
\begin{enumerate}
\item Montre que pour tout réel $m$, $f_m$ est une application linéaire du plan vectoriel dans lui-même.
\item On suppose dans cette question que $m=2$.
\begin{enumerate}
\item Écris l'expression de $f_2(x\,;y)$.
\item Montre que $f_2$ est un automorphisme du plan.
\item Détermine l'antécédent du vecteur $\vec{u}(4\,;5)$ par $f_2$.
\end{enumerate}
\item On suppose dans cette question que $m=1$.
\begin{enumerate}
\item Détermine $\operatorname{Ker}(f_1)$ et donne-en une base.
\item Détermine $\operatorname{Im}(f_1)$ et donne-en une base.
\item $f_1$ est-elle bijective ? Justifie.
\end{enumerate}
\item Détermine, en fonction du réel $m$, les valeurs pour lesquelles $f_m$ est un automorphisme du plan.
\item Pour chacune des valeurs de $m$ pour lesquelles $f_m$ n'est pas un automorphisme, précise $\operatorname{Ker}(f_m)$ et $\operatorname{Im}(f_m)$ par une équation et une base.
\end{enumerate}
\end{exercice}

\begin{exercice}[Type Bac : symétrie, projection et composition]
Le plan est muni d'une base $(\vec{i},\vec{j})$. Un technicien de la SONAREL, chargé d'implanter des lampadaires le long d'une avenue rectiligne de Yaoundé, modélise l'avenue par la droite $D$ d'équation $x-y=0$ et le trottoir par la droite $\Delta$ d'équation $x+y=0$. On considère les deux applications du plan définies par :
\[ s(x\,;y)=(y\,;\,x) \qquad \text{et} \qquad p(x\,;y)=\left(\dfrac{x+y}{2}\,;\,\dfrac{x+y}{2}\right). \]
\begin{enumerate}
\item Montre que $s$ et $p$ sont des applications linéaires du plan vectoriel.
\item Étude de $s$ :
\begin{enumerate}
\item Détermine $\operatorname{Ker}(s)$ et $\operatorname{Im}(s)$.
\item Montre que $s$ est un automorphisme du plan.
\item Calcule $s\circ s(x\,;y)$. Que constates-tu ? (On dit que $s$ est une symétrie vectorielle.)
\end{enumerate}
\item Étude de $p$ :
\begin{enumerate}
\item Détermine $\operatorname{Ker}(p)$ et donne-en une base. Quelle droite du plan reconnais-tu ?
\item Détermine $\operatorname{Im}(p)$ et donne-en une base. Quelle droite du plan reconnais-tu ?
\item Calcule $p\circ p(x\,;y)$. Que constates-tu ? (On dit que $p$ est une projection vectorielle.)
\item $p$ est-elle bijective ? Justifie de deux façons différentes.
\end{enumerate}
\item On pose $g=s\circ p$.
\begin{enumerate}
\item Détermine l'écriture analytique de $g$.
\item $g$ est-elle un automorphisme du plan ? Justifie.
\end{enumerate}
\item Interprète géométriquement, pour le technicien, le rôle des applications $s$ et $p$ par rapport aux droites $D$ et $\Delta$.
\end{enumerate}
\end{exercice}

\newpage

\coursec{Corrigés des exercices}

\begin{corrige}[Exercice 1 — QCM : image du vecteur nul]
\textbf{Réponse : b)} $f(\vec{0})=\vec{0}$.

\textbf{Justification :} pour toute application linéaire $f$ et tout vecteur $\vec{u}$ :
\[ f(\vec{0}) = f(0\cdot\vec{u}) = 0\cdot f(\vec{u}) = \vec{0}. \]
C'est une conséquence immédiate de la conservation de la multiplication par un scalaire. C'est d'ailleurs un test rapide : si $f(\vec{0})\neq\vec{0}$, l'application n'est \emph{pas} linéaire.
\end{corrige}

\begin{corrige}[Exercice 2 — QCM : image d'un vecteur]
\textbf{Réponse : a)} $(4\,;-1)$.

\textbf{Calcul :} avec $f(x\,;y)=(3x-y\,;\,x+2y)$ et $\vec{u}(1\,;-1)$ :
\[ f(1\,;-1) = \bigl(3\times 1 - (-1)\,;\,1 + 2\times(-1)\bigr) = (3+1\,;\,1-2) = (4\,;-1). \]
\end{corrige}

\begin{corrige}[Exercice 3 — QCM : noyau]
\textbf{Réponse : b)} la droite d'équation $x+y=0$.

\textbf{Justification :} $\operatorname{Ker}(f)$ est l'ensemble des $(x,y)$ tels que
\[ \begin{cases} x+y=0 \\ 2x+2y=0 \end{cases} \]
La deuxième équation est le double de la première : le système se réduit à $x+y=0$. Le noyau est donc la droite vectorielle d'équation $x+y=0$, c'est-à-dire $\operatorname{Vect}\bigl((1\,;-1)\bigr)$. Il n'est ni réduit à $\{\vec{0}\}$ ni égal au plan entier.
\end{corrige}

\begin{corrige}[Exercice 4 — QCM : automorphisme]
\textbf{Réponse : c)} $\operatorname{Ker}(f)=\{\vec{0}\}$.

\textbf{Justification :} un automorphisme est un endomorphisme \textbf{bijectif}. Pour une application linéaire du plan dans lui-même, on a les équivalences :
\[ f \text{ bijective} \iff f \text{ injective} \iff \operatorname{Ker}(f)=\{\vec{0}\} \iff \operatorname{Im}(f)=\text{plan entier} \iff ad-bc\neq 0. \]
La réponse d) est vraie pour \emph{toute} application linéaire (même non bijective). Les réponses a) et b) décrivent des applications de rang 1, ni injectives ni surjectives.
\end{corrige}

\begin{corrige}[Exercice 5 — Vrai ou faux]
\begin{itemize}
\item[a)] \textbf{VRAI.} $f(\vec{0})=f(0\cdot\vec{0})=0\cdot f(\vec{0})=\vec{0}$ par linéarité.
\item[b)] \textbf{FAUX.} Si $\operatorname{Ker}(f)$ est une droite, alors $\dim\operatorname{Ker}(f)=1$, donc $\dim\operatorname{Im}(f)=2-1=1$ (théorème du rang) : l'image est une droite, pas le plan, donc $f$ n'est pas surjective. Contre-exemple : $f(x,y)=(x,0)$, de noyau l'axe des ordonnées et d'image l'axe des abscisses.
\item[c)] \textbf{FAUX.} $ad-bc=0$ signifie seulement que $f$ n'est \emph{pas bijective}. Contre-exemple : $f(x,y)=(x,0)$, avec $a=1$, $b=0$, $c=0$, $d=0$, donc $ad-bc=0$, et pourtant $f(1,0)=(1,0)\neq\vec{0}$ : $f$ n'est pas l'application nulle.
\item[d)] \textbf{VRAI.} Soit $f$ et $g$ deux automorphismes. La composée de deux applications linéaires est linéaire, et la composée de deux bijections est une bijection. Donc $f\circ g$ est un endomorphisme bijectif du plan : c'est un automorphisme.
\end{itemize}
\end{corrige}

\begin{corrige}[Exercice 6 — Reconnaître une application linéaire]
Critère pratique : $f$ est linéaire si et seulement si $f(0,0)=(0,0)$ et $f(x,y)$ s'écrit $(ax+by\,;\,cx+dy)$ (coordonnées polynomiales de degré 1, sans terme constant ni produit $xy$ ni carré).

\begin{itemize}
\item $f_1(x,y)=(2x-3y\,;\,x+y)$ : \textbf{linéaire} (forme $ax+by\,;\,cx+dy$ avec $a=2$, $b=-3$, $c=1$, $d=1$).
\item $f_2(x,y)=(x^2\,;\,y)$ : \textbf{non linéaire}. Contre-exemple : $f_2(2,0)=(4,0)$ alors que $2f_2(1,0)=2(1,0)=(2,0)$. Le terme $x^2$ brise l'homogénéité.
\item $f_3(x,y)=(x+1\,;\,y)$ : \textbf{non linéaire} car $f_3(0,0)=(1,0)\neq(0,0)$ (terme constant).
\item $f_4(x,y)=(x-y\,;\,0)$ : \textbf{linéaire} ($a=1$, $b=-1$, $c=d=0$).
\item $f_5(x,y)=(y\,;\,x)$ : \textbf{linéaire} ($a=0$, $b=1$, $c=1$, $d=0$) ; c'est la symétrie par rapport à la droite $y=x$.
\item $f_6(x,y)=(3x\,;\,xy)$ : \textbf{non linéaire}. Contre-exemple : $f_6(1,1)=(3,1)$ alors que $f_6(1,0)+f_6(0,1)=(3,0)+(0,0)=(3,0)$. Le produit $xy$ brise l'additivité.
\end{itemize}
\end{corrige}

\begin{attention}[Point de vigilance — Exercice 6]
La condition $f(\vec{0})=\vec{0}$ est \emph{nécessaire} mais pas \emph{suffisante} : $f_2$ et $f_6$ la vérifient sans être linéaires. Il faut vérifier les deux axiomes (additivité et homogénéité), ou exhiber un contre-exemple précis.
\end{attention}

\begin{corrige}[Exercice 7 — Calculs directs]
\begin{enumerate}
\item $f(\vec{i})=f(1,0)=(1-0\,;\,3+0)=(1\,;\,3)=\vec{i}+3\vec{j}$.\\
$f(\vec{j})=f(0,1)=(0-2\,;\,0+1)=(-2\,;\,1)=-2\vec{i}+\vec{j}$.\\
$f(2\vec{i}-\vec{j})=f(2,-1)=\bigl(2-2\times(-1)\,;\,3\times 2+(-1)\bigr)=(4\,;\,5)$.\\
Contrôle par linéarité : $2f(\vec{i})-f(\vec{j})=2(1,3)-(-2,1)=(2,6)+(2,-1)=(4,5)$. Cohérent.
\item $f(4,-2)=\bigl(4-2\times(-2)\,;\,12-2\bigr)=(8\,;\,10)$.\\
\textbf{Remarque :} $(4,-2)=2(2,-1)$, donc par linéarité $f(4,-2)=2f(2,-1)=2(4,5)=(8,10)$ : on retrouve le résultat sans refaire tout le calcul.
\item On résout $f(x,y)=(5,1)$ :
\[ \begin{cases} x-2y=5 \\ 3x+y=1 \end{cases} \iff \begin{cases} x=5+2y \\ 3(5+2y)+y=1 \end{cases} \iff \begin{cases} x=5+2y \\ 7y=-14 \end{cases} \iff \begin{cases} x=1 \\ y=-2. \end{cases} \]
Le vecteur $\vec{v}(5\,;1)$ admet un \textbf{unique} antécédent : $\vec{u}(1\,;-2)$. Vérification : $f(1,-2)=(1+4\,;\,3-2)=(5,1)$. \checkmark
\end{enumerate}
\end{corrige}

\begin{corrige}[Exercice 8 — Écriture analytique et image d'un vecteur]
\begin{enumerate}
\item Par linéarité, pour $\vec{u}=x\vec{i}+y\vec{j}$ :
\begin{align*}
f(\vec{u}) &= x\,f(\vec{i})+y\,f(\vec{j}) = x(\vec{i}-2\vec{j})+y(3\vec{i}+\vec{j})\\
&= (x+3y)\vec{i}+(-2x+y)\vec{j}.
\end{align*}
Donc $f(x\,;y)=(x+3y\,;\,-2x+y)$.
\item $f(2,-5)=\bigl(2+3\times(-5)\,;\,-2\times 2+(-5)\bigr)=(2-15\,;\,-4-5)=(-13\,;\,-9)$.
\item La matrice de $f$ dans $(\vec{i},\vec{j})$ est $A=\begin{pmatrix} 1 & 3 \\ -2 & 1 \end{pmatrix}$ (les \emph{colonnes} sont les images de $\vec{i}$ et $\vec{j}$).
\[ ad-bc = 1\times 1 - 3\times(-2) = 1+6 = 7 \neq 0. \]
Donc $f$ est bijective : c'est un \textbf{automorphisme} du plan.
\end{enumerate}
\end{corrige}

\begin{attention}[Point de vigilance — Exercice 8]
Dans la matrice, les images des vecteurs de base se rangent en \textbf{colonnes} : $f(\vec{i})=(1,-2)$ donne la première colonne, $f(\vec{j})=(3,1)$ la deuxième. Une transposition changera le signe du déterminant... ou pire.
\end{attention}

\begin{corrige}[Exercice 9 — Noyau et image]
\begin{enumerate}
\item $\operatorname{Ker}(f)$ : on résout
\[ \begin{cases} 2x-4y=0 \\ -x+2y=0 \end{cases} \iff x=2y \]
(les deux équations sont proportionnelles). Donc
\[ \operatorname{Ker}(f)=\{(2t\,;\,t)\mid t\in\mathbb{R}\}=\operatorname{Vect}\bigl((2\,;\,1)\bigr). \]
Base : $\bigl\{(2\,;\,1)\bigr\}$ ; équation cartésienne : $x-2y=0$.
\item $\operatorname{Im}(f)=\{(2x-4y\,;\,-x+2y)\mid x,y\in\mathbb{R}\}$. En posant $v=-x+2y$, la première coordonnée vaut $-2v$ :
\[ \operatorname{Im}(f)=\{(-2v\,;\,v)\mid v\in\mathbb{R}\}=\operatorname{Vect}\bigl((-2\,;\,1)\bigr). \]
Base : $\bigl\{(-2\,;\,1)\bigr\}$ ; équation cartésienne : $x+2y=0$.
\item Les vecteurs directeurs $(2,1)$ et $(-2,1)$ ne sont pas colinéaires ($2\times 1-1\times(-2)=4\neq 0$) : les deux droites sont \textbf{distinctes}. Leur produit scalaire vaut $2\times(-2)+1\times 1=-3\neq 0$ : elles ne sont pas perpendiculaires. Géométriquement, $f$ « écrase » le plan sur la droite $\operatorname{Im}(f)$ en suivant la direction de $\operatorname{Ker}(f)$ : c'est une \textbf{projection oblique} (non orthogonale) sur $\operatorname{Im}(f)$ parallèlement à $\operatorname{Ker}(f)$.
\item \textbf{Non}, $f$ n'est pas bijective. Trois justifications possibles :
\begin{itemize}
\item $\operatorname{Ker}(f)\neq\{\vec{0}\}$ : $f$ n'est pas injective ;
\item $\operatorname{Im}(f)\neq\mathbb{R}^2$ : $f$ n'est pas surjective ;
\item $ad-bc=2\times 2-(-4)\times(-1)=4-4=0$.
\end{itemize}
On vérifie au passage le théorème du rang : $\dim\operatorname{Ker}(f)+\dim\operatorname{Im}(f)=1+1=2$. \checkmark
\end{enumerate}
\end{corrige}

\begin{corrige}[Exercice 10 — Étude de la bijectivité de deux façons]
\begin{enumerate}
\item Soit $\vec{u}(x_1,y_1)$, $\vec{v}(x_2,y_2)$ et $\lambda\in\mathbb{R}$.
\begin{align*}
f(\vec{u}+\vec{v}) &= \bigl((x_1+x_2)+3(y_1+y_2)\,;\,2(x_1+x_2)-(y_1+y_2)\bigr)\\
&= (x_1+3y_1\,;\,2x_1-y_1)+(x_2+3y_2\,;\,2x_2-y_2) = f(\vec{u})+f(\vec{v}) ;\\
f(\lambda\vec{u}) &= (\lambda x_1+3\lambda y_1\,;\,2\lambda x_1-\lambda y_1) = \lambda f(\vec{u}).
\end{align*}
Donc $f$ est linéaire. (On pouvait aussi invoquer directement la forme $(ax+by\,;\,cx+dy)$.)
\item \textbf{Noyau :} $\begin{cases} x+3y=0 \\ 2x-y=0 \end{cases}$. La deuxième équation donne $y=2x$ ; en reportant : $x+6x=7x=0$, donc $x=0$ puis $y=0$. Ainsi $\operatorname{Ker}(f)=\{\vec{0}\}$.\\
\textbf{Image :} par le théorème du rang, $\dim\operatorname{Im}(f)=2-0=2$, donc $\operatorname{Im}(f)=\mathbb{R}^2$ (le plan entier).
\item \textbf{Oui, $f$ est bijective} (automorphisme) :
\begin{itemize}
\item \emph{Par le noyau :} $\operatorname{Ker}(f)=\{\vec{0}\}$ donc $f$ est injective ; un endomorphisme injectif du plan est bijectif.
\item \emph{Par l'écriture analytique :} $ad-bc=1\times(-1)-3\times 2=-1-6=-7\neq 0$, donc $f$ est bijective.
\end{itemize}
\end{enumerate}
\end{corrige}

\begin{corrige}[Exercice 11 — Projection vectorielle]
\begin{enumerate}
\item On projette sur $D:y=x$ parallèlement à l'axe des ordonnées (verticalement) : l'abscisse est conservée, et l'image est le point de $D$ ayant cette abscisse.
\begin{itemize}
\item $\vec{i}=(1,0)$ : le point de $D$ d'abscisse $1$ est $(1,1)$, donc $p(\vec{i})=(1\,;\,1)=\vec{i}+\vec{j}$.
\item $\vec{j}=(0,1)$ : le point de $D$ d'abscisse $0$ est $(0,0)$, donc $p(\vec{j})=\vec{0}$.
\end{itemize}
\item Par linéarité : $p(x\vec{i}+y\vec{j})=x\,p(\vec{i})+y\,p(\vec{j})=x(\vec{i}+\vec{j})$, d'où
\[ p(x\,;y)=(x\,;\,x). \]
\item $\operatorname{Ker}(p)=\{(x,y)\mid x=0\}$ : c'est l'\textbf{axe des ordonnées}, la direction de projection ; base $\{(0\,;\,1)\}$.\\
$\operatorname{Im}(p)=\{(x,x)\mid x\in\mathbb{R}\}$ : c'est la \textbf{droite $D$} sur laquelle on projette ; base $\{(1\,;\,1)\}$.
\item $p\circ p(x,y)=p(x,x)=(x,x)=p(x,y)$, donc $p\circ p=p$. Ce résultat était \textbf{prévisible} : une projection est idempotente — un vecteur déjà sur $D$ est invariant par une nouvelle projection sur $D$.
\item \textbf{Non} : $\operatorname{Ker}(p)\neq\{\vec{0}\}$ (non injective) et $\operatorname{Im}(p)\neq\mathbb{R}^2$ (non surjective). Une projection (non triviale) n'est jamais un automorphisme.
\end{enumerate}
\end{corrige}

\begin{corrige}[Exercice 12 — Type Bac : la société « TransLittoral »]
\begin{enumerate}
\item $f(x,y)=(x+3y\,;\,2x-y)$ est de la forme $(ax+by\,;\,cx+dy)$ avec $a=1$, $b=3$, $c=2$, $d=-1$ : c'est une application linéaire du plan dans lui-même.
\item $f(7,-1)=\bigl(7+3\times(-1)\,;\,2\times 7-(-1)\bigr)=(4\,;\,15)$. Le vecteur corrigé est $(4\,;\,15)$.
\item $\operatorname{Ker}(f)$ : $\begin{cases} x+3y=0 \\ 2x-y=0 \end{cases}$. De $y=2x$ on tire $x+6x=7x=0$, donc $x=y=0$. Ainsi $\operatorname{Ker}(f)=\{\vec{0}\}$ (dimension $0$, pas de base non vide).
\item Théorème du rang : $\dim\operatorname{Im}(f)=2-0=2$, donc $\operatorname{Im}(f)=\mathbb{R}^2$.
\item \textbf{Oui, $f$ est un automorphisme :}
\begin{enumerate}
\item \emph{Par le noyau :} $\operatorname{Ker}(f)=\{\vec{0}\}$ donc $f$ est injective, donc bijective (endomorphisme du plan).
\item \emph{Par l'écriture analytique :} $ad-bc=1\times(-1)-3\times 2=-7\neq 0$, donc $f$ est bijective.
\end{enumerate}
\item On résout $f(x,y)=(10,6)$ :
\[ \begin{cases} x+3y=10 \\ 2x-y=6 \end{cases} \iff \begin{cases} y=2x-6 \\ x+3(2x-6)=10 \end{cases} \iff \begin{cases} 7x=28 \\ y=2x-6 \end{cases} \iff \begin{cases} x=4 \\ y=2. \end{cases} \]
Antécédent unique : $(4\,;\,2)$. \textbf{Interprétation :} $f$ étant bijective, tout trajet corrigé provient d'\textbf{un seul} trajet réel : l'ingénieur peut remonter sans ambiguïté du trajet corrigé au trajet initial.
\end{enumerate}

\textbf{Barème indicatif (sur 6 pts) :} 1 : 0,5 pt ; 2 : 0,5 pt ; 3 : 1 pt ; 4 : 1 pt ; 5 : 1,5 pt (0,75 par méthode) ; 6 : 1,5 pt (1 pt calcul + 0,5 pt interprétation).
\end{corrige}

\begin{attention}[Points de vigilance — Bac (Exercice 12)]
\begin{itemize}
\item En 5.a), il faut citer explicitement : « $\operatorname{Ker}(f)=\{\vec{0}\}$ donc $f$ injective, donc bijective car endomorphisme du plan ».
\item En 5.b), écrire le calcul $ad-bc$ en entier et conclure par « $\neq 0$ donc bijective » : le déterminant seul sans conclusion ne vaut pas la totalité des points.
\item En 6), l'interprétation en langage courant est attendue : ne pas s'arrêter au couple $(4\,;\,2)$.
\end{itemize}
\end{attention}

\begin{corrige}[Exercice 13 — Type Bac : famille d'applications à paramètre]
\begin{enumerate}
\item Pour tout réel $m$, $f_m(x,y)=(mx+y\,;\,x+my)$ est de la forme $(ax+by\,;\,cx+dy)$ avec $a=m$, $b=1$, $c=1$, $d=m$ : $f_m$ est linéaire.
\item Cas $m=2$ :
\begin{enumerate}
\item $f_2(x,y)=(2x+y\,;\,x+2y)$.
\item $ad-bc=2\times 2-1\times 1=3\neq 0$ : $f_2$ est bijective, c'est un automorphisme.
\item On résout $\begin{cases} 2x+y=4 \\ x+2y=5 \end{cases}$ : $y=4-2x$, puis $x+2(4-2x)=5 \iff -3x=-3 \iff x=1$, d'où $y=2$. Antécédent : $(1\,;\,2)$. Vérification : $f_2(1,2)=(4,5)$. \checkmark
\end{enumerate}
\item Cas $m=1$ : $f_1(x,y)=(x+y\,;\,x+y)$.
\begin{enumerate}
\item $\operatorname{Ker}(f_1)=\{(x,y)\mid x+y=0\}$ : droite d'équation $x+y=0$, base $\bigl\{(1\,;\,-1)\bigr\}$.
\item $\operatorname{Im}(f_1)=\{(t,t)\mid t\in\mathbb{R}\}$ : droite d'équation $y=x$, base $\bigl\{(1\,;\,1)\bigr\}$.
\item \textbf{Non} : $\operatorname{Ker}(f_1)\neq\{\vec{0}\}$ (non injective), $\operatorname{Im}(f_1)\neq\mathbb{R}^2$ (non surjective), et $ad-bc=1-1=0$.
\end{enumerate}
\item $f_m$ est un automorphisme $\iff ad-bc\neq 0 \iff m^2-1\neq 0 \iff (m-1)(m+1)\neq 0 \iff m\neq 1$ et $m\neq -1$.
\item Valeurs exclues : $m=1$ et $m=-1$.
\begin{itemize}
\item $m=1$ : $\operatorname{Ker}(f_1)$ : $x+y=0$, base $\{(1,-1)\}$ ; $\operatorname{Im}(f_1)$ : $y=x$, base $\{(1,1)\}$.
\item $m=-1$ : $f_{-1}(x,y)=(-x+y\,;\,x-y)$. $\operatorname{Ker}(f_{-1})$ : $y=x$, base $\{(1,1)\}$ ; $\operatorname{Im}(f_{-1})=\{(t,-t)\mid t\in\mathbb{R}\}$ : $y=-x$, base $\{(1,-1)\}$.
\end{itemize}
Remarque : dans les deux cas, noyau et image sont des droites perpendiculaires, échangées entre $m=1$ et $m=-1$.
\end{enumerate}

\textbf{Barème indicatif (sur 7 pts) :} 1 : 0,5 pt ; 2 : 1,5 pt (0,5 par sous-question) ; 3 : 2 pts ; 4 : 1,5 pt ; 5 : 1,5 pt.
\end{corrige}

\begin{attention}[Points de vigilance — Bac (Exercice 13)]
\begin{itemize}
\item Question 4 : la condition porte sur $m^2-1\neq 0$ ; factoriser en $(m-1)(m+1)$ pour ne pas oublier la valeur $m=-1$.
\item Question 5 : traiter \textbf{les deux} valeurs $m=1$ et $m=-1$ séparément, avec équation \emph{et} base à chaque fois.
\item Ne pas confondre « $f_m$ n'est pas un automorphisme » avec « $f_m$ est nulle » : pour $m=\pm 1$, $f_m$ est de rang 1.
\end{itemize}
\end{attention}

\begin{corrige}[Exercice 14 — Type Bac : symétrie, projection et composition]
\begin{enumerate}
\item $s(x,y)=(y,x)=(0\cdot x+1\cdot y\,;\,1\cdot x+0\cdot y)$ et $p(x,y)=\left(\tfrac12 x+\tfrac12 y\,;\,\tfrac12 x+\tfrac12 y\right)$ sont de la forme $(ax+by\,;\,cx+dy)$ : ce sont des applications linéaires.
\item Étude de $s$ :
\begin{enumerate}
\item $\operatorname{Ker}(s)=\{(x,y)\mid (y,x)=(0,0)\}=\{\vec{0}\}$. $\operatorname{Im}(s)=\mathbb{R}^2$ car tout couple $(u,v)$ est l'image de $(v,u)$.
\item $\operatorname{Ker}(s)=\{\vec{0}\}$ donc $s$ est injective, donc bijective : $s$ est un \textbf{automorphisme}. (Autre preuve : $ad-bc=0\times 0-1\times 1=-1\neq 0$.)
\item $s\circ s(x,y)=s(y,x)=(x,y)$ : $s\circ s=\mathrm{Id}$. $s$ est une \textbf{symétrie vectorielle} (involution) : géométriquement, la symétrie orthogonale par rapport à $D:y=x$.
\end{enumerate}
\item Étude de $p$ :
\begin{enumerate}
\item $\operatorname{Ker}(p)=\{(x,y)\mid x+y=0\}$ : c'est la droite $\Delta$ (le trottoir) ; base $\{(1\,;\,-1)\}$.
\item $\operatorname{Im}(p)=\{(t,t)\mid t\in\mathbb{R}\}$ : c'est la droite $D$ (l'avenue) ; base $\{(1\,;\,1)\}$.
\item $p\circ p(x,y)=p\left(\tfrac{x+y}{2},\tfrac{x+y}{2}\right)=\left(\tfrac{x+y}{2},\tfrac{x+y}{2}\right)=p(x,y)$ : $p\circ p=p$. $p$ est une \textbf{projection vectorielle} (idempotente) : projection sur $D$ parallèlement à $\Delta$.
\item \textbf{Non} : (i) $\operatorname{Ker}(p)=\Delta\neq\{\vec{0}\}$ donc non injective ; (ii) $\operatorname{Im}(p)=D\neq\mathbb{R}^2$ donc non surjective (et $ad-bc=\tfrac14-\tfrac14=0$).
\end{enumerate}
\item $g=s\circ p$ :
\begin{enumerate}
\item $g(x,y)=s\left(\tfrac{x+y}{2},\tfrac{x+y}{2}\right)=\left(\tfrac{x+y}{2},\tfrac{x+y}{2}\right)=p(x,y)$ : donc $g=p$.
\item \textbf{Non} : $g=p$ n'est pas bijective (noyau $\Delta$, image $D$, déterminant nul).
\end{enumerate}
\item \textbf{Interprétation :} $D:x-y=0$ est l'avenue, $\Delta:x+y=0$ le trottoir (perpendiculaire à l'avenue). $s$ est la symétrie orthogonale d'axe l'avenue $D$ : les points de $D$ sont invariants, ceux de $\Delta$ sont envoyés sur leur opposé. $p$ est la projection orthogonale sur l'avenue $D$ parallèlement au trottoir $\Delta$ : le noyau est le trottoir, l'image l'avenue. Enfin $s\circ p=p$ : projeter sur l'avenue puis symétriser par rapport à l'avenue ne change rien, puisque tout point de l'avenue est invariant par $s$.
\end{enumerate}

\textbf{Barème indicatif (sur 8 pts) :} 1 : 0,5 pt ; 2 : 2 pts ; 3 : 3 pts ; 4 : 1,5 pt ; 5 : 1 pt.
\end{corrige}

\begin{attention}[Points de vigilance — Bac (Exercice 14)]
\begin{itemize}
\item Pour $s\circ s=\mathrm{Id}$ et $p\circ p=p$, rédiger le calcul complet de la composée, pas seulement la conclusion.
\item En 4.b), le fait que $g=p$ ne dispense pas de justifier la non-bijectivité (noyau, image ou déterminant).
\item En 5), nommer précisément les droites ($D$ avenue, $\Delta$ trottoir) et le type de chaque transformation : symétrie orthogonale d'axe $D$, projection sur $D$ parallèlement à $\Delta$.
\end{itemize}
\end{attention}

\newpage

\coursec{Fiche bilan}

\coursec{Applications linéaires du plan}

\courssub{Définition et premiers exemples d'applications linéaires}

\begin{experience}[Intuition : transformer le plan sans le déchirer]
Imagine que tu as une feuille de papier quadrillé (le plan vectoriel $\mathcal{V}$) et que tu l'étires, le compresses, le fais tourner ou le projettes sur une droite, \emph{sans jamais le plier ni le déchirer}. L'origine $O$ reste fixe, les droites restent des droites, le parallélisme est conservé. C'est exactement ce que font les \cle{applications linéaires} : elles respectent la structure vectorielle (addition et multiplication par un réel). Au Cameroun, on utilise ces transformations pour la cartographie (projection de la Terre sur une carte), le traitement d'images numériques (rotation, redimensionnement chez MTN/Orange) ou la robotique (calcul de trajectoires).
\end{experience}

\begin{definition}[Application linéaire (Endomorphisme)]
Soit $\mathcal{V}$ le plan vectoriel muni d'un repère. Une application $f : \mathcal{V} \to \mathcal{V}$ est dite \cle{linéaire} (ou \cle{endomorphisme} de $\mathcal{V}$) si elle vérifie les deux propriétés suivantes pour \textbf{tous} vecteurs $\vec{u}, \vec{v} \in \mathcal{V}$ et \textbf{tout} réel $\lambda \in \mathbb{R}$ :
\begin{enumerate}
  \item \textbf{Additivité} : $f(\vec{u} + \vec{v}) = f(\vec{u}) + f(\vec{v})$.
  \item \textbf{Homogénéité} : $f(\lambda \vec{u}) = \lambda f(\vec{u})$.
\end{enumerate}
On note $\mathcal{L}(\mathcal{V})$ l'ensemble des endomorphismes de $\mathcal{V}$.
\end{definition}

\begin{propriete}[Propriétés immédiates et critère unique]
Soit $f : \mathcal{V} \to \mathcal{V}$ une application linéaire.
\begin{itemize}
  \item $f(\vec{0}) = \vec{0}$ (l'origine est invariante).
  \item $f(-\vec{u}) = -f(\vec{u})$.
  \item \textbf{Critère unique (souvent plus pratique) :} $f$ est linéaire $\iff$ pour tous $\vec{u}, \vec{v} \in \mathcal{V}$ et tous $\lambda, \mu \in \mathbb{R}$,
  \[ f(\lambda \vec{u} + \mu \vec{v}) = \lambda f(\vec{u}) + \mu f(\vec{v}). \]
\end{itemize}
\end{propriete}

\begin{attention}[Piège classique : Translation et applications affines]
Une \textbf{translation} $t_{\vec{a}}(\vec{u}) = \vec{u} + \vec{a}$ (avec $\vec{a} \neq \vec{0}$) n'est \textbf{pas} linéaire car $t_{\vec{a}}(\vec{0}) = \vec{a} \neq \vec{0}$.
Plus généralement, une application de la forme $f(\vec{u}) = A\vec{u} + \vec{b}$ (avec $\vec{b} \neq \vec{0}$) est \cle{affine}, pas linéaire. \textbf{Au Bac :} Vérifie toujours $f(\vec{0})=\vec{0}$ en premier ; si c'est faux, l'application n'est pas linéaire (gain de temps garanti).
\end{attention}

\begin{exemplebox}[Exemples géométriques fondamentaux (à connaître par cœur)]
Dans la base canonique $(\vec{i}, \vec{j})$, soit $\vec{u} = x\vec{i} + y\vec{j}$.
\begin{itemize}
  \item \textbf{Homothétie vectorielle} de rapport $k \in \mathbb{R}$ : $h_k(\vec{u}) = k\vec{u}$. Matrice $\begin{pmatrix} k & 0 \\ 0 & k \end{pmatrix}$.
  \item \textbf{Symétrie orthogonale} par rapport à l'axe $(Ox)$ : $s_x(x\vec{i}+y\vec{j}) = x\vec{i} - y\vec{j}$. Matrice $\begin{pmatrix} 1 & 0 \\ 0 & -1 \end{pmatrix}$.
  \item \textbf{Projection orthogonale} sur l'axe $(Ox)$ : $p_x(x\vec{i}+y\vec{j}) = x\vec{i}$. Matrice $\begin{pmatrix} 1 & 0 \\ 0 & 0 \end{pmatrix}$.
  \item \textbf{Rotation vectorielle} d'angle $\alpha$ : $r_\alpha(x\vec{i}+y\vec{j}) = (x\cos\alpha - y\sin\alpha)\vec{i} + (x\sin\alpha + y\cos\alpha)\vec{j}$. Matrice $\begin{pmatrix} \cos\alpha & -\sin\alpha \\ \sin\alpha & \cos\alpha \end{pmatrix}$.
\end{itemize}
Tous ces exemples vérifient $f(\lambda\vec{u}+\mu\vec{v}) = \lambda f(\vec{u})+\mu f(\vec{v})$.
\end{exemplebox}

\begin{methode}[Prouver qu'une application est linéaire]
\begin{enumerate}
  \item \textbf{Test rapide :} Vérifier $f(\vec{0}) = \vec{0}$. Si faux $\to$ \textbf{Non linéaire} (arrêt).
  \item \textbf{Preuve formelle :} Prendre $\vec{u}, \vec{v}$ quelconques et $\lambda, \mu \in \mathbb{R}$.
  \item Calculer $f(\lambda\vec{u} + \mu\vec{v})$ en utilisant la définition de $f$.
  \item Calculer $\lambda f(\vec{u}) + \mu f(\vec{v})$.
  \item Comparer : si égalité $\to$ \textbf{Linéaire}, sinon \textbf{Non linéaire}.
  \item \textbf{Astuce :} Si $f$ s'écrit sous forme analytique $f(x,y) = (ax+by, cx+dy)$ \textbf{sans terme constant}, elle est linéaire.
\end{enumerate}
\end{methode}

\begin{exoresolu}[Reconnaissance de linéarité]
Soient $f_1, f_2 : \mathbb{R}^2 \to \mathbb{R}^2$ définies par :
$f_1(x,y) = (2x+y, x-3y)$ et $f_2(x,y) = (x^2, y+1)$.
\tcblower
\textbf{$f_1$ :} $f_1(0,0) = (0,0)$. Forme analytique sans terme constant $\Rightarrow$ \textbf{Linéaire}.
\textbf{$f_2$ :} $f_2(0,0) = (0,1) \neq (0,0)$ $\Rightarrow$ \textbf{Non linéaire}. (On s'arrête là).
\end{exoresolu}

\courssub{Endomorphismes, isomorphismes, automorphismes}

\begin{definition}[Vocabulaire précis pour le Bac]
\begin{itemize}
  \item \cle{Endomorphisme} : Application linéaire d'un espace vectoriel dans lui-même ($f: \mathcal{V} \to \mathcal{V}$).
  \item \cle{Isomorphisme} : Application linéaire \textbf{bijective} entre deux espaces vectoriels (ici $\mathcal{V} \to \mathcal{V}$).
  \item \cle{Automorphisme} : Endomorphisme bijectif. C'est un isomorphisme de $\mathcal{V}$ sur $\mathcal{V}$.
\end{itemize}
\textbf{Attention :} Tout automorphisme est un endomorphisme, mais l'inverse est faux (ex: projection).
\end{definition}

\begin{propriete}[Groupe linéaire]
L'ensemble des automorphismes de $\mathcal{V}$, noté $GL(\mathcal{V})$ ou $GL_2(\mathbb{R})$, est un \textbf{groupe} pour la composition :
\begin{itemize}
  \item Stable par composition : $f, g$ automorphismes $\implies f \circ g$ automorphisme.
  \item Tout automorphisme $f$ admet une application réciproque $f^{-1}$ qui est \textbf{aussi} un automorphisme (donc linéaire).
\end{itemize}
\end{propriete}

\courssub{Expression d'une application linéaire dans une base}

\begin{propriete}[Matrice d'un endomorphisme dans une base]
Soit $\mathcal{B} = (\vec{e_1}, \vec{e_2})$ une base de $\mathcal{V}$ et $f \in \mathcal{L}(\mathcal{V})$.
Il existe quatre réels $a, b, c, d$ tels que :
\[ f(\vec{e_1}) = a\vec{e_1} + c\vec{e_2} \quad\text{et}\quad f(\vec{e_2}) = b\vec{e_1} + d\vec{e_2}. \]
La \cle{matrice de $f$ dans la base $\mathcal{B}$} est :
\[ M_{\mathcal{B}}(f) = \begin{pmatrix} a & b \\ c & d \end{pmatrix}. \]
\textbf{Règle d'or :} Les \textbf{colonnes} de la matrice sont les coordonnées des images des vecteurs de la base.
Pour tout vecteur $\vec{u} = x\vec{e_1} + y\vec{e_2}$, on a l'\cle{écriture analytique} :
\[ f(\vec{u}) = (ax+by)\vec{e_1} + (cx+dy)\vec{e_2} \quad\Longleftrightarrow\quad \begin{pmatrix} x' \\ y' \end{pmatrix} = \begin{pmatrix} a & b \\ c & d \end{pmatrix} \begin{pmatrix} x \\ y \end{pmatrix}. \]
\end{propriete}

\begin{exemplebox}[Déterminer la matrice et calculer une image]
Soit $f$ l'endomorphisme tel que dans la base $\mathcal{B}=(\vec{i},\vec{j})$ :
$f(\vec{i}) = 2\vec{i} + \vec{j}$ et $f(\vec{j}) = -\vec{i} + 3\vec{j}$.
\tcblower
\textbf{1. Matrice :} Les colonnes sont $f(\vec{i})$ et $f(\vec{j})$.
$M_{\mathcal{B}}(f) = \begin{pmatrix} 2 & -1 \\ 1 & 3 \end{pmatrix}$.
\textbf{2. Image de $\vec{u} = 3\vec{i} - 2\vec{j}$ :}
$\begin{pmatrix} x' \\ y' \end{pmatrix} = \begin{pmatrix} 2 & -1 \\ 1 & 3 \end{pmatrix} \begin{pmatrix} 3 \\ -2 \end{pmatrix} = \begin{pmatrix} 6+2 \\ 3-6 \end{pmatrix} = \begin{pmatrix} 8 \\ -3 \end{pmatrix}$.
Donc $f(\vec{u}) = 8\vec{i} - 3\vec{j}$.
\end{exemplebox}

\begin{methode}[Expression analytique et matrice dans une base $\mathcal{B}=(\vec{u},\vec{v})$]
\begin{enumerate}
  \item Calculer $f(\vec{u})$ et $f(\vec{v})$ et les exprimer dans $\mathcal{B}$ : $f(\vec{u}) = a\vec{u}+c\vec{v}$, $f(\vec{v}) = b\vec{u}+d\vec{v}$.
  \item Écrire la matrice $M = \begin{pmatrix} a & b \\ c & d \end{pmatrix}$.
  \item Pour tout $\vec{x} = x\vec{u}+y\vec{v}$, $f(\vec{x}) = (ax+by)\vec{u} + (cx+dy)\vec{v}$.
\end{enumerate}
\end{methode}

\courssub{Noyau d'une application linéaire}

\begin{definition}[Noyau]
Le \cle{noyau} de $f$ est l'ensemble des vecteurs envoyés sur $\vec{0}$ :
\[ \ker f = \{ \vec{u} \in \mathcal{V} \mid f(\vec{u}) = \vec{0} \}. \]
\end{definition}

\begin{propriete}[Structure du noyau]
$\ker f$ est un \textbf{sous-espace vectoriel} de $\mathcal{V}$ (donc stable par addition et multiplication par un réel).
Sa dimension (le \cle{défaut} de $f$) ne peut être que $0$, $1$ ou $2$ :
\begin{itemize}
  \item $\dim(\ker f) = 0 \iff \ker f = \{\vec{0}\}$ ($f$ injective).
  \item $\dim(\ker f) = 1 \iff \ker f$ est une \textbf{droite vectorielle} (noyau non trivial).
  \item $\dim(\ker f) = 2 \iff \ker f = \mathcal{V}$ ($f$ est l'application nulle).
\end{itemize}
\end{propriete}

\begin{methode}[Déterminer $\ker f$ et une base]
Soit $M = \begin{pmatrix} a & b \\ c & d \end{pmatrix}$ la matrice de $f$ dans une base $\mathcal{B}$.
\begin{enumerate}
  \item Résoudre le système homogène : $\begin{cases} ax + by = 0 \\ cx + dy = 0 \end{cases}$.
  \item \textbf{Cas 1 :} Solution unique $(0,0)$ $\implies \ker f = \{\vec{0}\}$. Base : $\varnothing$ (ou "pas de base non vide").
  \item \textbf{Cas 2 :} Une infinité de solutions (équations proportionnelles). Exprimer $(x,y)$ en fonction d'un paramètre $t$ : $(x,y) = t(x_0, y_0)$.
  \item \textbf{Conclusion :} $\ker f = \operatorname{Vect}(\vec{u_0})$ où $\vec{u_0} = x_0\vec{e_1}+y_0\vec{e_2}$. Une base est $(\vec{u_0})$.
\end{enumerate}
\end{methode}

\begin{exoresolu}[Calcul de noyau]
Soit $f$ de matrice $M = \begin{pmatrix} 2 & -1 \\ 4 & -2 \end{pmatrix}$ dans la base canonique.
\tcblower
Système : $\begin{cases} 2x - y = 0 \\ 4x - 2y = 0 \end{cases} \iff y = 2x$.
Les solutions sont $(x, 2x) = x(1, 2)$.
$\ker f = \operatorname{Vect}(\vec{i} + 2\vec{j})$. Une base : $(\vec{i} + 2\vec{j})$. $\dim(\ker f) = 1$.
\end{exoresolu}

\courssub{Image d'une application linéaire et stabilité}

\begin{definition}[Image]
L'\cle{image} de $f$ est l'ensemble des vecteurs atteints :
\[ \operatorname{Im} f = \{ f(\vec{u}) \mid \vec{u} \in \mathcal{V} \}. \]
\end{definition}

\begin{propriete}[Structure de l'image et théorème du rang]
$\operatorname{Im} f$ est un \textbf{sous-espace vectoriel} de $\mathcal{V}$ (stable par addition et multiplication par un réel).
Si $M = \begin{pmatrix} a & b \\ c & d \end{pmatrix}$ dans la base $\mathcal{B}=(\vec{e_1},\vec{e_2})$, alors :
\[ \operatorname{Im} f = \operatorname{Vect}\big( f(\vec{e_1}), f(\vec{e_2}) \big) = \operatorname{Vect}\left( \begin{pmatrix} a \\ c \end{pmatrix}, \begin{pmatrix} b \\ d \end{pmatrix} \right). \]
\textbf{Théorème du rang (Dimension) :} $\dim(\ker f) + \dim(\operatorname{Im} f) = \dim(\mathcal{V}) = 2$.
Le \cle{rang} de $f$ est $\dim(\operatorname{Im} f)$ ($0, 1$ ou $2$).
\end{propriete}

\begin{methode}[Déterminer $\operatorname{Im} f$, une base et une équation caractéristique]
\begin{enumerate}
  \item Les vecteurs générateurs sont les \textbf{colonnes} de la matrice $M$ : $\vec{c_1} = \begin{pmatrix} a \\ c \end{pmatrix}$, $\vec{c_2} = \begin{pmatrix} b \\ d \end{pmatrix}$.
  \item \textbf{Si $\vec{c_1}$ et $\vec{c_2}$ ne sont pas colinéaires (rang 2) :} $\operatorname{Im} f = \mathcal{V}$. Base : $(\vec{c_1}, \vec{c_2})$.
  \item \textbf{Si $\vec{c_1}$ et $\vec{c_2}$ sont colinéaires et non nuls (rang 1) :} $\operatorname{Im} f$ est la droite vectorielle dirigée par $\vec{c_1}$ (ou $\vec{c_2}$). Base : $(\vec{c_1})$.
    \begin{itemize}
      \item \textbf{Équation caractéristique (cartésienne) :} Un vecteur $\begin{pmatrix} x \\ y \end{pmatrix}$ appartient à $\operatorname{Im} f$ s'il est colinéaire à $\vec{c_1}$, c'est-à-dire si son déterminant avec $\vec{c_1}$ est nul : $\det(\vec{c_1}, \vec{x}) = 0 \iff cx - ay = 0$ (si $\vec{c_1} = \binom{a}{c}$).
    \end{itemize}
  \item \textbf{Si $M$ est la matrice nulle (rang 0) :} $\operatorname{Im} f = \{\vec{0}\}$.
\end{enumerate}
\end{methode}

\begin{exemplebox}[Image et équation caractéristique]
$f$ de matrice $M = \begin{pmatrix} 1 & 2 \\ 2 & 4 \end{pmatrix}$.
\tcblower
Colonnes : $\vec{c_1} = \binom{1}{2}$, $\vec{c_2} = \binom{2}{4} = 2\vec{c_1}$. Colinéaires $\implies$ Rang 1.
$\operatorname{Im} f = \operatorname{Vect}(\binom{1}{2})$ (droite $y=2x$).
Base : $\left( \binom{1}{2} \right)$.
Équation caractéristique : $\det\left( \binom{1}{2}, \binom{x}{y} \right) = 1\cdot y - 2\cdot x = 0 \iff y - 2x = 0$.
\end{exemplebox}

\begin{attention}[Erreurs fréquentes sur l'image]
\begin{itemize}
  \item Confondre les colonnes et les lignes de la matrice pour l'image. \textbf{Image = espace engendré par les COLONNES}.
  \item Oublier de vérifier si les colonnes sont colinéaires avant de donner une base (deux vecteurs colinéaires ne forment pas une base).
  \item Pour l'équation caractéristique : utiliser le déterminant avec un vecteur directeur de la droite image, pas avec un vecteur quelconque.
\end{itemize}
\end{attention}

\courssub{Caractérisation de la bijectivité et synthèse}

\begin{propriete}[Caractérisations de la bijectivité (Théorème fondamental)]
Pour un endomorphisme $f$ de matrice $M = \begin{pmatrix} a & b \\ c & d \end{pmatrix}$ dans une base, les propositions suivantes sont \textbf{équivalentes} :
\begin{enumerate}
  \item $f$ est \textbf{bijective} (automorphisme).
  \item $f$ est \textbf{injective} $\iff \ker f = \{\vec{0}\}$.
  \item $f$ est \textbf{surjective} $\iff \operatorname{Im} f = \mathcal{V}$.
  \item $\det(M) = ad - bc \neq 0$.
  \item Les colonnes de $M$ sont \textbf{linéairement indépendantes} (forment une base de $\mathcal{V}$).
  \item Les lignes de $M$ sont \textbf{linéairement indépendantes}.
\end{enumerate}
Si $f$ est bijective, sa réciproque $f^{-1}$ est linéaire et sa matrice dans la même base est $M^{-1} = \frac{1}{ad-bc} \begin{pmatrix} d & -b \\ -c & a \end{pmatrix}$.
\end{propriete}

\begin{methode}[Étudier la bijectivité d'un endomorphisme]
\begin{enumerate}
  \item Écrire la matrice $M = \begin{pmatrix} a & b \\ c & d \end{pmatrix}$ de $f$ dans la base donnée.
  \item Calculer le déterminant $\Delta = ad - bc$.
  \item \textbf{Si $\Delta \neq 0$ :} $f$ est bijective (automorphisme). On peut calculer $M^{-1}$ si demandé.
  \item \textbf{Si $\Delta = 0$ :} $f$ n'est pas bijective. Préciser : $\ker f \neq \{\vec{0}\}$ (non injective) et $\operatorname{Im} f \neq \mathcal{V}$ (non surjective). Calculer $\ker f$ et $\operatorname{Im} f$ pour compléter l'étude.
\end{enumerate}
\end{methode}

\begin{savaistu}[Applications concrètes au Cameroun et au-delà]
\begin{itemize}
  \item \textbf{Télécoms (MTN/Orange) :} Le codage/décodage de signaux (OFDM, QAM) utilise des transformations linéaires (matrices de Fourier, matrices de Hadamard) pour multiplexer les données.
  \item \textbf{Cartographie / SIG :} Le passage de coordonnées GPS (WGS84) à une projection plane locale (UTM, Lambert) pour les cartes du MINHDU ou l'électrification rurale (AER) fait intervenir des applications affines/linéaires locales.
  \item \textbf{Infographie / Jeux vidéo :} Rotation, zoom, déformation des sprites 2D sont des multiplications matricielles (endomorphismes) appliquées 60 fois par seconde.
  \item \textbf{Économie / Input-Output (Leontief) :} Modélisation des flux inter-secteurs (agriculture, industrie, services) par un système linéaire $X = AX + D$.
\end{itemize}
\end{savaistu}

\begin{popfigure}
\begin{center}
\begin{tikzpicture}[mindmap, grow cyclic, every node/.style={concept, text=white, font=\small\bfseries},
  root concept/.append style={concept color=popdark, font=\normalsize\bfseries},
  level 1/.append style={level distance=3.4cm, sibling angle=60, font=\footnotesize},
  level 2/.append style={level distance=2.1cm, font=\scriptsize}]
\node[root concept]{Applications linéaires du plan}
  child[concept color=popblue]{ node {Définition}
    child { node {$f(\vec{u}+\vec{v})=f(\vec{u})+f(\vec{v})$} }
    child { node {$f(\lambda\vec{u})=\lambda f(\vec{u})$} }
    child { node {$f(\vec{0})=\vec{0}$} }
  }
  child[concept color=poporange]{ node {Exemples géométriques}
    child { node {Homothéties $k\vec{u}$} }
    child { node {Symétries, Projections} }
    child { node {Rotations vectorielles} }
  }
  child[concept color=popgreen]{ node {Matrice dans une base}
    child { node {Colonnes = $f(\vec{e_1}), f(\vec{e_2})$} }
    child { node[align=center] {$M=\begin{pmatrix}a&b\\c&d\end{pmatrix}$} }
    child { node {$f(x,y)=(ax+by, cx+dy)$} }
  }
  child[concept color=poppurple]{ node {Noyau $\ker f$}
    child { node {Système homogène $M\vec{X}=0$} }
    child { node {Sous-espace : $\{0\}$, droite, plan} }
    child { node {Base : vecteur directeur} }
  }
  child[concept color=poppink]{ node {Image $\operatorname{Im}f$}
    child { node {Engendré par colonnes de $M$} }
    child { node {Rang = dim(Im) = 0, 1 ou 2} }
    child { node {Éq. cartésienne si droite} }
  }
  child[concept color=popgold]{ node {Bijectivité $\iff$}
    child { node {$\ker f = \{\vec{0}\}$} }
    child { node {$\operatorname{Im}f = \mathcal{V}$} }
    child { node {$\det(M) = ad-bc \neq 0$} }
  };
\end{tikzpicture}
\end{center}
\legende{Synthèse visuelle : les 6 piliers de la leçon M19.}
\end{popfigure}

\begin{aretenir}[Checklist finale avant le Baccalauréat]
\begin{itemize}
  \item[$\square$] Je sais énoncer la définition (additivité + homogénéité) et le critère unique $f(\lambda\vec{u}+\mu\vec{v})=\lambda f(\vec{u})+\mu f(\vec{v})$.
  \item[$\square$] Je sais que $f(\vec{0})=\vec{0}$ est une condition \textbf{nécessaire} (test rapide d'élimination).
  \item[$\square$] Je connais les 4 exemples géométriques types : homothétie, symétrie, projection, rotation (et leurs matrices).
  \item[$\square$] Je distingue \textbf{endomorphisme} (linéaire $\mathcal{V}\to\mathcal{V}$), \textbf{isomorphisme} (linéaire bijective) et \textbf{automorphisme} (endo bijectif).
  \item[$\square$] Je sais construire la matrice de $f$ dans une base $\mathcal{B}$ : \textbf{colonnes = coordonnées de $f(\vec{e_1}), f(\vec{e_2})$ dans $\mathcal{B}$}.
  \item[$\square$] Je calcule l'image d'un vecteur $\vec{u}=x\vec{e_1}+y\vec{e_2}$ par $f(\vec{u}) = x f(\vec{e_1}) + y f(\vec{e_2})$ (produit matrice-vecteur).
  \item[$\square$] Je détermine $\ker f$ en résolvant $M\vec{X}=\vec{0}$ et j'en donne une base (ou $\varnothing$ si noyau nul).
  \item[$\square$] Je détermine $\operatorname{Im}f = \operatorname{Vect}(\text{colonnes de } M)$, j'en extrais une base et, si c'est une droite, j'en donne l'équation cartésienne (déterminant nul).
  \item[$\square$] Je justifie que $\ker f$ et $\operatorname{Im}f$ sont stables par addition et multiplication par un réel (ce sont des sous-espaces vectoriels).
  \item[$\square$] Je maîtrise les \textbf{3 caractérisations équivalentes de la bijectivité} : $\ker f=\{\vec{0}\}$ / $\operatorname{Im}f=\mathcal{V}$ / $\det(M)\neq 0$.
  \item[$\square$] Je calcule le déterminant $ad-bc$ (attention au signe $-$ !) et j'en déduis l'inversibilité.
  \item[$\square$] Je rédige proprement : j'annonce la base utilisée, j'écris les systèmes, je conclus par une phrase.
\end{itemize}
\end{aretenir}

\begin{exoresolu}[Exercice type Bac complet]
Soit $f$ l'endomorphisme de $\mathbb{R}^2$ défini dans la base canonique $\mathcal{B}=(\vec{i},\vec{j})$ par :
$f(x,y) = (x+2y, 2x+4y)$.
\begin{enumerate}
  \item Montrer que $f$ est linéaire et donner sa matrice $M$.
  \item Déterminer $\ker f$ et $\operatorname{Im}f$. Donner une base de chacun et une équation cartésienne de $\operatorname{Im}f$.
  \item $f$ est-elle bijective ? Justifier.
  \item Soit $\vec{u} = (3, -1)$. Calculer $f(\vec{u})$.
\end{enumerate}
\tcblower
\textbf{1.} $f(0,0)=(0,0)$. Forme analytique sans terme constant $\implies$ Linéaire.
$f(\vec{i}) = f(1,0) = (1,2) = \vec{i}+2\vec{j}$.
$f(\vec{j}) = f(0,1) = (2,4) = 2\vec{i}+4\vec{j}$.
$M = \begin{pmatrix} 1 & 2 \\ 2 & 4 \end{pmatrix}$.

\textbf{2.} \textbf{Noyau :} $\begin{cases} x+2y=0 \\ 2x+4y=0 \end{cases} \iff x = -2y$.
Solutions : $(-2y, y) = y(-2, 1)$.
$\ker f = \operatorname{Vect}(-2\vec{i}+\vec{j})$. Base : $(-2\vec{i}+\vec{j})$. $\dim=1$.

\textbf{Image :} Colonnes de $M$ : $\vec{c_1}=\binom{1}{2}$, $\vec{c_2}=\binom{2}{4}=2\vec{c_1}$. Colinéaires.
$\operatorname{Im}f = \operatorname{Vect}(\binom{1}{2})$ (droite vectorielle $y=2x$).
Base : $(\vec{i}+2\vec{j})$.
Équation caractéristique : $\det(\binom{1}{2}, \binom{x}{y}) = 1\cdot y - 2\cdot x = 0 \iff y - 2x = 0$.

\textbf{3.} $\det(M) = 1\times 4 - 2\times 2 = 0$.
$f$ n'est \textbf{pas bijective} (ni injective car $\ker f \neq \{\vec{0}\}$, ni surjective car $\operatorname{Im}f \neq \mathbb{R}^2$).

\textbf{4.} $f(3,-1) = (3+2(-1), 2\times 3 + 4(-1)) = (1, 2)$.
Ou par matrice : $M \binom{3}{-1} = \binom{1}{2}$.
\end{exoresolu}

\begin{exercice}[Entraînement personnel]
\begin{enumerate}
  \item Soit $g(x,y) = (2x-y+1, x+3y)$. $g$ est-elle linéaire ? Justifier.
  \item Soit $h$ l'endomorphisme de matrice $M = \begin{pmatrix} 3 & 1 \\ -1 & 2 \end{pmatrix}$ dans la base canonique.
  \begin{enumerate}
    \item Calculer $\det(M)$. $h$ est-elle bijective ?
    \item Déterminer $h^{-1}(4, 1)$ (le vecteur $\vec{u}$ tel que $h(\vec{u})=(4,1)$).
  \end{enumerate}
  \item Dans une base $\mathcal{B}=(\vec{u},\vec{v})$, l'endomorphisme $p$ a pour matrice $P = \begin{pmatrix} 2 & -1 \\ 4 & -2 \end{pmatrix}$.
  \begin{enumerate}
    \item Étudier la bijectivité de $p$.
    \item Déterminer $\ker p$ et $\operatorname{Im}p$ (bases et équation pour l'image).
    \item Vérifier la relation de dimension : $\dim(\ker p) + \dim(\operatorname{Im}p) = 2$.
  \end{enumerate}
\end{enumerate}
\end{exercice}

\begin{corrige}[Exercice n°1]
$g(0,0) = (1,0) \neq (0,0)$ donc $g$ n'est \textbf{pas linéaire} (c'est une application affine).
\end{corrige}

\begin{corrige}[Exercice n°2]
\begin{enumerate}
  \item $\det(M) = 3\times 2 - 1\times(-1) = 6+1 = 7 \neq 0$. $h$ est \textbf{bijective} (automorphisme).
  \item $M^{-1} = \frac{1}{7}\begin{pmatrix} 2 & -1 \\ 1 & 3 \end{pmatrix}$.
  $h^{-1}(4,1) = M^{-1} \binom{4}{1} = \frac{1}{7}\binom{2\times4 - 1}{4 + 3} = \frac{1}{7}\binom{7}{7} = \binom{1}{1}$.
  Donc $h^{-1}(4,1) = \vec{i}+\vec{j}$.
\end{enumerate}
\end{corrige}

\begin{corrige}[Exercice n°3]
\begin{enumerate}
  \item $\det(P) = 2\times(-2) - (-1)\times 4 = -4 + 4 = 0$. $p$ n'est \textbf{pas bijective}.
  \item \textbf{Noyau :} $\begin{cases} 2x - y = 0 \\ 4x - 2y = 0 \end{cases} \iff y = 2x$.
  Vecteurs $(x, 2x) = x(1, 2)$. $\ker p = \operatorname{Vect}(\vec{u}+2\vec{v})$. Base : $(\vec{u}+2\vec{v})$. $\dim=1$.
  \item \textbf{Image :} Colonnes $\binom{2}{4}$ et $\binom{-1}{-2}$. Colinéaires (2ème = $-\frac{1}{2}$ 1ère).
  $\operatorname{Im}p = \operatorname{Vect}(2\vec{u}+4\vec{v}) = \operatorname{Vect}(\vec{u}+2\vec{v})$.
  Base : $(\vec{u}+2\vec{v})$.
  Équation cartésienne : $\det(\binom{2}{4}, \binom{x}{y}) = 2y - 4x = 0 \iff y - 2x = 0$.
  \item $\dim(\ker p) + \dim(\operatorname{Im}p) = 1 + 1 = 2$. Relation vérifiée.
\end{enumerate}
\end{corrige}

\findecours

\end{document}
